\documentclass[11pt]{amsart}
\IfFileExists{lmodern.sty}%
  {\usepackage[T1]{fontenc}\usepackage{lmodern}\usepackage{microtype}}%
  {\usepackage[expansion=false]{microtype}}
\usepackage{amssymb,mathtools}
\usepackage{booktabs,array}
\usepackage{enumitem}
\emergencystretch=2.5em
\usepackage[colorlinks=true,linkcolor=blue,citecolor=blue,urlcolor=blue]{hyperref}

% ---------------------------------------------------------------------------
% Location of the Lean formalization.
\newcommand{\FormalizationURL}{https://github.com/slitsey/catalan-irrationality-measure}
\newcommand{\FormalizationLink}{\url{\FormalizationURL}}
% ---------------------------------------------------------------------------

\newtheorem{theorem}{Theorem}[section]
\newtheorem{lemma}[theorem]{Lemma}
\newtheorem{proposition}[theorem]{Proposition}
\newtheorem{corollary}[theorem]{Corollary}
\theoremstyle{definition}
\newtheorem{definition}[theorem]{Definition}
\theoremstyle{remark}
\newtheorem{remark}[theorem]{Remark}
\numberwithin{equation}{section}
\setcounter{tocdepth}{1}

\newcommand{\Z}{\mathbb Z}
\newcommand{\Q}{\mathbb Q}
\newcommand{\R}{\mathbb R}
\newcommand{\C}{\mathbb C}
\newcommand{\F}{\mathbb F}
\newcommand{\N}{\mathbb N}
\newcommand{\den}{\operatorname{den}}
\newcommand{\rank}{\operatorname{rank}}
\newcommand{\corank}{\operatorname{corank}}
\newcommand{\lean}[1]{\nolinkurl{#1}}
\renewcommand{\Re}{\operatorname{Re}}
\renewcommand{\Im}{\operatorname{Im}}
\newcommand{\Lw}{\mathsf L}          % Loewner matrix
\newcommand{\cO}{\mathcal O}
\newcommand{\Id}{\mathrm{Id}}
\newcommand{\err}{\operatorname{err}}
\newcommand{\cI}{\mathcal I}

\newcommand{\MuLamU}{31.2876}
\newcommand{\MuThirdsU}{9.4998}
\newcommand{\MuFifthsU}{23.8272}
\newcommand{\MuSeventhsU}{1806.1719}
\newcommand{\MuChi}{13.7046}
\newcommand{\MuHalvesU}{6.5253}
\newcommand{\MuPlaneA}{20.6622}
\newcommand{\MuPlaneB}{335.20}
\newcommand{\MuZetaThree}{6.0382}
\newcommand{\MuDipA}{14}
\newcommand{\MuDipB}{15.359}
\newcommand{\MuDipC}{18.161}
\newcommand{\LamFour}{3.5149036}  \newcommand{\DFour}{3.7447213}  \newcommand{\gFour}{7.1904273}  \newcommand{\ExpFour}{31.28751}
\newcommand{\LamThree}{3.3104458} \newcommand{\DThree}{4.0291249} \newcommand{\gThree}{6.8272548} \newcommand{\ExpThree}{9.499725}
\newcommand{\LamFive}{3.5128488}  \newcommand{\DFive}{3.8135797}  \newcommand{\gFive}{7.1655663}  \newcommand{\ExpFive}{23.82717}
\newcommand{\LamSeven}{3.6391158643} \newcommand{\DSeven}{3.6432250312} \newcommand{\gSeven}{7.4218615374} \newcommand{\ExpSeven}{1806.172}
\newcommand{\LamHalf}{2.9923824}  \newcommand{\DHalf}{3.9785713}  \newcommand{\gHalf}{6.4350874}  \newcommand{\ExpHalf}{6.525207}
\newcommand{\LamChi}{4.0959278}   \newcommand{\DChi}{4.6435117}   \newcommand{\gChi}{7.5043742}  \newcommand{\ExpChi}{13.70452}
\newcommand{\LamPl}{3.2227428}    \newcommand{\DPl}{3.5503512}    \newcommand{\gPl}{6.7690926}   \newcommand{\ExpPl}{20.66214}
\newcommand{\LamDA}{3.0668400}    \newcommand{\ExpDA}{13.99987}
\newcommand{\LamDB}{3.1096062}    \newcommand{\ExpDB}{15.35830}
\newcommand{\LamDC}{3.1776039}    \newcommand{\ExpDC}{18.16000}
\newcommand{\LamPlB}{3.4768376}   \newcommand{\DPlB}{3.4982052}   \newcommand{\gPlB}{7.1600631}  \newcommand{\ExpPlB}{335.0886}
\newcommand{\LamZ}{2.912839}     \newcommand{\ExpZ}{6.038184}
\newcommand{\MuChiExact}{13.704516}

\begin{document}

\title[Irrationality measures for $L(2,\chi_{-3})$ and other trigamma differences]{Irrationality measures for
$L(2,\chi_{-3})$ and other differences of trigamma values}

\author{Sean Litsey}
\email{sean.litsey.math@gmail.com}

\date{November 8, 2026}

\subjclass[2020]{Primary 11J82; Secondary 11J72, 15A15, 31A15, 68V20}

\keywords{Irrationality exponent, trigamma function, $L(2,\chi_{-3})$, $\zeta(3)$, L\"owner matrices, $p$-adic valuations,
Hermite interpolation, Andr\'eief's identity, logarithmic energy, formal verification, Lean}

\begin{abstract}
Let $L=L(2,\chi_{-3})$. We prove that $\mu(aL+b\pi^2)\le\MuPlaneA$ for all rationals $a$ and $b$, not both zero. For
particular directions in this plane the bound is better. Among them are
\[
 \mu(L)\le\MuChi,\qquad \mu\Bigl(L\pm\frac{\pi^2}9\Bigr)\le\MuThirdsU,\qquad
 \mu\Bigl(L\pm\frac{2\pi^2}{27}\Bigr)\le\MuDipA .
\]
The best published bound for these numbers was $24781$, due to Calegari, Dimitrov and Tang. We also prove
$\mu(a\zeta(3)+b\pi^2)\le\MuPlaneB$ for all rationals $a$ and $b$, not both zero. This second plane contains the
value $\psi''(1)=-2\zeta(3)$ of the next derivative. For $ab\ne0$ these numbers are irrational by a theorem of
Q.~Tang. As far as we know, no irrationality measure for them was known. Finally we prove
\[
 \mu\Bigl(G\pm\frac{5\pi^2}{48}\Bigr)\le\MuLamU,
\]
together with $\mu(\psi'(\frac k5)-\pi^2/6)\le\MuFifthsU$ and $\mu(\psi'(\frac k7)-\pi^2/6)\le\MuSeventhsU$ for $k$
prime to the denominator. For these numbers we know of no other proof of irrationality.

The approximations are determinants of rectangular L\"owner matrices of rational combinations of the digamma
function and its derivative, at integer nodes. Their denominators are bounded prime by prime. On the first plane one
table of local bounds and one decay certificate serve every point. On the second plane they serve every point except
the multiples of $\pi^2$. Nonvanishing at rational
points follows from Hermite interpolation in characteristic zero. The size of the determinants is bounded by
Andr\'eief's identity and an energy estimate. The proofs are formalized in Lean~4.
\end{abstract}

\maketitle
\tableofcontents

%% =====================================================================
\section{Introduction}\label{sec:intro}

\subsection{The results}
The values $\psi'(s)=\sum_{n\ge0}(n+s)^{-2}$ of the trigamma function at rational points $s=k/q$ are among the
simplest constants of analysis. Each lies in the $\Q$-span of $\pi^2$ and the real and imaginary parts of the values
$L(2,\chi)$ for the Dirichlet characters $\chi$ modulo $q$. For $q\le2$ they are rational multiples of $\pi^2$. Already for $q=3$ and
$q=4$ they involve $L(2,\chi_{-3})$ and Catalan's constant $G=L(2,\chi_{-4})$, whose irrationality was proved only in
2024 and 2026 \cite{CDT2024,OpenAI2026}. Throughout the paper we work with the differences
\[
 \Psi(s,t)=\psi'(s)-\psi'(t)=\sum_{n=0}^\infty\Bigl(\frac1{(n+s)^2}-\frac1{(n+t)^2}\Bigr)\qquad(0<s<t\le1).
\]

The \emph{irrationality exponent} $\mu(\xi)$ of a real number $\xi$ is the infimum of the real numbers $\mu$ such that
$|\xi-a/b|>b^{-\mu}$ for all integers $a$ and all sufficiently large integers $b$. It does not change under
$\xi\mapsto r\xi+r'$ with $r,r'\in\Q$ and $r\ne0$. A finite bound for $\mu(\xi)$ is an irrationality measure for
$\xi$, and it shows that $\xi$ is irrational.

We write
\[
 L=L(2,\chi_{-3})=\sum_{n\ge0}\Bigl(\frac1{(3n+1)^2}-\frac1{(3n+2)^2}\Bigr)=0.78130\,24128\ldots
\]
The trigamma values with denominators $2$, $3$ and $6$ all lie in the $\Q$-span of $1$, $L$ and $\pi^2$:
\begin{gather*}
 \psi'(\tfrac13)=\frac{2\pi^2}3+\frac92L,\qquad \psi'(\tfrac23)=\frac{2\pi^2}3-\frac92L,\\
 \psi'(\tfrac16)=2\pi^2+\frac{45}2L,\qquad \psi'(\tfrac56)=2\pi^2-\frac{45}2L,\\
 \psi'(\tfrac12)=\frac{\pi^2}2,\qquad \psi'(1)=\frac{\pi^2}6 .
\end{gather*}
The first pair follows from the reflection formula $\psi'(s)+\psi'(1-s)=\pi^2/\sin^2\pi s$ and from
$\psi'(\frac13)-\psi'(\frac23)=9L$. For the second pair, $\psi'(\frac16)-\psi'(\frac56)=36\sum_n((6n+1)^{-2}-(6n+5)^{-2})$,
and splitting the series for $L$ by residues modulo $6$ shows that this sum is $\frac54L$. So every difference
$\Psi(s,t)$ with $s,t\in\{\frac16,\frac13,\frac12,\frac23,\frac56,1\}$ is a rational point of the plane spanned by $L$
and $\pi^2$ modulo $\Q$. This is our first plane.

\begin{theorem}[first plane]\label{thm:plane1}
Let $a$ and $b$ be rational numbers, not both zero. Then $aL+b\pi^2$ is irrational and
\[
 \mu(aL+b\pi^2)\le\MuPlaneA .
\]
\end{theorem}

At some points of the plane we prove better bounds. These bounds do not follow from Theorem~\ref{thm:plane1}. Each has
its own proof, with its own frame and geometry or with extra local bounds. The point $L$ uses the function $F_\chi$ with
three digamma terms, its own geometry, its own local bounds and its own decay certificate (Section~\ref{sec:chi3-quarter}). The
points $L\pm\pi^2/9$ use the differences $F_{1/3}$ and $F_{2/3}$ of two digamma functions, with their own tables for
$\Lambda$ and their own certificate. The points in (iii)--(v) below use the functions of the plane, but with further
local bounds that hold only at these weights. They come from the rank of the reduction modulo $p$, in an
assignment form (Section~\ref{ssec:levels}). We state the points in the
notation $\Psi$, with the classical form beside it.

\begin{theorem}[special points of the first plane]\label{thm:points1}
The following bounds hold.
\begin{enumerate}[label=(\roman*),leftmargin=*]
\item $\mu(L)\le\MuChi$. Here $L=\frac19\Psi(\frac13,\frac23)$.
\item $\mu(L\pm\frac{\pi^2}9)\le\MuThirdsU$. Here $\Psi(\frac13,1)=\frac92(L+\frac{\pi^2}9)$ and
 $\Psi(\frac23,1)=-\frac92(L-\frac{\pi^2}9)$.
\item $\mu(\Psi(\frac16,\frac13))\le\MuDipA$ and $\mu(\Psi(\frac23,\frac56))\le\MuDipA$. Here
 $\Psi(\frac16,\frac13)=18(L+\frac{2\pi^2}{27})$ and $\Psi(\frac23,\frac56)=18(L-\frac{2\pi^2}{27})$.
\item $\mu(\Psi(\frac13,\frac12))\le\MuDipB$ and $\mu(\Psi(\frac12,\frac23))\le\MuDipB$. Here
 $\Psi(\frac13,\frac12)=\frac92(L+\frac{\pi^2}{27})$ and $\Psi(\frac12,\frac23)=\frac92(L-\frac{\pi^2}{27})$.
\item $\mu(L\pm\frac{\pi^2}{18})\le\MuDipC$. Here $L+\frac{\pi^2}{18}=\frac1{18}(4\Psi(\frac13,\frac12)+\Psi(\frac12,1))$
 and $L-\frac{\pi^2}{18}=\frac1{18}(4\Psi(\frac12,\frac23)-\Psi(\frac12,1))$.
\end{enumerate}
\end{theorem}

The second plane is spanned by $\zeta(3)$ and $\pi^2$ modulo $\Q$. It contains the values $\psi''(1)=-2\zeta(3)$ and
$\psi'(1)$, $\psi'(\frac12)$, and so the numbers $\zeta(2)+\zeta(3)$ and $\zeta(3)-r\zeta(2)$ for rational $r$.

\begin{theorem}[second plane]\label{thm:plane2}
Let $a$ and $b$ be rational numbers, not both zero. Then
\[
 \mu(a\zeta(3)+b\pi^2)\le\MuPlaneB .
\]
\end{theorem}

For $ab\ne0$ this is, as far as we know, the first irrationality measure for these numbers. Their irrationality is due
to Q.~Tang, who proved in Lean that $1$, $\zeta(2)$ and $\zeta(3)$ are linearly independent over $\Q$ \cite{Tang26}.
At the two ends of the plane we also prove the following bounds, each by its own argument. For $\zeta(3)$ the function
is $\psi'(z+1)$, with local bounds of its own and the decay certificate of $\pi^2$ (\S~\ref{ssec:zeta3}). For $\pi^2$
the function is $F_{1/2}$, with its own table for $\Lambda$. We include them for completeness. Both are weaker than the bounds
$\mu(\zeta(3))<5.4197$ and $\mu(\pi^2)<4.62675$ of Leibtag et al.\ \cite{LBKWSK26}.

\begin{theorem}[the ends of the second plane]\label{thm:ends}
$\mu(\zeta(3))\le\MuZetaThree$ and $\mu(\pi^2)\le\MuHalvesU$.
\end{theorem}

The last three results concern numbers outside both planes. Here
\[
 G=L(2,\chi_{-4})=\sum_{n\ge0}\Bigl(\frac1{(4n+1)^2}-\frac1{(4n+3)^2}\Bigr)=0.91596\,55941\ldots,
\]
and $\psi'(\frac14)=\pi^2+8G$ and $\psi'(\frac34)=\pi^2-8G$, so that
$\Psi(\frac14,1)=8(G+\frac{5\pi^2}{48})$ and $\Psi(\frac34,1)=-8(G-\frac{5\pi^2}{48})$. For $q=5$ and $q=7$ there is
no single classical constant. For $\gcd(k,q)=1$ the number $\Psi(\frac kq,1)$ lies in the $\Q$-span of $\pi^2$ and
the real and imaginary parts of the values $L(2,\chi)$ for the Dirichlet characters $\chi$ modulo $q$, since
\[
 \psi'(\tfrac kq)=q^2\sum_{n\equiv k\ (q)}\frac1{n^2}=\frac{q^2}{\varphi(q)}\sum_{\chi\bmod q}\bar\chi(k)L(2,\chi).
\]
For $q=5$ and $q=7$ some of the coefficients $\bar\chi(k)$ are not real, and the conjugate terms pair up.

\begin{theorem}\label{thm:main}\label{thm:offplane}
The following bounds hold:
\begin{enumerate}[label=(\roman*),leftmargin=*]
\item $\mu(G\pm\frac{5\pi^2}{48})\le\MuLamU$.
\item $\mu(\psi'(\frac k5)-\frac{\pi^2}6)\le\MuFifthsU$ for $k=1,2,3,4$.
\item $\mu(\psi'(\frac k7)-\frac{\pi^2}6)\le\MuSeventhsU$ for $k=1,\dots,6$.
\end{enumerate}
\end{theorem}

Each bound $\mu(\xi)\le\mu_0$ above is proved in the stronger form $|\xi-a/b|>b^{-\mu_0}$ for all integers $a$ and
all sufficiently large integers $b$. Table~\ref{tab:main} collects the results.

\begin{table}[ht]
\centering
{
\begin{tabular}{@{}llll@{}}
\toprule
number & form in $\Psi$ & $\mu\le$ & best published before\\
\midrule
$aL+b\pi^2$ & -- & $\MuPlaneA$ & $24781$ \cite{CDT2025}\\
$L$ & $\frac19\Psi(\frac13,\frac23)$ & $\MuChi$ & $24781$\\
$L\pm\frac{\pi^2}9$ & $\Psi(\frac13,1)$, $\Psi(\frac23,1)$ & $\MuThirdsU$ & $24781$\\
$L\pm\frac{2\pi^2}{27}$ & $\Psi(\frac16,\frac13)$, $\Psi(\frac23,\frac56)$ & $\MuDipA$ & $24781$\\
$L\pm\frac{\pi^2}{27}$ & $\Psi(\frac13,\frac12)$, $\Psi(\frac12,\frac23)$ & $\MuDipB$ & $24781$\\
$L\pm\frac{\pi^2}{18}$ & -- & $\MuDipC$ & $24781$\\
\addlinespace
$a\zeta(3)+b\pi^2$, $ab\ne0$ & -- & $\MuPlaneB$ & none\\
$\zeta(3)$ & -- & $\MuZetaThree$ & $5.4197$ \cite{LBKWSK26}\\
$\pi^2$ & $3\Psi(\frac12,1)$ & $\MuHalvesU$ & $4.62675$ \cite{LBKWSK26}\\
\addlinespace
$G\pm\frac{5\pi^2}{48}$ & $\Psi(\frac14,1)$, $\Psi(\frac34,1)$ & $\MuLamU$ & none\\
$\psi'(\frac k5)-\frac{\pi^2}6$ & $\Psi(\frac k5,1)$ & $\MuFifthsU$ & none\\
$\psi'(\frac k7)-\frac{\pi^2}6$ & $\Psi(\frac k7,1)$ & $\MuSeventhsU$ & none\\
\bottomrule
\end{tabular}}
\caption{The results of this paper. The first column is a number in common notation, the second is a form of the
same number in the notation $\Psi$ up to a rational factor, where one exists. In the first row $a$ and $b$ are
rationals, not both zero. The bound $\MuPlaneB$ also holds for all such $a\zeta(3)+b\pi^2$. ``None'' means that we know of no irrationality measure. A bound of about $9.7\cdot10^3$
for the first plane follows from an unrefereed supplement to \cite{LBKWSK26}, as explained in
\S\ref{ssec:background}.}\label{tab:main}
\end{table}

Several numbers named in the literature lie on the first plane but at none of the points of
Theorem~\ref{thm:points1}. For them Theorem~\ref{thm:plane1} gives $\mu\le\MuPlaneA$ in place of $24781$. These are
$\psi'(\frac13)=\frac92(L+\frac{4\pi^2}{27})$, $\psi'(\frac23)$, $\psi'(\frac16)=\frac{45}2(L+\frac{4\pi^2}{45})$ and
$\psi'(\frac56)$, whose irrationality Calegari, Dimitrov and Tang state explicitly \cite[Corollary~B]{CDT2024}, and
$\Psi(\frac16,1)=\psi'(\frac16)-\zeta(2)=\frac{45}2(L+\frac{11\pi^2}{135})$ and $\Psi(\frac56,1)$. The same holds for the
remaining sixths, $\Psi(\frac16,\frac12)=\frac{45}2(L+\frac{\pi^2}{15})$, $\Psi(\frac16,\frac23)=27(L+\frac{4\pi^2}{81})$
and their mirrors $\Psi(\frac12,\frac56)$, $\Psi(\frac13,\frac56)$. Some numbers with similar names are not on the
plane. The Clausen value $\mathrm{Cl}_2(\frac\pi3)=\frac{3\sqrt3}4L$ and the Mahler measure
$m(1+x+y)=\frac{3\sqrt3}{4\pi}L$ carry the factors $\sqrt3$ and $\sqrt3/\pi$, and our results say nothing about them.
On the second plane, Theorem~\ref{thm:plane2} gives, as far as we know, the first measures for $\zeta(2)+\zeta(3)$ and
for $\zeta(3)-r\zeta(2)$ with $r\in\Q^\times$. It gives none for the ratio $\zeta(3)/\pi^2$, whose irrationality also
follows from Q.~Tang's theorem \cite{Lai2026}.

All the theorems above have been verified in the Lean~4 proof assistant, together with every theorem, proposition and
lemma used in their proofs. The formal proofs have no unproved steps and use only the standard axioms.
Section~\ref{sec:lean} describes them. The remarks and Section~\ref{sec:limits} discuss the method and are not part of
the formal proofs. The formalization is available at \FormalizationLink. The use of AI in this work is disclosed at the
end of the paper.

\subsection{Background}\label{ssec:background}
\emph{The value $L(2,\chi_{-3})$ and the first plane.} Calegari, Dimitrov and Tang \cite{CDT2024} proved that $1$, $\pi^2$ and $L$ are linearly independent
over $\Q$. Their proof applies a new arithmetic holonomy bound to a construction of Zagier. They list the
irrationality of $\psi'(\frac13)$, $\psi'(\frac23)$, $\psi'(\frac16)$ and $\psi'(\frac56)$ as a corollary, and remark that
$\psi'(\frac n6)$ is irrational for every $n$. In \cite{CDT2025} they made holonomy bounds effective. Theorem~5.1 there gives
$\mu(L)\le24781$. The remark that follows it states that the same theorem and proof hold for any $\Q$-linear
combination of $L$ and $\pi^2$. So $24781$ is the best published bound at every point of the first plane.
Theorem~\ref{thm:plane1} improves it to $\MuPlaneA$ everywhere, and Theorem~\ref{thm:points1} improves it further at
five pairs of points. Leibtag et al.\ \cite{LBKWSK26} find numerically that Ap\'ery-like approximations from conservative
matrix fields have positive approximation exponents for $L\pm\pi^2/9$. A supplement to their paper uses one such
field to give an Ap\'ery-style proof of the theorem of Calegari, Dimitrov and Tang. For every point of the plane it
proves an approximation exponent $\delta>0.000103$. With the bound $\mu\le1+1/\delta$ of their \S2 this gives
$\mu\lesssim9.7\cdot10^3$, a consequence not stated there. The supplement is unrefereed and is not part of the arXiv
version.

\emph{The values $\zeta(2)$, $\zeta(3)$ and the second plane.} Ap\'ery proved that $\zeta(2)$ and $\zeta(3)$ are
irrational \cite{Apery79}, and Beukers gave his integral proofs \cite{Beukers79}. Rhin and Viola obtained
$\mu(\zeta(2))<5.441243$ and $\mu(\zeta(3))<5.513891$ with permutation groups of integrals \cite{RhinViola96,RhinViola01}.
Zudilin lowered the bound for $\zeta(2)$ to $5.095412$ \cite{Zudilin14}. Giraudin lowered the bounds to $5.0193784$ for
$\zeta(2)$ \cite{Giraudin26} and to $5.5138800$ for $\zeta(3)$ \cite{Giraudin26z3}. Leibtag, Barkan, Kalisch, Weinbaum,
Shalyt and Kaminer \cite{LBKWSK26} showed that the classical constructions are one-dimensional trajectories in
conservative matrix fields. By searching over other trajectories they obtained $\mu(\pi^2)<4.62675$ and
$\mu(\zeta(3))<5.4197$. They also proved the uniform bound $\mu(A\zeta(3)+B\zeta(4))<66.3$ on the plane spanned by
$\zeta(3)$ and $\pi^4$, which is the closest analogue of Theorem~\ref{thm:plane2}. The linear forms in $1$, $\zeta(2)$
and $\zeta(3)$ have a longer history. Dauguet and Zudilin \cite{DauguetZudilin14} bounded such forms from below when
the coefficients satisfy divisibility and size conditions. They noted that this does not give the linear independence of
$1$, $\zeta(2)$, $\zeta(3)$, which remained open. For a fixed number $a\zeta(3)+b\zeta(2)$ with $ab\ne0$ their theorem
gives no measure. In 2026 Q.~Tang \cite{Tang26} proved the linear independence in Lean. This proof is a formalization
with no accompanying written paper. Lai
\cite{Lai2026} then proved that $1$, $\zeta(3)$ and $\zeta(4)$ are linearly independent. As far as we know, no
irrationality measure for any $a\zeta(3)+b\pi^2$ with $ab\ne0$ was known before Theorem~\ref{thm:plane2}.

\emph{Catalan's constant.} In October 2026 OpenAI \cite{OpenAI2026} proved that $G$ is irrational and formalized the
proof in Lean. Sun's preprint \cite{Sun2026} had claimed a proof of the irrationality of $G$. Wachs \cite{Wachs2026}
reports that its main estimate is not justified. Finite irrationality exponents for $G$ were obtained after OpenAI's
proof by Jha \cite{Jha2026}, Carreira \cite{Carreira2026}, Bojnal \cite{Bojnal2026}, Xu \cite{Xu2026} and the author
\cite{Litsey2026}. Carreira \cite{Carreira2026b} also records attempts to carry the determinant method to $\beta(4)$. For
the arithmetic of $G$ and related values before these results see
\cite{RivoalZudilin2003,Zudilin2003,Zudilin02Cheb,KrattRivoal08,KrattZudilin19} and the survey \cite{Zudilin19Survey}.
Eskandari, Murty and Nemoto \cite{EMN2025} construct a $2$-dimensional mixed motive over $\Q$ with period $G$. They use
it to obtain linear forms in $1$ and $G$. Eskandari \cite{Eskandari2026} improves the exponent of rational approximations
to $G$ of Rivoal--Zudilin type from $0.52$ to $0.62$. These are explicit approximations, not a lower bound for all
rationals. Our searches found no proof that $G-5\pi^2/48$ or $G+5\pi^2/48$ is irrational. This does not follow from the
irrationality of $G$, since it is not known whether $1$, $G$ and $\pi^2$ are linearly independent over $\Q$.

\emph{The numbers $\psi'(\frac k5)-\pi^2/6$ and $\psi'(\frac k7)-\pi^2/6$.} Our searches found no irrationality result for
them. Prellberg \cite{Prellberg26} proved that $\mathrm{Li}_2(1/5)$ is irrational, but that is a different number.

\emph{Recent results.} Since September 2026 a series of new results has appeared in this area. Several of them were
obtained with AI systems. OpenAI proved $\mu(\pi)=2$, with a Lean formalization \cite{OpenAIPi2026}. This improves the
bound $\mu(\pi)<7.1033$ of Zeilberger and Zudilin \cite{ZeilbergerZudilin20}. Liu, Jiang and Zhang proved $\mu(\log r)=2$
for every rational $r>0$ with $r\ne1$ \cite{LiuJiangZhang26}, improving the classical bound $\mu(\log2)<3.5746$ of
Marcovecchio \cite{Marcovecchio09}. Their result also covers nonzero rational combinations of such logarithms. Prellberg
\cite{Prellberg26} proved that $\mathrm{Li}_2(-1/4)$, $\mathrm{Li}_2(1/5)$ and $\mathrm{Li}_2(-1/3)$ are irrational. These
are three cases left open in \cite{CDT2024}. The proof uses the integrals of Rhin--Viola and Viola--Zudilin. Fauzan
\cite{Fauzan2026} announced a proof that $\zeta(5)$ is irrational, using rationally normalized Hankel determinants.

\emph{Determinant methods.} Irrationality proofs built on determinants of large size have a precedent in Zudilin's
determinantal criterion with Hankel determinants \cite{Zudilin}. The preprint of Liu, Jiang and Zhang on logarithms
\cite{LiuJiangZhang26} uses a separated-weight interpolation determinant. The arguments of Fauzan and Lai
\cite{Fauzan2026,Lai2026} use Hankel determinants. Brown \cite{Brown26} proposes a multivariate
transfinite-diameter test for Mellin-type integrals and applies it to $\pi^2$. The framework we use is OpenAI's
\cite{OpenAI2026}. In it the decay of the determinants is controlled by Cauchy determinants and Andr\'eief's identity,
and the denominators are controlled prime by prime. Our determinants are L\"owner matrices, that is, matrices of
divided differences. They are a matrix-valued analogue of Pad\'e and Hermite--Pad\'e approximation
\cite{NikishinSorokin91}. Formal proofs of irrationality results in Lean also include the irrationality of $\zeta(3)$ by
Liu, Zhang and Zhi \cite{LZZ25} and the linear independence of $1$, $\zeta(2)$ and $\zeta(3)$ by Q.~Tang
\cite{Tang26}.

\subsection{The companion paper}
Catalan's constant is also a difference of trigamma values, $G=\frac1{16}\Psi(\frac14,\frac34)$, and the construction
applies to this pair. For it, however, our bound for the denominators exceeds our bound for the decay, so we obtain no
measure for $G$. The companion paper \cite{Litsey2026} refines OpenAI's determinant instead and proves
$\mu(G)\le270$. Our computations for $(\frac14,\frac34)$ are not part of the formal proof.

\subsection{What a measure requires}
Let $P(x)$ be a polynomial with rational coefficients for which $P(\xi)$ is small, and suppose that $|\xi-a/b|$ is very
small. Three estimates turn this into a lower bound for $|\xi-a/b|$.
\begin{enumerate}[label=(\roman*)]
\item \emph{Denominators.} The coefficients of $P$ have a common denominator of size $e^{(\Lambda+o(1))r^2}$.
\item \emph{Smallness on a neighbourhood.} $|P(x)|\le e^{-(D-o(1))r^2}$ for every $x$ with $|x-\xi|\le e^{-\rho r}$,
 not only at $\xi$.
\item \emph{Nonvanishing at the point.} $P(a/b)\ne0$ for some member of the family whose size can be chosen in terms
 of $b$.
\end{enumerate}
In OpenAI's argument for $G$ the size of the determinant is tied to a prime $p$, and nonvanishing is proved modulo $p$.
Here nonvanishing is proved in characteristic zero, so the size $r$ is free. A plane theorem needs more. A number
$aL+b\pi^2$ is $\xi=F'(0)$ for a combination $F$ of three digamma functions whose weights depend on $(a,b)$. The three
estimates must then hold with constants that do not depend on the weights. For the second plane the function $F$ has a
pole of order two, so the matrices are no longer Cauchy-like at every pole, and both the arithmetic and the analysis
need a new argument.

\subsection{New ideas}
The proofs rest on the following ideas. Section~\ref{sec:overview} explains the first five in more detail.
\begin{enumerate}[leftmargin=*]
\item \emph{L\"owner determinants of digamma combinations.} The approximations are determinants of L\"owner matrices,
 the matrices of divided differences, of a combination $F$ of $\psi(z+s)$ and $\psi'(z+s)$ at integer nodes. The rows
 and columns share a long block of nodes, which carries the variable. A further border of row nodes and column nodes
 alternates beyond it. The degree in the variable stays small, while the decay is that of a larger node set.
\item \emph{Denominators from local bounds.} For each prime we prove a list of upper bounds for the valuation of the
 common denominator. Each is an explicit integer expression in a few integers attached to the prime. Integrated
 against $\log p$, their minimum gives the denominator exponent, and checking it reduces to finitely many comparisons of
 rational numbers.
\item \emph{Local bounds that do not see the weights.} For integer weights the polar data at a prime $p$ depend only on
 the shifts. All our local bounds for the first plane read only these data. So one table of bounds serves every point of
 the plane. A prime dividing a weight can only lower the true valuation, and a weight zero removes a pole without
 changing the bounds.
\item \emph{A decay certificate uniform in the weights.} The energy inequalities behind the decay estimate contain no
 weight. The weights enter only through a constant factor in the external field, which costs $e^{O(r)}$. A weight
 zero only removes atoms. One trial measure on one contour therefore proves the decay for the whole first plane, and
 another one for the whole second plane.
\item \emph{Nonvanishing in characteristic zero.} A kernel vector of the matrix at a rational point gives a solution of a
 Hermite interpolation problem. A discrete Wronskian bounds the corank, because one pole of the increment of $F$ is
 isolated in its class modulo $\Z$. The bound depends only on the number of poles, counted with order.
\item \emph{Level returns and a gauge for the special points.} At a prime $p$, the partial sums of $F$ step by the
 residues at its poles. For special weights this level returns to an earlier value, and whole blocks of entries become
 $p$-integral. A row operation inside each residue class modulo $p$, by a $p$-integral multiple of a common row,
 gives some further rows the level of the return. This turns the rank of the reduction modulo $p$ into an assignment
 bound, and it gives the points of Theorem~\ref{thm:points1}(iii)--(v).
\item \emph{Elimination of complete periods for a double pole.} For the second plane the Schur complement of the
 L\"owner matrix over a union of complete periods modulo $p$ is again of L\"owner type, with generators that are
 constant on half-periods. This follows from the displacement identity of L\"owner matrices and from
 $\sum_{v\in\F_p^\times}v^{-1}=0$. A uniform pivot on complete double periods extends it to primes smaller than the border.
 The required nondegeneracy modulo $p$ is the Wronskian argument of item~5 in characteristic $p$.
\item \emph{Jets reduced to simple poles.} At a double pole the Gram representation of the L\"owner matrix carries a
 derivative. A rank-one splitting of this term, and Cauchy's estimate on small circles, bound every jet determinant by
 the Cauchy determinant of a fixed function with simple poles. The decay certificate of that function then applies.
\end{enumerate}
The constants of the proofs are listed in Table~\ref{tab:data}.

\subsection{Organization}
Section~\ref{sec:overview} explains the construction and the main ideas without technicalities, and states the general
deduction. Section~\ref{sec:construction} gives the construction. Sections~\ref{sec:content}, \ref{sec:nonvanishing}
and~\ref{sec:decay} prove the three estimates for differences of two digamma functions. They also prove nonvanishing and
decay for general combinations. Section~\ref{sec:plane1} proves Theorem~\ref{thm:plane1} and the special points of
Theorem~\ref{thm:points1}, and Section~\ref{sec:chi3-quarter} treats the point $L$ itself. Section~\ref{sec:plane2}
proves Theorem~\ref{thm:plane2}. Section~\ref{sec:proof} completes the proofs and gives the data.
Section~\ref{sec:limits} discusses the limits of the method. Section~\ref{sec:lean} describes the formalization. The
disclosure of AI use follows it.

\subsection{Notation and conventions}
We write $\psi$ for the digamma function, the logarithmic derivative of Euler's gamma function, and $v_p$ for the
$p$-adic valuation. For $P\in\Q[x]$, $\den P$ is the least positive integer whose product with $P$ lies in $\Z[x]$. We
write $[0,r)=\{0,1,\dots,r-1\}$. For finite sets $X_0$, $Y_0$ of integers and a matrix $A$ indexed by integers,
$A[X_0,Y_0]$ is the submatrix with rows $X_0$ and columns $Y_0$. We write $\|y\|$ for the distance from a real number
$y$ to $\Z$, and $\|m\|_p=p\|m/p\|$.
%% =====================================================================
\section{Overview of the method}\label{sec:overview}

\subsection{L\"owner matrices of digamma differences}
Let $F$ be a function analytic near a finite set of distinct points. The \emph{L\"owner matrix} \cite{Loewner34} of $F$
is the matrix of divided differences indexed by the points:
\[
 \Lw_F[t,t']=\frac{F(t')-F(t)}{t'-t}\ (t\ne t'),\qquad \Lw_F[t,t]=F'(t).
\]
It is a Cauchy-like structured matrix. Hermite's formula gives
$\Lw_F=\frac1{2\pi i}\oint F(z)\,\upsilon(z)\upsilon(z)^{\top}dz$ with $\upsilon_t(z)=1/(z-t)$. A representation of $F$ as
a Stieltjes or contour integral therefore exhibits $\Lw_F$ as a Gram matrix of Cauchy kernels.

For a rational shift $0<s<1$ with denominator $q$ we take the digamma difference $F_s$. For the pair
$(\frac13,\frac23)$ we take a combination $F_\chi$ of three digamma functions. The nodes are integers. We put
\begin{equation}\label{eq:frames}
\begin{gathered}
 F_s(z)=\psi(z+1)-\psi(z+s),\qquad \omega=\tfrac1{q^2},\quad \xi_s=-\tfrac1{q^2}\Psi(s,1),\\
 F_{\chi}(z)=-\tfrac13\psi(z+\tfrac13)+\tfrac12\psi(z+\tfrac12)-\tfrac16\psi(z+1),\\
 \omega=-\tfrac16,\quad \xi_\chi=\tfrac1{36}\Psi(\tfrac13,\tfrac23).
\end{gathered}
\end{equation}
In each case $\xi=\omega\,F'(0)$. The differences $F(t')-F(t)$ and $F'(t)-F'(0)$ are rational for integers $t,t'\ge0$.
Let $R$ and $Q$ be finite sets of nonnegative integer nodes, which index the rows and the columns. Then
\begin{equation}\label{eq:shift}
 \omega\,\Lw_F[R,Q]=M[R,Q]+\xi\,\Id_{R,Q},\qquad \Id_{R,Q}[t,t']=\mathbf 1[t=t'],
\end{equation}
with an explicit symmetric function $M$ on $\Z_{\ge0}^2$ with rational values, given in \S\ref{sec:construction}.
Replacing $\xi$ by a variable $x$ gives the matrix polynomial $\Phi(x)=M[R,Q]+x\Id_{R,Q}$. Its determinant $P(x)$ has
rational coefficients and degree at most $|R\cap Q|$, and $P(\xi)=\det(\omega\,\Lw_F[R,Q])$.

Suppose that $|\xi-a/b|$ is very small. Then $P(a/b)$ is very close to $P(\xi)$, which is tiny. On the other hand
$b^{|R\cap Q|}\den(P)\,P(a/b)$ is an integer. If it is not zero, we get a contradiction as soon as
\[
 \underbrace{\log\den P}_{\text{denominators}}+|R\cap Q|\log b\ <\ \underbrace{-\log|P(a/b)|}_{\text{decay}} .
\]
Three estimates are therefore needed. The \emph{arithmetic estimate} is an upper bound for $\log\den P$. The
\emph{analytic estimate}, or decay, is an upper bound for $|P(x)|$ near $\xi$. The \emph{nonvanishing} is that of
$P(a/b)$ for some member of a family of determinants whose size can be chosen in terms of $b$. Each estimate is proved once, for every shift,
and the families differ only in a few parameters.

\subsection{Rectangular L\"owner matrices with an alternating border}
We take a long common block $C=[0,r)$, which carries the variable $x$, and add $m$ row nodes and $m$ column nodes that
alternate beyond it:
\[
 X'=\{r,r+2,\dots,r+2m-2\},\qquad Y'=\{r+1,r+3,\dots,r+2m-1\}.
\]
The rows are $R=C\sqcup X'$ and the columns are $Q=C\sqcup Y'$ (Figure~\ref{fig:geometry}). The determinant
$P(x)=\det\Phi(x)$ has degree at most $r$, since the variable occurs only on the diagonal of the common block. So the
factor $b^{\deg P}$ in the contradiction is that of $r$ nodes, while the matrix has size $r+m$ and its decay is that of
about $(1+c/2)r$ nodes, where $c=2m/r$ is the length of the border relative to $r$. The denominators grow too, but less,
because border rows and border columns never share a node. To guarantee nonvanishing we also use the principal minors
that delete a bounded number of common nodes. A larger $c$ raises $D$ and $\Lambda$ together, and we chose $c$
numerically to minimise $\rho/(D-\Lambda)$. Its value for each constant is given in Table~\ref{tab:data}.

\begin{figure}[ht]
\centering
\setlength{\unitlength}{0.42cm}
\begin{picture}(20,13)(-6,-1)
 \put(0,0){\framebox(13,11){}}
 \put(8,0){\line(0,1){11}}
 \put(0,3){\line(1,0){13}}
 \multiput(0.5,10.5)(1,-1){8}{\makebox(0,0){$\scriptstyle x$}}
 \put(4,11.6){\makebox(0,0){$\scriptstyle 0,1,\dots,r-1$}}
 \put(10.5,11.6){\makebox(0,0){$\scriptstyle r+1,\,r+3,\,r+5$}}
 \put(-0.3,7){\makebox(0,0)[r]{$\scriptstyle 0,\dots,r-1$}}
 \put(-0.3,1.5){\makebox(0,0)[r]{$\scriptstyle r,\,r+2,\,r+4$}}
 \put(4,12.5){\makebox(0,0){$C$}}
 \put(10.5,12.5){\makebox(0,0){$Y'$}}
 \put(-4.8,7){\makebox(0,0)[r]{$C$}}
 \put(-4.8,1.5){\makebox(0,0)[r]{$X'$}}
 \put(4,1.5){\makebox(0,0){$M[X',C]$}}
 \put(10.5,7){\makebox(0,0){$M[C,Y']$}}
 \put(10.5,1.5){\makebox(0,0){$M[X',Y']$}}
\end{picture}
\caption{The geometry for $r=8$ and $m=3$. Rows are indexed by $C\sqcup X'$ and columns by $C\sqcup Y'$. The variable
$x$ occurs only on the diagonal of the common block. The border rows $X'=\{r,r+2,r+4\}$ and border columns
$Y'=\{r+1,r+3,r+5\}$ alternate and never meet the diagonal.}\label{fig:geometry}
\end{figure}

The interleaving of $X'$ and $Y'$ is essential. It keeps the two Cauchy determinants in Andr\'eief's formula comparable
(\S\ref{sec:decay}). It also gives the reduction of the matrix modulo a prime an alternating structure. Several of the
arithmetic bounds of \S\ref{sec:content} rest on this structure.

\subsection{Denominators from local bounds}
As in the work of Rhin and Viola \cite{RhinViola96,RhinViola01}, we bound the $p$-adic valuation of the common
denominator prime by prime and sum the bounds with the prime number theorem in arithmetic progressions. For $p\asymp r$
the valuation depends on $r$, $m$, $p$ and the first pole $k_p$, the least index $k_p\ge1$ at which the second digamma
term has a pole modulo $p$. One has $k_p=\kappa p+O(1)$ with $\kappa=j/q$, where the residue $j\equiv qs\,p^{-1}\pmod q$
is the \emph{root class} of $p$. The function $F_\chi$ has two first poles, and on each class of primes modulo $6$ both
are affine functions of one index of the same kind (Section~\ref{sec:chi3-quarter}).

The local bounds come from three sources: the valuations of the entries, the corank of the reduction modulo $p$, and
cancellations at second and third order in the Smith normal form. Each is an explicit piecewise linear expression in
$(r,m,p,k_p)$, valid in a range of $p/r$ described by linear inequalities. Replacing $(r,m,p,k_p)$ by
$(1,\alpha,\theta,\kappa\theta)$, with $\alpha=c/2$, turns it into a piecewise affine function of $\theta=p/r$, with an
explicitly bounded error. The \emph{denominator exponent} is
\[
 \Lambda=\frac1{\varphi(q)}\sum_j\int_0^\infty\mathcal P_j(\theta)\,d\theta,
\]
where $\mathcal P_j$ is the minimum of these functions for the root class $j$. A table of affine pieces above
$\mathcal P_j$ certifies $\Lambda$ by finitely many comparisons of rational numbers (Theorem~\ref{thm:content}). The
bounds depend on the shift only through the root class, so one table serves every shift with denominator $q$, and the
numbers $\Psi(s,1)$ for these shifts receive the same exponent.

\subsection{Primes smaller than the border}
For $p\le2m$ some residue classes modulo $p$ contain two border nodes, a row $t$ and a column $t+p$. The bounds above
are not sharp for these primes. Each such pair gives a pivot that is a $p$-adic unit, because the interval $(t,t+p]$
contains one pole of each digamma term and their contributions cannot cancel. Eliminating these pivots leaves a matrix
of the same shape with a border of $d=\|m\|_p$ rows and columns. To this smaller matrix we apply the separation
argument used for large primes, at third order. Two identities for bordered inverses show that the elimination does not
disturb the second-order terms (Lemma~\ref{lem:bordered}). A rational interpolation problem of type $(n,n)$ shows that
the remaining border nodes stay separated (Lemma~\ref{lem:harmonic}).

\subsection{Nonvanishing}
Let $x_0$ be a rational point. The kernel of $\Phi(x_0)$ is identified with a space of pairs of polynomials that satisfy
Hermite interpolation conditions at the nodes. Two such pairs have vanishing cross product, by a count of degrees and
zeros. A pair of low degree would give a rational antidifference of the increment of $F$. No such
antidifference exists, because a pole of the increment is isolated in its class modulo $\Z$. This bounds the corank of
$\Phi(x_0)$ by half the number of poles of $F$ in a period, counted with order. A principal restriction that
deletes at most that many common nodes is then nonsingular. This is Theorem~\ref{thm:nvframe}. No prime enters, so
the size $r$ can be chosen freely in terms of $b$.

\subsection{Decay}
Moving Hermite's contour to a polygonal path $\Gamma$ in the left half-plane writes $\omega\,\Lw_F[R,Q]$ as a Gram matrix
$\int\upsilon_t\upsilon_{t'}\,d\Omega$. Here $\Omega$ is a complex measure. It has signed atoms at poles of $F$ and an
exponentially small density on $\Gamma$. Andr\'eief's identity with two families of functions, one for the rows and one
for the columns, gives
\[
 |\det\Lw_F[R,Q]|\le\frac1{n_*!}\int\bigl|\mathsf C_R(z)\bigr|\,\bigl|\mathsf C_Q(z)\bigr|\prod_{i=1}^{n_*}d|\Omega|(z_i),
\]
where $n_*$ is the size of the matrix and $\mathsf C_R$, $\mathsf C_Q$ are Cauchy determinants. Taking logarithms turns
the integrand into a discrete logarithmic energy. The free points $z_i$ interact with the nodes in an external field.

We bound this energy from below by comparing it with a \emph{trial measure} $\sigma$, a density that is polynomial on
each of finitely many segments. The comparison uses the conditional positive definiteness of the logarithmic kernel.
What is left is a lower bound for one explicit function $\Upsilon$ on the contour, on the segment carrying the atoms and
on the interval of the nodes. We prove it on a finite cover of these sets by cells, each checked in exact rational
arithmetic. One theorem, Theorem~\ref{thm:decay}, gives the decay for every geometry parameter $c\le1$ and every
function $F$ whose poles are separated from the contour. Each family then supplies only $\sigma$, the polygon and the
cover.

\subsection{The deduction}
Write $\eta=D-\Lambda$, where $D$ is the \emph{decay rate}, and let $\rho$ be the \emph{neighbourhood exponent} of the
decay estimate: the decay holds at every point within $e^{-\rho r}$ of $\xi$. Section~\ref{sec:proof} combines the
three estimates into
\[
 \mu(\xi)\le\frac\rho\eta=\frac\rho{D-\Lambda}.
\]
This is the bound $d\rho/\eta$ of the companion paper \cite{Litsey2026} with degree $d=1$ per unit of scale.

\subsection{Formal verification}
Every step is formalized in Lean~4 with Mathlib \cite{Lean4,Mathlib}. This includes the local arithmetic bounds and the
conditions on the tables for $\Lambda$, and the interpolation argument for nonvanishing. It also includes the contour
representation, the cell checks of the decay estimate, and the prime number theorem in arithmetic progressions. The
formal proofs share one core, which
contains rectangular determinants, Gauss valuations, the local bounds, the tables for $\Lambda$, the selection in
characteristic zero, the decay theorem and the deduction. A family supplies its function $F$, its geometry parameter and
its numerical tables. The written proofs below are at the level of a research article. For long computations we point
to the formal proof, which is the basis for the correctness of our theorems.

%% =====================================================================
\section{The construction}\label{sec:construction}

\subsection{The entries of $M$}
Let $0<s<1$ be rational with denominator $q$. Put $u=q(1-s)$, so that $k-1+s=(qk-u)/q$. For the function $F_s$ of
\eqref{eq:frames}, the rational matrix $M$ of \eqref{eq:shift} is as follows. Here $0\le t<t'$ are integers, and
$M[t',t]=M[t,t']$.
\begin{align*}
 M[t,t']&=\frac1{q^2(t'-t)}\sum_{k=t+1}^{t'}\Bigl(\frac1k-\frac q{qk-u}\Bigr),\\
 M[t,t]&=-\frac1{q^2}\sum_{k=1}^{t}\Bigl(\frac1{k^2}-\frac{q^2}{(qk-u)^2}\Bigr).
\end{align*}
For $s=\frac34$ this is $M[t,t']=-\frac1{t'-t}\sum_{k=t+1}^{t'}\frac1{16k(4k-1)}$. For the function $F_\chi$,
\[
 M[t,t']=\frac1{6(t'-t)}\sum_{k=t+1}^{t'}g(k),\qquad M[t,t]=\frac16\sum_{k=1}^{t}g'(k),
\]
with $g(k)=\frac1{3k-2}-\frac1{2k-1}+\frac1{6k}$. Then
$-\frac16g'(k)=\frac1{2(3k-2)^2}-\frac1{3(2k-1)^2}+\frac1{36k^2}$.

We write each function as $F=\sum_{\varsigma\in\mathcal T}w_\varsigma\psi(z+\varsigma)$, where $\mathcal T$ is the set of
shifts of $F$ and $\sum_\varsigma w_\varsigma=0$. For $F_s$, $\mathcal T=\{1,s\}$ with $w_1=1$ and $w_s=-1$. For
$F_\chi$, $\mathcal T=\{\frac13,\frac12,1\}$. The number $|\mathcal T|$ of digamma terms is $2$ for $F_s$ and $3$ for
$F_\chi$.

\begin{lemma}\label{lem:loewner}
For each function of \eqref{eq:frames}, $\omega\,\Lw_F[R,Q]=M[R,Q]+\xi\Id_{R,Q}$ for all finite sets
$R,Q\subset\Z_{\ge0}$, with $\omega$ and $\xi$ as in \eqref{eq:frames}.
\end{lemma}
\begin{proof}
From $\psi(w+1)-\psi(w)=1/w$ one gets $F_s(k)-F_s(k-1)=\frac1k-\frac1{k-1+s}=\frac1k-\frac q{qk-u}$. Telescoping gives the
off-diagonal entries. In the same way $\psi'(w+1)-\psi'(w)=-1/w^2$ gives $F_s'(t)-F_s'(0)$. Finally
$\omega F_s'(0)=q^{-2}(\psi'(1)-\psi'(s))=-q^{-2}\Psi(s,1)=\xi_s$. For $F_\chi$ the entries are computed in the same way.
For the value at $0$, the multiplication formula $\sum_{i<n}\psi'(z+\frac in)=n^2\psi'(nz)$ gives
$\psi'(\frac12)=3\psi'(1)$ and $\psi'(\frac13)+\psi'(\frac23)=8\psi'(1)$. Hence
$F_\chi'(0)=-\frac13\psi'(\frac13)+\frac12\psi'(\frac12)-\frac16\psi'(1)=-\frac16\Psi(\frac13,\frac23)$, and
$-\frac16F_\chi'(0)=\frac1{36}\Psi(\frac13,\frac23)=\xi_\chi$.
\end{proof}

\subsection{Combinations with several terms}\label{ssec:frames}
The planes need combinations with more than two terms, and the second plane needs the trigamma function itself. We
call a \emph{frame} a function
\[
 F(z)=\sum_{s\in S}\bigl(a_s\psi(z+s)+b_s\psi'(z+s)\bigr),
\]
where $S\subset\Q\cap(0,1]$ is a finite set of shifts, the weights $a_s,b_s$ are rational, not all zero, and
$\sum_sa_s=0$. The \emph{order} of a shift is $2$ if $b_s\ne0$, $1$ if $b_s=0\ne a_s$, and $0$ if $a_s=b_s=0$. A
shift of order $0$ carries no pole of $F$. We keep such shifts in $S$, because the planes are families of frames on a
fixed set of shifts, and some members of a family have a weight zero. We write $d$ for the sum of the orders,
the number of poles of $F$ in a period counted with order. The balance $\sum_sa_s=0$ gives $F(z)=O(1/|z|)$ in the right
half-plane. Put $\xi=F'(0)$ and
\[
 \delta(z)=F(z+1)-F(z)=\sum_{s\in S}\Bigl(\frac{a_s}{z+s}-\frac{b_s}{(z+s)^2}\Bigr),
\]
and put $f_t=\sum_{k<t}\delta(k)$ and $e_t=-\sum_{k<t}\delta'(k)$.
Then $f_t=F(t)-F(0)$ and $F'(t)=\xi-e_t$ for integers $t\ge0$. So the L\"owner matrix of $F$ at nonnegative integer
nodes is $\Lw_F=M+\xi\Id$, with $M[t,t']=(f_{t'}-f_t)/(t'-t)$ for $t\ne t'$ and $M[t,t]=-e_t$. The function $F_s$ of
\eqref{eq:frames} is the frame with $S=\{s,1\}$, $a_1=1$, $a_s=-1$, multiplied by $\omega$. The geometries, the members
$P_{r,I}$ and Lemma~\ref{lem:integrality} below apply to every frame without change.

\emph{The frames of the first plane.} Let $S^+=\{\frac13,\frac12,1\}$ and $S^-=\{\frac23,\frac12,1\}$, and write $s_\pm$
for the first shift. For integers $a,b,c$, not all zero, with $a+b+c=0$ put
\[
 F^\pm_{a,b,c}(z)=a\psi(z+s_\pm)+b\psi(z+\tfrac12)+c\psi(z+1).
\]
These are the \emph{frames on $S^\pm$}. All shifts have order at most $1$, so $d\le3$, with $d=3$ when $abc\ne0$. The
function $F_\chi$ of \eqref{eq:frames} is $-\frac16F^+_{2,-3,1}$.

\begin{lemma}\label{lem:planeconst}
For $F^\pm_{a,b,c}$ one has $\xi=\pm\frac{9a}2L+(\frac a2+\frac b3)\pi^2$. If $a\ne0$ and $\tau=(3a+2b)/(27a)$, this
is $\xi=\frac{9a}2(L+\tau\pi^2)$ on $S^+$ and $\xi=-\frac{9a}2(L-\tau\pi^2)$ on $S^-$. For all rationals $a_0$ and
$b_0$, not both zero, there is a frame on $S^+$ or on $S^-$ whose number $\xi$ is a nonzero rational multiple of
$a_0L+b_0\pi^2$.
\end{lemma}
\begin{proof}
On $S^+$, $\xi=a\psi'(\frac13)+b\psi'(\frac12)+c\psi'(1)$. Insert the values of \S\ref{sec:intro} and $c=-a-b$. The
coefficient of $L$ is $\frac{9a}2$ and that of $\pi^2$ is $\frac{2a}3+\frac b2-\frac{a+b}6=\frac a2+\frac b3$. On $S^-$
the sign of $L$ changes. If $a_0=0$, the frame $F^+_{0,3,-3}$ has $\xi=\pi^2$. Its shift $\frac13$ has weight zero. If
$a_0\ne0$, put $t=b_0/a_0$. Let $a=2\den(t)$, $b=a(27t-3)/2$ and
$c=-a-b$. These are integers with $\tau=t$. Here $a\ne0$, $b=0$ only for $t=\frac19$, and $c=0$ only for
$t=\frac1{27}$. For these two values we use $S^-$ with the frames $(1,-3,2)$ and $(1,-2,1)$, for which $\tau=-\frac19$
and $\tau=-\frac1{27}$.
\end{proof}

\emph{The frames of the second plane.} For integers $v$ and $u$, not both zero, put
\[
 F^{(2)}_{v,u}(z)=v\,\psi'(z+1)+u\bigl(\psi(z+1)-\psi(z+\tfrac12)\bigr).
\]
The shift $1$ has order $2$ if $v\ne0$ and order $1$ if $v=0$. The shift $\frac12$ has order $1$ if $u\ne0$ and
order $0$ if $u=0$. So $d=3$ for $uv\ne0$ and $d=2$ otherwise. Since $\psi''(1)=-2\zeta(3)$ and
$\psi'(1)-\psi'(\frac12)=-\frac{\pi^2}3$,
\[
 \xi=-2v\,\zeta(3)-\frac u3\pi^2 .
\]
Given rationals $a$ and $b$, not both zero, choose a rational $\lambda\ne0$ for which $v=\lambda a$ and $u=6\lambda b$
are integers. Then $\xi=-2\lambda(a\zeta(3)+b\pi^2)$, and $u\ne0$ exactly when $b\ne0$.
The frame $F^{(2)}_{1,0}=\psi'(z+1)$ gives $\zeta(3)$. Its increment is $\delta(k)=-(k+1)^{-2}$, so
$M[t,t']=-(t'-t)^{-1}\sum_{k=t+1}^{t'}k^{-2}$ and $M[t,t]=\sum_{k\le t}2k^{-3}$.

Since $\mu$ is invariant under rational affine maps, the theorems of the introduction are statements about the
numbers $\xi$ of these frames and of \eqref{eq:frames}.

\subsection{Geometries}
\begin{definition}\label{def:geometry}
Let $r,m\ge0$. The \emph{geometry} $\mathcal G(r,m)$ consists of the common nodes $C=[0,r)$, the row nodes
$X'=\{r+2i:0\le i<m\}$ and the column nodes $Y'=\{r+2i+1:0\le i<m\}$. Rows are indexed by $R=C\sqcup X'$ and columns by
$Q=C\sqcup Y'$. Let $I\subseteq C$ be a set of deleted common nodes. The matrix
\[
 \Phi_{\mathcal G,I}(x)=\bigl(M[t,t']+x\,\mathbf 1[t=t'\in C\setminus I]\bigr)_{t\in R\setminus I,\ t'\in Q\setminus I}
\]
is square of size $|C\setminus I|+m$. Put $P_{\mathcal G,I}(x)=\det\Phi_{\mathcal G,I}(x)\in\Q[x]$. We write
$\Phi_{\mathcal G}$ and $P_{\mathcal G}$ when $I=\emptyset$, and we write $\bar N=r+2m$ for the number of nodes.
\end{definition}

By Lemma~\ref{lem:loewner}, $\Phi_{\mathcal G}(\xi+h)=\omega\,\Lw_F[R,Q]+h\Id_{R,Q}$ for real $h$. The geometry with
$m=0$ is the square L\"owner matrix of size $r$.

\begin{definition}\label{def:altgeom}
Let $0<c\le1$ and $\alpha=c/2$. The \emph{alternating geometry} of parameter $c$ at an even scale $r$ is
$\mathcal G(r,m(r))$ with $m(r)=\lceil\alpha r-\frac12\rceil$. Then $|m(r)-\alpha r|\le\frac12$, and the border
$[r,r+2m)$ has length $cr+O(1)$. For a depth $K\ge0$, its \emph{members} are the matrices
$\Phi_{r,I}=\Phi_{\mathcal G,I}$ with $|I|\le K$. We put $P_{r,I}=\det\Phi_{r,I}$.
\end{definition}

For the functions $F_s$ we use $c=\frac3{10},\frac15,\frac3{10},\frac29,\frac16$ for $q=2,4,3,5,7$. For the frames of the first plane we use $c=\frac16$, for those of the second plane $c=\frac6{41}$, and for $\zeta(3)$ $c=\frac3{10}$. For $F_\chi$ we use
$c=\frac12$, as explained in Section~\ref{sec:chi3-quarter}. All $o(1)$ terms below are uniform over the members as
$r\to\infty$.

\begin{lemma}[integrality]\label{lem:integrality}
$P_{\mathcal G,I}$ has degree at most $|C\setminus I|$. If $a$ and $b\ge1$ are integers with
$P_{\mathcal G,I}(a/b)\ne0$, then $|\den(P_{\mathcal G,I})\,P_{\mathcal G,I}(a/b)|\ge b^{-|C\setminus I|}$.
\end{lemma}
\begin{proof}
The variable occurs only in the $|C\setminus I|$ diagonal entries indexed by $C\setminus I$. So each Leibniz term has
degree at most $|C\setminus I|$. Then $b^{|C\setminus I|}\den(P_{\mathcal G,I})\,P_{\mathcal G,I}(a/b)$ is a nonzero
integer.
\end{proof}

\emph{Example.} Take $s=\frac34$, $r=2$ and $m=1$, so that $C=\{0,1\}$, $X'=\{2\}$ and $Y'=\{3\}$. Then
\[
 \Phi_{\mathcal G}(x)=\begin{pmatrix}
 x & -\frac1{48} & -\frac{67}{7392}\\[2pt]
 -\frac1{48} & x+\frac7{144} & -\frac{47}{14784}\\[2pt]
 -\frac{17}{1344} & -\frac1{224} & -\frac1{528}
 \end{pmatrix}
\]
and
\[
 P_{\mathcal G}(x)=-\frac{x^2}{528}-\frac{823\,x}{3725568}-\frac{239}{37158912}.
\]
The variable occurs only at the common nodes $0$ and $1$, so the degree is $|C|=2$ although the matrix has size $3$. When $M[X',Y']$ is invertible, a Schur complement shows that $P_{\mathcal G}$
is, up to a constant factor, the characteristic polynomial of a rational matrix of size $|C|$, as in the square case.
We do not use this factorisation.

%% =====================================================================
\section{Arithmetic: the denominators}\label{sec:content}

In this section we bound the common denominator of the coefficients of the members $P_{r,I}$ for the functions $F_s$.
The method is that of Rhin and Viola \cite{RhinViola96,RhinViola01}: the $p$-adic valuation of the denominator is bounded
prime by prime, and the bounds are summed against $\log p$ with the prime number theorem in arithmetic progressions.
What is new is the form of the local bounds. Each is an explicit integer expression in four integers attached to the
prime, and their minimum, integrated, gives the denominator exponent.

The local bounds come from four mechanisms, of increasing depth.
\begin{enumerate}[label=(\alph*)]
\item \emph{Valuations of the entries.} The assignment bounds (T1)--(T5) of \S\ref{sssec:trop} use only the
 valuations of single entries and LP duality.
\item \emph{Corank modulo $p$.} The bounds (LP1), (LP2) and (R) of \S\S\ref{sssec:lp} and~\ref{sssec:rank} compute the
 corank of the reduction of a rescaled matrix.
\item \emph{Second order.} The bounds (W) and (S) of \S\S\ref{sssec:window} and~\ref{ssec:sep} use the pairing of kernel
 vectors modulo $p^2$ (\S\ref{ssec:second}).
\item \emph{Third order.} The bounds (U$_n$) of \S\ref{sssec:wrap} treat the primes smaller than the border.
\end{enumerate}
Comparing the local bounds with a table of affine pieces is a finite computation with rational numbers, and
Theorem~\ref{thm:content} turns an admissible table into the arithmetic estimate. Section~\ref{sec:chi3-quarter} treats the
function $F_\chi$.

\subsection{The root class of a prime}
Fix the function $F_s$, with $s$ of denominator $q$, and the alternating geometry of parameter $c=2\alpha$.
Let $p>q$ be a prime. The indices $k\ge1$ at which the entries of $M$ acquire a denominator divisible by $p$ are the
multiples of $p$ and the solutions of $qk\equiv u\pmod p$. Let $k_p$ be the least of the latter. Then
\[
 qk_p=jp+u,\qquad 0<j<q,\qquad j\equiv qs\,p^{-1}\pmod q .
\]
We call $j$ the \emph{root class} of $p$ and put $\kappa=\kappa_j=j/q$. Then $k_p-\kappa p=u/q\in[\frac1q,1-\frac1q]$. So
the first pole is at distance less than one from $\kappa p$, uniformly in $s$. Throughout we assume $q(r+2m)<p^2$. Then
every pole below $\bar N$ is simple, and only the first power of $p$ matters. The primes with $p^2\le q(r+2m)$ and the prime
powers contribute $O(r^{3/2}\log r)$ in total.

\subsection{Gauss valuations, the assignment bound and the rank of the reduction}
Fix a prime $p$. Let $\cO\subset\Q(x)$ be the localisation of $\Z_{(p)}[x]$ at the prime ideal $(p)$. It is a discrete
valuation ring with uniformiser $p$ and residue field $\F_p(x)$. Its valuation $v$, restricted to $\Q[x]$, is the
\emph{Gauss valuation} $v(f)=\min_k v_p(f_k)$ for $f=\sum f_kx^k$. By Gauss's lemma $v(fg)=v(f)+v(g)$, and
$v_p(\den f)=\max(0,-v(f))$. For a member we write $e_p(r,I)=v_p(\den P_{r,I})$.

\begin{lemma}[assignment bound and LP duality]\label{lem:trop}
Let $A$ be a square matrix over $\Q[x]$ with rows $X_0$ and columns $Y_0$. Let $W_{tt'}=\max(0,-v(A_{tt'}))$, and let
$\tau(W)$ be the largest weight $\sum W_{tt'}$ of a perfect matching between $X_0$ and $Y_0$. Then
$v(\det A)\ge-\tau(W)$. If real numbers $u_t$, $u'_{t'}$ satisfy $W_{tt'}\le u_t+u'_{t'}$ for all $t\in X_0$ and
$t'\in Y_0$, then $\tau(W)\le\sum_tu_t+\sum_{t'}u'_{t'}$.
\end{lemma}
\begin{proof}
Each Leibniz term is a product of entries along a perfect matching. So its valuation is at least minus the weight of the
matching. The valuation of a sum is at least the minimum of the valuations of its terms. The second statement is weak
LP duality for the assignment problem. The weight of a perfect matching is at most $\sum_tu_t+\sum_{t'}u'_{t'}$.
\end{proof}

\begin{lemma}[rank of the reduction]\label{lem:rank}
Let $(u,u')$ be integer potentials with $W_{tt'}\le u_t+u'_{t'}$ for the matrix $\Phi_{r,I}$ at $p$, of value
$\tau'=\sum_tu_t+\sum_{t'}u'_{t'}$. Then $\widetilde\Phi=\operatorname{diag}(p^{u})\,\Phi_{r,I}\,
\operatorname{diag}(p^{u'})$ has entries in $\cO$, and
\[
 e_p(r,I)\le\tau'-\corank_{\F_p(x)}(\widetilde\Phi\bmod p).
\]
\end{lemma}
\begin{proof}
We have $v(\det\Phi_{r,I})=-\tau'+v(\det\widetilde\Phi)$. Over the discrete valuation ring $\cO$ the Smith normal form
shows $v(\det\widetilde\Phi)\ge\corank(\widetilde\Phi\bmod p)$.
\end{proof}

\subsection{Second-order information}\label{ssec:second}
The rank bounds below use a refinement of Lemma~\ref{lem:rank} at second order. Suppose that, after a unimodular change
of rows and columns over $\cO$, the first $n_1$ rows and the first $n_2$ columns of $\widetilde\Phi$ are divisible by
$p$. Suppose also that the $n_1\times n_2$ block where they meet is divisible by $p^2$. Dividing the first $n_1$ rows by
$p$ leaves the first $n_2$ columns divisible by $p$. So
\begin{equation}\label{eq:corner}
 v(\det\widetilde\Phi)\ge n_1+n_2 .
\end{equation}

The vectors that produce such rows and columns are kernel vectors of the reduction of $\widetilde\Phi$ on single residue
classes. The condition on the corner is a statement about their pairing at first order, which we now describe. Let
$\widetilde\Phi=p^2\Phi_{r,I}$, so that all potentials equal $1$. Since $q(r+2m)<p^2$, two distinct nodes $t\ne t'$ have
$v_p(t'-t)\ge1$ only if $t\equiv t'\pmod p$, and then $v_p(t'-t)=1$. Hence $\widetilde\Phi$ has entries in $\cO$. Its
reduction is block diagonal over the residue classes modulo $p$, which we call the \emph{fibres}. For nodes $t$, $t'$
in distinct fibres,
\begin{equation}\label{eq:firstorder}
 \widetilde\Phi_{tt'}\equiv\frac{p\,(V(t')-V(t))}{t'-t}\pmod{p^2}.
\end{equation}
Here $V(t)\in\F_p$ is the residue of $p$ times the partial sum $\sum_{k\le t}(\frac1k-\frac q{qk-u})$, up to the
normalising factor $q^{-2}$. The function $V$ is a step function of $t$ which jumps only at the poles $t\equiv0$ and
$t\equiv k_p$. Since $t'-t$ is prime to $p$, only its residue matters in \eqref{eq:firstorder}.

For a vector $\mathbf y$ indexed by nodes, its \emph{moment pair} is
\[
 \hat{\mathbf y}=\Bigl(\sum_ty_t,\ \sum_ty_tV(t)\Bigr).
\]

\begin{lemma}[pairing]\label{lem:pairing}
Let $\mathbf y$ be supported on the rows of one fibre and $\mathbf y'$ on the columns of a different fibre. Let $t$ and
$t'$ be nodes of the first and of the second fibre. Then
\[
 \mathbf y^{\top}\widetilde\Phi\,\mathbf y'\equiv\frac p{t'-t}\,
 \det\begin{pmatrix}\hat{\mathbf y}\\ \hat{\mathbf y}'\end{pmatrix}\pmod{p^2}.
\]
In particular the pairing vanishes modulo $p^2$ whenever the moment pairs of $\mathbf y$ and $\mathbf y'$ are
proportional.
\end{lemma}
\begin{proof}
Expand $\mathbf y^{\top}\widetilde\Phi\,\mathbf y'$ with \eqref{eq:firstorder}.
\end{proof}

The moment pairs are never both zero. This follows from the displacement structure of L\"owner matrices.

\begin{lemma}[moment lemma]\label{lem:moment}
Let $A$ be a matrix with rows indexed by a set $X_0$ of distinct integers and columns by a set $Y_0$, with
$|X_0|=|Y_0|+1$. Suppose that its off-diagonal entries are $A(t,t')=(\mathcal L(t')-\mathcal L(t))/(t'-t)$ for some
function $\mathcal L$, and that its diagonal entries, on $X_0\cap Y_0$, are arbitrary. Then $A$ has a nonzero left kernel
vector $\mathbf y$ with $(\sum_ty_t,\sum_ty_t\mathcal L(t))\ne(0,0)$.
\end{lemma}
\begin{proof}
Write $\operatorname{diag}(X_0)$ and $\operatorname{diag}(Y_0)$ for the diagonal matrices of the indices. Then
$\operatorname{diag}(X_0)A-A\operatorname{diag}(Y_0)=\mathcal L_{X_0}\mathbf 1^{\top}-\mathbf 1\mathcal L_{Y_0}^{\top}$,
because the diagonal entries cancel on common indices. The left kernel $\mathcal V$ of $A$ is nonzero. Suppose that
every $\mathbf y\in\mathcal V$ had zero moment pair. Multiplying on the left by $\mathbf y^{\top}$ would give
$\mathbf y^{\top}\operatorname{diag}(X_0)A=0$. So $\operatorname{diag}(X_0)$ would preserve $\mathcal V$. As this matrix
is diagonal with distinct entries, $\mathcal V$ would contain a standard basis vector. That vector is supported on one
index $t$, and its moment pair is $(1,\mathcal L(t))\ne(0,0)$. This is a contradiction.
\end{proof}

\subsection{The local bounds}
We now list the local bounds. Each is a statement of the following form. Let $P_{r,I}$ be a member with $|I|\le K$, and
let $p>p_0$ be a prime with $q(r+2m)<p^2$ that satisfies the stated conditions. Then
\begin{equation}\label{eq:localform}
 e_p(r,I)\le\tau(r,m,p,k_p)+\epsilon_\tau\,|I| .
\end{equation}
Here $\tau$ is the expression in the table and $\epsilon_\tau\in\{0,1\}$ is given in its last column. The
constant $p_0$ depends only on $F$. We write $(\cdot)_+=\max(\cdot,0)$, $\|m\|_p$ for the distance of $m$ to
$p\Z$, and $\bar u=(k_p+p-r)_+$.

\begin{theorem}[local bounds]\label{thm:local}
The following bounds of the form \eqref{eq:localform} hold.
\begin{center}
{\renewcommand{\arraystretch}{1.2}
\begin{tabular}{@{}lp{0.47\textwidth}p{0.27\textwidth}c@{}}
\toprule
 & bound $\tau$ & conditions & $\epsilon_\tau$\\
\midrule
(T1) & $2(r+m)$ & none & $0$\\
(T1$'$) & $2r+2m-\|m\|_p$ & none & $0$\\
(T2) & $2(r-\bar u)_++2m+\min(\bar u,2m)$ & $2m<p\le r$ & $0$\\
(T3) & $2r+m$ & $2m<p\le r$ & $0$\\
(T4) & $2(r-k_p)_++2m+(2m-(p-r))_+$ & $r<p$ & $0$\\
(T5) & $(2m-(k_p-r)_+)_+$ & $\bar N\le p$, $r\le k_p$ & $0$\\
(LP1) & $\tau_{\rm T2}-(\min(\bar u,2m)-m)_+$ & $p\le r<p+k_p$, $\bar N\le2p$ & $1$\\
(LP2) & $2(r-k_p)+3m$ & $r<p<r+m$, $k_p\le r$, $\bar N\le2p$ & $1$\\
(W) & $2r+2m-(t_2-t_1)$ for each window $[t_1,t_2)$ of Theorem~\ref{thm:window} & as there & $0$\\
(R) & $\Pi_R$ of \eqref{eq:PiR} & $2m<p$ & $0$\\
(S) & $\Pi_S$ of \eqref{eq:PiS} & $2m\le p$ & $1$\\
(U$_n$) & $\Pi_{U,n}$ of \eqref{eq:PiU}, $n\ge0$ & $np\le m<(n+1)p$, and $p\le2m$ if $n=0$ & $0$\\
\bottomrule
\end{tabular}}
\end{center}
\end{theorem}

The rest of this subsection proves Theorem~\ref{thm:local}, family by family.

\subsubsection{The assignment bounds (T1)--(T5)}\label{sssec:trop}
These follow from Lemma~\ref{lem:trop} with explicit potentials. Consider an off-diagonal entry $M[t,t']$. By the
ultrametric inequality, $-v_p(M[t,t'])$ is at most the number of the following two events that occur: the interval
$(t,t']$ contains a pole, and $p$ divides $t'-t$. A diagonal entry $M[t,t]$ has $-v_p\le2$ if $t\ge k_p$ and $-v_p\le0$
if $t<k_p$. The bound (T1) is the trivial bound, with potentials $1$ on every row and column.

For (T1$'$) the potentials depend on the residue class modulo $p$. On the rows of a class they are $\frac12$,
$\frac32$ or $1$, according as $X'$ has more, fewer or as many nodes in the class as $Y'$. On the columns they are
$\frac32$, $\frac12$ or $1$ in the same three cases. The border alternates between $X'$ and $Y'$. So the imbalance is
at most one per class, and the number of unbalanced classes is $2\|m\|_p$.

For (T2)--(T5) one starts from the potential $\mathbf 1[t\ge k_p]$ on the common nodes $t$, which accounts for the diagonal.
One adds $1$, or $1+\mathbf 1[t\ge p]$, on the border nodes $t$. Then one raises the potentials of the common nodes that
have a congruent partner in the border. In each range one checks $W\le u+u'$ entrywise. Two facts are used. Two distinct
nodes not congruent modulo $p$ have weight at most $1$. Two border nodes are never congruent once $2m<p$. For (T5) we
have $\bar N\le p$ and $r\le k_p$. Then no common node accounts for the diagonal, and only the border nodes beyond $k_p$ carry
potential $1$. So (T5) vanishes once $k_p\ge \bar N$, which happens for $p/r>\max(1+2\alpha,1/\kappa)+o(1)$. All cases are
checked in the formal proof.

\subsubsection{Primes near $r$: the bounds (LP1) and (LP2)}\label{sssec:lp}
For $p$ close to $r$ the reduction of the rescaled matrix has a large block of rank one per residue class. Its corank
can be computed. We give the proof of (LP2). The bound (LP1) applies to the primes with $p\le r<p+k_p$ and $\bar N\le2p$. It
is proved by the same method. With the potentials of (T2), eliminating an invertible block of common nodes brings the
reduction to the form $\begin{psmallmatrix}\tilde x\Id&\mathcal J_Y\\ \mathcal J_X&0\end{psmallmatrix}$, where $\tilde x$
is the variable $x$ plus a constant. In this form the border nodes below $p+k_p$ contribute rows and columns
proportional to unit vectors. The Frobenius--K\"onig theorem bounds its rank. The saving over (T2) is
$(\min(2m,p+k_p-r)-m)_+$, where $\min(2m,p+k_p-r)$ is the number of border nodes below $p+k_p$. We refer to the formal
proof for the details.

\begin{theorem}\label{thm:lp2}
Assume $r<p<r+m$, $k=k_p\le r$ and $\bar N=r+2m\le2p$, and that every pole below $\bar N$ is simple. Then
$v(\det\Phi_{r,\emptyset})\ge-(2(r-k)+3m)$. The same holds for every member, with $\epsilon_\tau=1$.
\end{theorem}
\begin{proof}
Use the potentials $u(t)=\mathbf 1[k\le t]$ for common nodes $t<r$ and $u(t)=1+\mathbf 1[p\le t]$ for border nodes $t$.
Their value is $\tau'=2(r-k)+2m+(\bar N-p)$. Since $\bar N\le2p$, the poles below $\bar N$ are $k$, $p$, and $p+k$ if $p+k<\bar N$. Each is
simple. The residues of $p$ times the increment $\frac1t-\frac q{qt-u}$ at these poles are $-q/j$, $1$ and $-q/(q+j)$.

\emph{Rows and columns beyond $p+k$ vanish.} Let $t\in[p+k,\bar N)$ be a border node, so that $u(t)=2$. Let $t'$ be a node of
the other kind. If $t'<k$, then $u(t')=0$. The partial sum over $(t',t]$ then has valuation at least $-1$, and
$t-t'\in(p,2p)$ is prime to $p$. Otherwise $u(t')\ge1$ and $v_p(M[t',t])\ge-2$. In both cases the rescaled entry is
divisible by $p$.

\emph{The remaining entries.} For nodes below $p+k$ the only poles are $k$ and $p$. Put $C_0=[0,k)$ and $C_1=[k,r)$.
Recall that $\omega=q^{-2}$ is the normalisation of \eqref{eq:frames}. The reduction has the following blocks. On
$C_0\times C_0$ it is the regular part plus $x$ on the diagonal. On $C_0\times C_1$ the entries are
$-\omega q/(j(t-t'))$ with $t'\in C_0$ and $t\in C_1$. On $C_1\times C_1$ it is the scalar matrix $\omega(q/j)^2\Id$. A
border node $t\in[r,p)$ against $t'\in C_0$ gives $\omega q/(j(t-t'))$. A border node $t\in[p,p+k)$ against $t'\in C_0$
gives $\omega(1-q/j)$ if $t'=t-p$, and $0$ otherwise. All other border entries vanish.

\emph{Rank.} Delete the zero rows and columns found above. This does not change the rank. Eliminate the scalar block on
$C_1$. The Schur complement on $C_0$ is $\tilde x\Id$ plus a constant multiple of the matrix with entries
$\sum_{i=1}^{p-r}\frac1{(t+i)(t'+i)}$ for $t,t'\in C_0$. Here $\tilde x$ is $x$ plus a constant, and the computation uses
the reflection $H_{p-1-n}\equiv H_n\pmod p$ of harmonic numbers. The border nodes $t\in[r,p)$ are indexed by $i=p-t$.
They are all present because $\bar N>p$. Even offsets are rows and odd offsets are columns, and together they cover
$1\le i\le p-r$. So row and column operations remove that multiple and leave
$\begin{psmallmatrix}\tilde x\Id&\mathcal J_Y\\ \mathcal J_X&0\end{psmallmatrix}$. Here it is essential that $x$ is an
indeterminate, so that $\tilde x$ is a unit. A Schur complement on $\tilde x\Id$ then shows that the rank is
$k+\rank(\mathcal J_X\mathcal J_Y)$. In $\mathcal J_X\mathcal J_Y$ the block between the high rows and the high columns is
zero, since their offsets $t-p$ have opposite parities. So $\rank(\mathcal J_X\mathcal J_Y)\le p-r$. In total the rank of
the reduction is at most $(r-k)+k+(p-r)=p$, and the corank is at least $r+m-p$. Lemma~\ref{lem:rank} gives
$v(\det)\ge(r+m-p)-\tau'=-(2(r-k)+3m)$. Deleting a common node changes the potentials by at most one.
\end{proof}

Under the hypothesis $\bar N\le p+k$ the third pole does not occur, and the proof is shorter. The border nodes beyond $p+k$
only contribute rows and columns that vanish modulo $p$. This observation allows the whole range $\bar N\le2p$, which is
automatic on the interval $r<p<r+m$.

\subsubsection{Pole-free windows: the bound (W)}\label{sssec:window}
Like the rank bounds, the bound (W) does not come from the valuations of the entries alone. Between two consecutive
poles the step function $V$ of \eqref{eq:firstorder} is constant. The reductions of the fibres of nearby residue classes
are then translates of one another. Kernel vectors can be transported from one fibre to the next. Their pairings vanish
modulo $p^2$ by Lemma~\ref{lem:pairing}. The corner bound \eqref{eq:corner} turns this into a saving at second order.

A \emph{cut} is an integer $t$ with $t\equiv0$, $t\equiv k_p$ or $t\equiv r\pmod p$. A window $[t_1,t_2)$ of integers is
\emph{cut-free} if no cut lies in the open interval $(t_1,t_2)$.

\begin{theorem}\label{thm:window}
Assume that every pole below $\bar N$ is simple and $\bar N<p^2$. Let $n\ge0$ with $np\le m$, and put $\bar N_n=r+2(m-np)$. Let
$[t_1,t_2)$ be a cut-free window with $\max(r,\bar N_n-p)\le t_1\le t_2\le\min(\bar N_n,r+p)$. Then every member satisfies
\[
 e_p(r,I)\le2r+2m-(t_2-t_1).
\]
\end{theorem}
For $n=0$ the windows lie in the border above $\max(r,\bar N-p)$, and the bound applies to primes $p>m$. For
$n=\lfloor m/p\rfloor\ge1$ it applies to the small primes $p\le m$. For these primes only (T1) and (T1$'$) are
available among the assignment bounds. Read modulo $p$, the window range is then the arc of residues that receive an odd
number of border nodes. The saving over (T1$'$) is the amount, if any, by which the length of the longest cut-free run
of this arc exceeds $\|m\|_p$. For other $n$ the window range is empty. The proof uses $F$ only through its poles.
For a function with several digamma terms the same statement holds when the cuts are the residues of all the poles and
of $r$.

\begin{proof}
Let $\widetilde\Phi=p^2\Phi$, with independent diagonal variables on the common nodes. Every member is a coefficient of
this determinant, so a bound for it bounds every member. Let $i<i'$ be two positions of the window. They lie in
different residue classes. The translation $t\mapsto t+(i'-i)$ maps the nodes of the fibre of $i$ bijectively onto those
of the fibre of $i'$. It maps common nodes to common nodes, since no $t\equiv r$ lies in $(i,i']$. It maps border nodes
to border nodes, since both classes lie in the window range. It maps $X'$ to $Y'$ when $i'-i$ is odd. There is no pole
in $(i,i']$. So the step function $V$ and the diagonal residues are unchanged, and so are the differences of nodes.
Hence the reduced block of the fibre of $i'$ equals that of $i$ when $i'-i$ is even, and its transpose when $i'-i$ is
odd.

In the window the fibres alternate between one more row than column and one more column than row. Take a nonzero left
kernel vector $\mathbf y$ of one block and transport it to every fibre of the window. It is a left kernel vector on the
fibres with a surplus row and a right kernel vector on the others. Its moment pair is the same everywhere. By
Lemma~\ref{lem:pairing} all pairings between these vectors vanish modulo $p^2$. There are $n_1$ such vectors on rows and
$n_2$ on columns, with $n_1+n_2=t_2-t_1$. Extending them to unimodular matrices, \eqref{eq:corner} gives
$v(\det\widetilde\Phi)\ge t_2-t_1$. Since $\det\widetilde\Phi=p^{2(r+m)}\det\Phi$, the theorem follows.
\end{proof}

\subsubsection{The rank bound (R)}\label{sssec:rank}
Let $2m<p$. Write $r=k_0p+r_p$ with $0\le r_p<p$. Note that $k_p=\lceil jp/q\rceil$, since $0<u<q$. Put
$\mathrm{arc}_0=[0,k_p-1]$ and $\mathrm{arc}_1=[k_p,p-1]$. For $k\ge0$ and $\iota\in\{0,1\}$ let
\[
\begin{gathered}
 n_{k,\iota}=\bigl|[r,r+2m)\cap[kp+\min\mathrm{arc}_\iota,\,kp+\max\mathrm{arc}_\iota]\bigr|,\\
 n_{\rm sq}=\bigl|[\max(0,r+2m-p),\min(r,k_p))\bigr|,
\end{gathered}
\]
and
\begin{multline}\label{eq:PiR}
 \Pi_R=2(r+m)\\-\max\bigl(m+n_{\rm sq},\ n_{0,0}+n_{1,0}+2n_{\rm sq},\ n_{k_0,0},\ n_{k_0,1},\ n_{k_0+1,0},\ n_{k_0+1,1}\bigr).
\end{multline}

\begin{proof}[Proof of the bound (R)]
Every node is $t=i+pk$ with residue $i\in[0,p)$ and digit $k$. Write $k(i)$ for the number of common nodes in the fibre
of $i$. The border has length $2m<p$. So this fibre contains the common digits $0,\dots,k(i)-1$ and at most one border
digit, namely $k(i)\in\{k_0,k_0+1\}$. This border node is a row or a column according to the parity of $i+pk(i)-r$. The
step function $V$ is constant on the fibre up to the bit $\iota(i)=\mathbf 1[i\ge k_p]$. Thus the reduced block of a fibre
depends only on $(k(i),\iota(i))$ and on the role of its border digit. Its entries have the L\"owner form of
Lemma~\ref{lem:moment}.

A fibre with a border row has a left kernel vector with nonzero moment pair. A fibre with a border column has such a
right kernel vector. For the digit $0$ below the pole the step function vanishes. So these fibres have moment pairs
proportional to $(1,0)$. A fibre of type $(1,0)$ without border digit has the zero block and contributes one row and one
column. Choosing one such vector per fibre annihilates $m+n_{\rm sq}$ rows and as many columns. This gives the saving
$m+n_{\rm sq}$.

Within a class of fibres of the same type $(k,\iota)$ all moment pairs are equal. So by Lemma~\ref{lem:pairing} the
pairings between the annihilated rows and columns of one class vanish modulo $p^2$. Then \eqref{eq:corner} gives a saving
equal to the number of annihilated coordinates of the class. This number is $n_{k,\iota}$ for $k\ge1$ and
$(k,\iota)\ne(1,0)$. It is $n_{0,0}+n_{1,0}+2n_{\rm sq}$ for the class with moment pair $(1,0)$. Taking the best class gives
\eqref{eq:PiR}. A deleted common node changes at most one fibre. Give up that fibre, and one further fibre to restore the
balance of rows and columns. The trivial certificate then costs at most two annihilated coordinates of each kind. The smaller
matrix compensates for them, so $\epsilon_\tau=0$.
\end{proof}

\subsubsection{Separated border nodes: the bound (S)}\label{ssec:sep}
The bound (S) uses a different source of cancellation. Choose a window $[t_0,t_v]$ of border nodes whose residues lie in
one arc between consecutive poles. Consider a border node outside it. Suppose the interval between the node and $t_0$
contains a singular index. Then the residues modulo $p$ of $p$ times the partial sums at the two nodes differ, because
two singular contributions in an interval of length less than $p$ cannot cancel. The leading coefficient in $x$ of the
bordered determinant of a fibre is this residue. So the cross product of the two bordered fibre determinants has a
nonzero leading coefficient. The border node is then eliminated against $t_0$ with no further information. Count the
residues of the window together with the residues of the eliminated border nodes. This gives an explicit bound, a
minimum over the four window types $(k,\iota)\in\{k_0,k_0+1\}\times\{0,1\}$, which we now state.

Let $2m\le p$, and let $k_0$, $r_p$, $\mathrm{arc}_0$, $\mathrm{arc}_1$ and $n_{k,\iota}$ be as in \S\ref{sssec:rank}.
For a type $(k,\iota)$ let $\mathcal K_{k,\iota}$ be the interval of the integers $kp+i$ with $i\in\mathrm{arc}_\iota$,
and with $i<r_p$ if $k=k_0+1$. The border nodes in $\mathcal K_{k,\iota}$ form the window, and there are $n_{k,\iota}$ of
them. Let $n_K$ be the number of border nodes in $\mathcal K_{k,\iota}$ and in its two translates
$\mathcal K_{k,\iota}\pm p$. Put
\[
 n_\emptyset=p-|\mathcal K_{k,\iota}|-(2m-n_K),\qquad n_+=(n_{k,\iota}-m)_+ .
\]
Since $2m\le p$, each residue class receives at most one border node. So $n_K$ is the number of border nodes whose
residues are those of $\mathcal K_{k,\iota}$. And $n_\emptyset$ is the number of the other residue classes that receive
no border node. Then
\begin{equation}\label{eq:PiS}
 \Pi_S=\min_{(k,\iota)}\bigl(2r+m-n_+-(n_+-n_\emptyset)_+\bigr),
\end{equation}
where a type whose window is empty contributes $2r+m$. The bound (S) is $e_p(r,I)\le\Pi_S+|I|$. Its proof carries out the
elimination above for every border node outside the window. We refer to the formal proof for it.

\subsubsection{Primes smaller than the border: the bound (U$_n$)}\label{sssec:wrap}
When $p\le2m$ the border meets some residue classes more than once, and the bound (S) is not available. The bound (U$_n$)
extends it to these primes. It works at third order in the Smith normal form. We show two things. After the elimination
of a block of unit pivots, the rescaled matrix has the Smith invariants of a smaller geometry, to which the separation
argument of (S) applies. The third-order information of that argument survives the elimination.

Write $m=np+\bar m$ with $0\le\bar m<p$, so that $n=\lfloor m/p\rfloor$. Put
\[
 d=\|m\|_p=\min(\bar m,p-\bar m),\qquad \bar r=r+m-d .
\]
For $i\in\{0,1,2\}$ and $\iota\in\{0,1\}$ let $Z_i=p\lfloor r/p\rfloor+(n+i)p$, and put
\[
 w_0=Z_i+\min\mathrm{arc}_\iota,\quad t_0=\max(\bar r,w_0),\quad t_v=\min(Z_i+\max\mathrm{arc}_\iota,\,\bar r+p-1).
\]
We call the geometry with common block $[0,\bar r)$ and border $[\bar r,\bar r+2d)$ the \emph{auxiliary geometry}
$(\bar r,d)$. In it, let
$n_+=\bigl(\min(t_v+1,\bar r+2d)-t_0-d\bigr)_+$. Let $n_K$ be the number of border nodes in $[w_0,t_v]$ and in its two
translates by $\pm p$, and let $n_\emptyset=p-(t_v+1-w_0)-(2d-n_K)$. Then
\begin{equation}\label{eq:PiU}
 \Pi_{U,n}=\min_{i,\iota}\Bigl(2(r+m)-d-n_+-(n_+-n_\emptyset)_+\Bigr),
\end{equation}
all subtractions being truncated. This is the expression \eqref{eq:PiS} for the auxiliary geometry $(\bar r,d)$, with the
windows $[w_0,t_v]$ in place of $\mathcal K_{k,\iota}$. Note that $2(r+m)-d=2\bar r+d$.

\begin{theorem}\label{thm:wrap}
Let $p>p_0$ be a prime with $q(r+2m)<p^2$, and let $n\ge0$ with $np\le m<(n+1)p$. Assume $p\le2m$ if $n=0$. Then every
member satisfies $e_p(r,I)\le\Pi_{U,n}$.
\end{theorem}

We first describe the proof for $n=0$, that is $m<p\le2m$. We then give the changes for $n\ge1$. As in the proof of
Theorem~\ref{thm:window} we work with $\widetilde\Phi=p^2\Phi$. We put an independent indeterminate $x_t$ on each common
diagonal entry. We work over the discrete valuation ring $\cO$ of Gauss valuations in these indeterminates. Every member
is a coefficient of $\det\widetilde\Phi$, up to the deletion of common rows and columns, which we treat at the end.

\emph{Pivots from pairs of border nodes.} For $n=0$ there are $2m-p$ residue classes that carry two border nodes $t<t+p$, one a row and the
other a column. There are $2d$ classes that carry one border node, and they lie in an interval of length $2d<p$.
Eliminate the common block. Its fibres are invertible because the $x_t$ are indeterminates. Consider a class with two
border nodes. Let $A$ be its common fibre, let $\mathbf a$ and $\mathbf a'$ be the couplings of its border column and
border row to $A$, and let $\varkappa$ be its border entry. After the elimination the class has the pivot
$\varpi=\varkappa-\mathbf a'^{\top}A^{-1}\mathbf a$. As the $x_t$ tend to infinity, $\varpi$ tends to a nonzero multiple
of the residue of $p$ times the increment of $F$ over $(t,t+p]$. This interval contains exactly one pole of each
digamma term. Let $k'p$ be the multiple of $p$ in it, and $t'$ the pole of the second term. If the two contributions
cancelled, then $q(t'-k'p)\equiv u$ with $|q(t'-k'p)-u|<q$. Hence $q(t'-k'p)=u$, which is impossible for $0<u<q$. So
these $2m-p$ pivots are units.

After them, the matrix of size $r+m=\bar r+d$ is equivalent over $\cO$ to $\operatorname{diag}(\Id_{\bar r},p\,\Xi)$,
with $\Xi$ of size $d$. Its Schur complement on the border has the expansion
\begin{equation}\label{eq:schur3}
 S=p\,S_1+p^2\bigl(S_2-\mathcal C_{UB}\mathcal B^{-1}\mathcal C_{BU}\bigr)+O(p^3).
\end{equation}
Here $\mathcal B$ is the diagonal of these pivots. The matrices $\mathcal C_{UB}$ and $\mathcal C_{BU}$ are the
first-order couplings between these pivots and the single border nodes. The correction term in \eqref{eq:schur3} is not
negligible. The following identities account for it.

\begin{lemma}[bordered inverses]\label{lem:bordered}
Let $A$ be an invertible square matrix over a commutative ring, and let $\mathbf u$ be a vector.
\begin{enumerate}[label=(\alph*),leftmargin=*]
\item Let $A'$ be obtained from $A$ by replacing one column by a vector $\mathbf w$, and suppose that $A'$ is invertible.
 Let $\mathbf u'$ be obtained from $\mathbf u$ by replacing the corresponding entry by $\mathbf u^\top A^{-1}\mathbf w$.
 Then $\mathbf u'^{\top}A'^{-1}\mathbf u=\mathbf u^\top A^{-1}\mathbf u$.
\item Suppose that the matrix
 \[
  A'=\begin{pmatrix}A&\mathbf a\\ \mathbf a'^\top&\varkappa\end{pmatrix}
 \]
 is invertible. Then
 $(\mathbf u,y)^\top A'^{-1}(\mathbf u,\mathbf a'^\top A^{-1}\mathbf u)=\mathbf u^\top A^{-1}\mathbf u$ for every
 scalar $y$.
\end{enumerate}
No symmetry of $A$ or $A'$ is assumed.
\end{lemma}
\begin{proof}
(a) Comparing columns gives $\mathbf u'^{\top}=\mathbf u^\top A^{-1}A'$. (b) We have
$A'\,(A^{-1}\mathbf u,0)^\top=(\mathbf u,\mathbf a'^\top A^{-1}\mathbf u)^\top$.
\end{proof}

\emph{Moments.} In the separation argument of (S), a border node is counted when its moment relative to a reference
node vanishes. That moment is the bilinear form $\mathbf u^\top A^{-1}\mathbf u$ in the inverse of the common fibre $A$
of its residue class. After the elimination, the common fibre of a class with two border nodes is replaced by a mixed
block $A'$ as in Lemma~\ref{lem:bordered}. Its lower node is a row and its upper node a column, or conversely. The
product added to $A'^{-1}$ is exactly the product subtracted in \eqref{eq:schur3}. So, by Lemma~\ref{lem:bordered}, the
moments of the counted classes are unchanged. The single border nodes outside the window are separated from it as in
(S), since two singular contributions in an interval of length $2d<p$ cannot cancel. Applied to the auxiliary geometry
$(\bar r,d)$, the count of (S) then gives \eqref{eq:PiU}. It uses the third-order analogue of \eqref{eq:corner}. Call
the \emph{coupling} of a residue class the residue modulo $p$ of its border entry minus its moment against the reference
node. Suppose that $n_+$ kernel vectors on each side survive, and that at most $n_0$ residue classes have a nonzero
coupling. Then
$\sum_i\min(v_i,3)\ge d+n_++(n_+-n_0)_+$, where $v_1,v_2,\dots$ are the valuations of the
Smith invariants. The first term $d$ is the factor $p$ of $p\,\Xi$.

\emph{Members.} Deleting $\ell$ common nodes lowers the trivial bound $2(r+m)$ by $2\ell$. Consider a class with one
deleted common node whose bordered system keeps its value at infinity $\vartheta$ (Lemma~\ref{lem:interp} below). Then
the class keeps its moment. Indeed, let $A$ be the common fibre and delete its last coordinate. The residual of the
smaller solution is supported on that coordinate. Since the two values at infinity are equal, its pairing with the data
vanishes. A class whose value at infinity changes gives a new separation instead. Classes with
two or more deleted nodes, and classes outside the window, cost at most one each. Let $k_1$, $k_2$ and $k'$ be the
numbers of classes of these three kinds. Then $k_1+2k_2+k'\le\ell$. An elementary inequality for the count
$n_++(n_+-n_0)_+$ shows that the loss is at most $2\ell$. Hence the bound holds for every
member with $\epsilon_\tau=0$.

\emph{Larger $n$.} Let $np\le m<(n+1)p$. Each class of the remaining border absorbs $n$ pairs of border nodes. Its block
after the elimination is a mixed block. Its rows and columns are the nodes of $2n+1$ consecutive quotients
$\beta,\beta+1,\dots,\beta+2n$. Here $\beta$ is the number of common nodes in the class, and the quotient of a node $t$ in
the class of a residue $i$ is $(t-i)/p$. The window counts are those of the first case at the scale $r+np$, translated by
$np$. The new difficulty is to show that a remaining border node separated from the window by a pole has a nonzero
coupling. The values at infinity are now those of $(n+1)\times(n+1)$ bordered systems, which we describe first.

\begin{lemma}[bordered systems and rational interpolation]\label{lem:interp}
Let $\mathbb K$ be a field and $\mathcal Z=\{0,1,\dots,2n\}$, with $(2n)!\ne0$ in $\mathbb K$. Let
$\mathcal L:\mathcal Z\to\mathbb K$, and let $\mathcal Z=X\sqcup Y$ with $|X|=n+1$ and $|Y|=n$. Consider the bordered
system
\[
 \sum_{y\in Y}\frac{\mathcal L(y)-\mathcal L(x)}{y-x}\,z_y+\vartheta=\mathcal L(x)\qquad(x\in X)
\]
in the unknowns $z_y$ and $\vartheta$, with data $\mathcal L$. Its solutions correspond bijectively to the pairs $(\mathcal N,\mathcal D)$ of
polynomials with $\mathcal D$ monic of degree $n$, $\deg\mathcal N\le n$ and $\mathcal N(t)=\mathcal L(t)\mathcal D(t)$
for $t\in\mathcal Z$. Under this correspondence $\vartheta$ is the coefficient of $t^n$ in $\mathcal N$. In particular the
system is nonsingular exactly when the interpolation problem of type $(n,n)$ has a unique solution. Then
$\vartheta=\lim_{t\to\infty}\mathcal N(t)/\mathcal D(t)$ depends only on $\mathcal L$ and not on the split. We call $\vartheta$ the \emph{value at infinity} of the system.
\end{lemma}
\begin{proof}
Let $\mathcal W(t)=\prod_{y\in Y}(t-y)$, and let $\mathcal W_y=\mathcal W/((t-y)\mathcal W'(y))$ be the Lagrange
polynomials of $Y$. Given a solution, put $\mathcal D=\mathcal W-\sum_yz_y\mathcal W'(y)\,\mathcal W_y$ and
$\mathcal N=\vartheta\mathcal W-\sum_yz_y\mathcal L(y)\mathcal W'(y)\,\mathcal W_y$. Then
$\mathcal N(y)=\mathcal L(y)\mathcal D(y)$ for $y\in Y$. Dividing $\mathcal N(x)-\mathcal L(x)\mathcal D(x)$ by
$\mathcal W(x)\ne0$ gives the equation of the row $x$. Conversely, Lagrange interpolation on $Y$ recovers
$z_y=-\mathcal D(y)/\mathcal W'(y)$ and $\vartheta$ from $(\mathcal N,\mathcal D)$.
\end{proof}

\begin{lemma}[harmonic data]\label{lem:harmonic}
Let $\mathcal A\in\mathbb K$ with $\mathcal A+i\ne0$ for $0\le i<2n$, and let
$\mathcal L_{\mathcal A}(t)=\sum_{0\le i<t}1/(\mathcal A+i)$. Let the characteristic of $\mathbb K$ be $0$ or larger than
$2n$. For every split, the bordered system of Lemma~\ref{lem:interp} with $\mathcal L=\mathcal L_{\mathcal A}$ is
nonsingular. Its value at infinity is $\vartheta=\mathcal L_{\mathcal A}(n)+\mathsf c_n$, with $\mathsf c_n\in\mathbb K$
depending only on $n$. Consequently, for the harmonic data $\mathcal L(t)=H_t$ on the quotients $\beta,\dots,\beta+2n$ the
value at infinity is $H_{\beta+n}+\mathsf c_n$. The values at infinity of two classes with $\beta$ and $\beta+1$ common nodes differ by
$1/(\beta+n+1)$.
\end{lemma}
\begin{proof}
Write $\Delta$ for the forward difference. For $l\ge1$ one has
$\Delta^l\mathcal L_{\mathcal A}(k)=(-1)^{l-1}(l-1)!/(\mathcal A+k)_l$, where $(y)_l$ is the rising factorial. One also
has $\Delta^l[\binom tk f(t)](0)=\binom lk\Delta^{l-k}f(k)$. By Lemma~\ref{lem:interp} a solution is a pair
$(\mathcal N,\mathcal D)$ with $\mathcal N=\mathcal L_{\mathcal A}\mathcal D$ on $\mathcal Z$. Since $\mathcal N$ has degree
at most $n$, this means $\Delta^l(\mathcal L_{\mathcal A}\mathcal D)(0)=0$ for $n<l\le2n$. Write
$\mathcal D=\sum_k\mathsf d_k\binom tk$ and insert the two formulas. After removing nonzero factors of the rows and
columns, these $n$ equations become a system for the coefficients $\mathsf d_k$ with $k<n$. Its matrix is the Cauchy
matrix $\bigl(1/(l-k)\bigr)_{n<l\le2n,\,0\le k<n}$, and its right-hand side comes from $\mathsf d_n=1$. The Cauchy matrix
is nonsingular. So $\mathcal D$ is unique, which is the nonsingularity.

The same computation shows that $\mathsf d_k/(\mathcal A+k)_{n-k}$ does not depend on $\mathcal A$. For $k<n$,
$\Delta^{n-k}\mathcal L_{\mathcal A}(k)$ is a constant multiple of $1/(\mathcal A+k)_{n-k}$. Hence
$\vartheta=\Delta^n(\mathcal L_{\mathcal A}\mathcal D)(0)/n!=\mathcal L_{\mathcal A}(n)+\mathsf c_n$ with
$\mathsf c_n$ independent of $\mathcal A$. For the harmonic data on the quotients $\beta+t$ we have
$H_{\beta+t}=H_\beta+\mathcal L_{\beta+1}(t)$. A constant added to $\mathcal L$ is added to $\vartheta$. So the value at infinity is
$H_\beta+\mathcal L_{\beta+1}(n)+\mathsf c_n=H_{\beta+n}+\mathsf c_n$.
\end{proof}

\begin{remark}\label{rem:hahn}
We note without proof that the interpolation problem of Lemma~\ref{lem:harmonic} can be solved in closed form. Its denominator is the polynomial
$\mathcal D_n(t)=\sum_{k=0}^n\binom nk\binom{2n-k}{n-k}(\mathcal A+k)_{n-k}\,t(t-1)\cdots(t-k+1)$, of Hahn type \cite{KLS}.
The constant is $\mathsf c_n=2H_n$. The determinant of the bordered system is
$\operatorname{sf}(n)^2\operatorname{sf}(n-1)^2\big/\bigl(\prod_{x\in X,y\in Y}|x-y|\cdot\prod_{i<n}(\mathcal A+i)_{n+1}\bigr)$,
with $\operatorname{sf}(n)=\prod_{i\le n}i!$. All its prime factors are at most $2n$ or divide some $\mathcal A+i$. These
closed forms are not used in the proof and are not part of the formalization.
\end{remark}

\emph{Separation for $n\ge1$.} Between a remaining border node and the reference node of the window lies a cut. Suppose
first that the cut is a pole of the second digamma term. Deform the shift to $s+pY$, with an indeterminate $Y$. Evaluate
at the value of $Y$ at which that pole becomes central. After multiplication by the vanishing factor, the bordered systems
of the two classes become block systems built from Cauchy matrices on the quotients. The square blocks are nonsingular
by Cauchy's determinant formula. A rectangular Cauchy matrix with $k+1$ rows and $k$ columns has a left kernel vector
with nonzero coordinate sum. Otherwise the rational function $\sum_iy_i/(T_i-T)$, formed from a kernel vector $(y_i)$
and the row nodes $T_i$, would vanish at the $k$ column nodes while its numerator has degree less than $k$. This gives the
two values at infinity the values $0$ and the weight of the term, which differ.

Suppose next that the cut is a multiple of $p$. Take instead the leading coefficient in $Y$. Every fractional pole term
tends to $0$, and the data become the harmonic numbers. The two classes have $\beta$ and $\beta+1$ common nodes.
Lemma~\ref{lem:harmonic} gives values at infinity that differ by $1/(\beta+n+1)$. This is a unit because $\beta+2n+1<p$, by
$q(r+2m)<p^2$.

In both cases the comparison is made on the cleared numerators of the original determinants, not on inverses. The result
returns to $Y=0$ by Gauss's lemma. The cleared denominators are products of the elementary denominators of the entries,
which are units at $Y=0$. So a lower bound for the Gauss valuation of the deformed determinant gives the same lower bound
for every coefficient of the original one. With these separations, the count and the treatment of the members are those
of the case $n=0$.

\subsection{Profiles}\label{ssec:profiles}
Let $\mathcal E$ be the set of integer expressions in $r,p,m,k$ built from integer constants with the following
operations: addition, multiplication by integer constants, truncated subtraction $(X-Y)_+$, $\min$, $\max$, and the
distance $\|X\|_p$ of an expression $X$ to $p\Z$. The quantities $p\lfloor r/p\rfloor$ and $r\bmod p$ are also allowed.
Every bound of Theorem~\ref{thm:local} lies in $\mathcal E$. For (R) and (S) this is \eqref{eq:PiR} and \eqref{eq:PiS}.
The ends of the intervals $\mathcal K_{k,\iota}$ are $kp=p\lfloor r/p\rfloor+(k-k_0)p$ plus $0$, $k_p$, $k_p-1$, $p-1$ or
$r_p-1$. For (U$_n$) this is \eqref{eq:PiU} with $d=\|m\|_p$ and $Z_i=p\lfloor r/p\rfloor+(n+i)p$. For (W) the window
ends are the cuts nearest to the ends of the window range, which are expressions of the same kind.

The \emph{homogenisation} $\hat\tau(\theta)$ of $\tau\in\mathcal E$ replaces $r,p,m,k$ by $1,\theta,\alpha,\kappa\theta$.
It replaces every constant by $0$, $\|X\|_p$ by $\theta\|\hat X/\theta\|$ and $p\lfloor r/p\rfloor$ by
$\theta\lfloor1/\theta\rfloor$. It is piecewise affine with rational breakpoints, finitely many in every interval bounded
away from $0$.

\begin{lemma}[homogenisation]\label{lem:homog}
For each $\tau\in\mathcal E$ there is an explicit constant $\err(\tau)$ such that
$|\tau(r,p,m(r),k_p)-r\hat\tau(p/r)|\le\err(\tau)$ for all $r,p\ge1$. It is computed recursively from
$\err(m)=\frac12$, $\err(k)=1-\frac1q$ and the rules $\err(\tau+\tau')=\err(\tau)+\err(\tau')$ and
$\err(\min(\tau,\tau'))\le\max(\err(\tau),\err(\tau'))$.
\end{lemma}
\begin{proof}
The real version of $\tau$, with the constants dropped, is positively homogeneous of degree one. Every operation is
$1$-Lipschitz in the maximum norm. The inputs differ from their homogeneous values by $|m-\alpha r|\le\frac12$ and
$|k_p-\kappa p|\le1-\frac1q$. Induct on $\tau$.
\end{proof}

The conditions of Theorem~\ref{thm:local} are linear inequalities in $r,p,m,k$, strict or not. They homogenise to
intervals of $\theta$. At an endpoint the error of such an inequality is a combination of $m-\alpha r$ and $k_p-\kappa p$.
With the rounding $m=\lceil\alpha r-\frac12\rceil$ two endpoints are \emph{exact}. These are $\theta=2\alpha$, from
$2m<p$, and $\theta=(1+2\alpha)/(1+\kappa)$, from $\bar N\le p+k$. At an exact endpoint the condition holds for every prime on
the correct side. The other endpoints are exact, or hold outside a strip of primes of bounded width around the endpoint.
On that strip another bound must be used.

\subsection{Tables for $\Lambda$ and the arithmetic estimate}\label{ssec:checker}
For the root class $j$ let $\mathcal P_j(\theta)$ be the minimum of the homogenised bounds $\hat\tau(\theta)$ over the
bounds $\tau$ of Theorem~\ref{thm:local} valid at $\theta$. A \emph{table for $\Lambda$}, for the function $F_s$ and
the geometry of parameter $c$, consists of the parameters $(q,\alpha,K)$ and the following data for every root class $j$:
\begin{itemize}
\item breakpoints $0<\theta_0<\theta_1<\dots<\theta_{\rm end}$;
\item on each interval $(\theta_{i-1},\theta_i]$, an affine function $\ell_i$, a constant $\err_i\ge0$ and a local bound
 $\tau_i$ of Theorem~\ref{thm:local}, with a second bound at each inexact endpoint of the range of $\tau_i$;
\item a constant $\ell_0$ for the initial interval $(0,\theta_0]$;
\end{itemize}
together with a rational number $\Lambda_{\rm fin}$. We call the table \emph{admissible} if it satisfies the following
conditions.
\begin{enumerate}[label=(C\arabic*),leftmargin=*]
\item For each shift of the table, every unit class modulo the modulus $q^*$ of the prime number theorem used is assigned
 a root class with a table.
\item On each piece $\ell_i$ is nonnegative, and the piece lies in the validity interval of its bound. Moreover
 $\hat\tau$ lies below $\ell_i$ at both ends and at both one-sided limits of each breakpoint inside the piece.
\item At each inexact endpoint the second bound is valid on both sides and lies below $\ell_i$ there. The constant
 $\err_i$ exceeds the error constants of Lemma~\ref{lem:homog}, the numbers $\epsilon_\tau K$, the strip widths times the slopes,
 and the trivial bound at the scales below an explicit threshold.
\item $\ell_0\ge2+2\alpha$ and $0<\theta_0\le1$. Also $\theta_{\rm end}\ge\max(q(1+2\alpha),p_0/2)$ and
 $\kappa_j\theta_{\rm end}\ge1+2\alpha$.
\item $\frac1{\varphi(q)}\sum_j\bigl(\ell_0\theta_0+\sum_i\int_{\theta_{i-1}}^{\theta_i}\ell_i(\theta)\,d\theta\bigr)
 =\Lambda_{\rm fin}$.
\end{enumerate}
Each condition is a finite set of comparisons of rational numbers.

\begin{theorem}[denominators]\label{thm:content}
Let $0<c\le1$, and let $F_s$ be the function with shift $s$ of denominator $q$. Suppose that an admissible table for $\Lambda$ of
depth $K$ has $s$ among its shifts. Then for every even $r$ and every member with $|I|\le K$,
\[
 \log\den P_{r,I}\le(\Lambda_{\rm fin}+o(1))\,r^2 .
\]
\end{theorem}
\begin{proof}
By definition $\den P_{r,I}=\prod_pp^{e_p(r,I)}$. Fix a piece $(\theta_{i-1},\theta_i]$ of the root class $j$. Let $p$ be
a prime in a unit class assigned to $j$, with $\theta_{i-1}r<p\le\theta_ir$. Let $r$ be beyond the threshold of (C3).
Then the conditions of the bound $\tau_i$ hold, either exactly or outside a strip where the second bound applies.
Lemma~\ref{lem:homog} with (C2) and (C3) gives $e_p(r,I)\le r\,\ell_i(p/r)+\err_i$. Below the threshold the trivial bound
is at most $\err_i$.

For $p\le\theta_0r$ the bound (T1) gives $e_p\le\ell_0r+O(1)$. The primes with $p^2\le q(r+2m)$, the prime powers and the
finitely many primes below $p_0$ contribute $O(r^{3/2}\log r)$. For $p>\theta_{\rm end}r$, (C4) gives $\bar N\le p$ and
$k_p\ge \bar N$, and (T5) gives $e_p=0$.

Fix a unit class modulo $q^*$. By the prime number theorem in arithmetic progressions,
\[
 \sum_{\substack{\theta_{i-1}r<p\le\theta_ir\\ p\text{ in the class}}}r\,\ell_i(p/r)\log p
 =\frac{r^2}{\varphi(q^*)}\int_{\theta_{i-1}}^{\theta_i}\ell_i(\theta)\,d\theta+o(r^2).
\]
The class $j$ receives $\varphi(q^*)/\varphi(q)$ unit classes modulo $q^*$, because the map that sends a unit to $qs$
times its inverse permutes the units modulo $q$. Summing and using (C5) gives the theorem. The constants $\err_i$
contribute $O(r)$.
\end{proof}

The tables of this paper are generated from the functions $\mathcal P_j$ themselves. The pieces are the maximal intervals
on which $\mathcal P_j$ is affine and attained by one bound. So $\Lambda_{\rm fin}=\frac1{\varphi(q)}\sum_j\int\mathcal P_j$
up to the initial interval. Every bound of Theorem~\ref{thm:local} is available to the tables: (T1)--(T5), (LP1), (LP2), (R), (S),
the windows (W) and the bounds (U$_n$). The windows with $n=0$ are used for $m<p$, between consecutive cuts.

On $(0,\alpha]$, that is for $p\le m$, the only other bounds are (T1), (T1$'$) and (U$_n$). With (T1$'$) alone the
profile there is the sawtooth $2+2\alpha-\theta\|\alpha/\theta\|$. Its integral is
\[
 \int_0^\alpha\Bigl(2+2\alpha-\theta\Bigl\|\frac\alpha\theta\Bigr\|\Bigr)d\theta=2\alpha+\frac{\alpha^2}2+2\alpha^2\log2,
\]
since $\int_1^\infty\|y\|y^{-3}\,dy=\frac32-2\log2$. The windows with $n=\lfloor m/p\rfloor\ge1$ lower this profile on the
strips $\alpha/(n+1)<\theta\le\alpha/n$. The tables use them for $n\le8$, again with each window running between two
consecutive cuts. For $q=3,4,5,7$ the tables also use (U$_0$) on $(\alpha,2\alpha]$ and (U$_n$) on the strips with
$1\le n\le24$. The table for $q=2$ uses only the other bounds. Beyond the strips used, a table resolves the sawtooth on each
interval $(\alpha/(n+1),\alpha/n]$ with $n<n_0$. It ends with a constant piece on $(0,\alpha/n_0]$, which costs $\alpha^2T_{n_0}$ with
$0<T_n\le1/(4n^2)$. The number $\Lambda_{\rm fin}$ is then rational, and (C5) is an identity between rational numbers.
The bounds depend on the shift $s$ only through the root class, and their profiles only through $(q,j,\alpha)$. So one
table serves all shifts with denominator $q$.

On small $\theta$ the pieces of a table repeat with a period in $1/\theta$. Write $\alpha=a/b$ in lowest terms, and
call $T_\nu(r,p,m,k)=(r+\nu bp,\,p,\,m+\nu ap,\,k)$ the \emph{translation of order} $\nu\ge0$. It keeps $p$, $k$ and
$m-\alpha r$, and hence the rounding of $m$. A list of local bounds is \emph{closed under translation}, with constant
$K_0$, if for every bound $\tau$ of the list and every $\nu$ there is a bound $\tau'$ of the list with
$\tau'(T_\nu y)=\tau(y)+\nu K_0bp$ for all $y=(r,p,m,k)$, whose conditions hold at $T_\nu y$ whenever those of $\tau$
hold at $y$.

\begin{theorem}[periodic tables]\label{thm:periodic}
Let a list of local bounds be closed under translation with constant $K_0$. Suppose that for every $r\ge R$ and every
prime $p$ with $sr<p\le tr$ some bound $\tau$ of the list has its conditions satisfied at $y=(r,p,m(r),k_p)$ and
$\tau(y)\le Ar+Bp+C$. Let $1/s_\nu=1/s+\nu b$ and $1/t_\nu=1/t+\nu b$. Then for every member at a scale
$r\ge Rt/t_\nu$ and every prime $p$ with $s_\nu r<p\le t_\nu r$,
\[
 e_p(r,I)\le Ar+\bigl(B+(K_0-A)\nu b\bigr)p+C .
\]
\end{theorem}
\begin{proof}
Put $r'=r-\nu bp$ and $m'=m(r)-\nu ap$. Then $r'/p=r/p-\nu b$, so $s<p/r'\le t$, and $r'\ge r\,t_\nu/t\ge R$. Also
$m'=m(r')$, since $m'-\alpha r'=m(r)-\alpha r$. So the hypothesis gives a bound $\tau$ whose conditions hold at
$y=(r',p,m',k_p)$, with $\tau(y)\le Ar'+Bp+C$. The conditions of its translate $\tau'$ hold at $T_\nu y=(r,p,m(r),k_p)$.
Since $\tau'$ is a local bound, $e_p(r,I)\le\tau(y)+\nu K_0bp\le A(r-\nu bp)+Bp+C+\nu K_0bp$.
\end{proof}

So a table for $\Lambda$ on a range where its bounds are closed under translation needs to be checked on one period of
$1/\theta$ only. The pieces of the other periods are its translates, and their integrals are again rational numbers.
The finitely many scales below the thresholds take the trivial bound, as in (C3).

%% =====================================================================
\section{Nonvanishing}\label{sec:nonvanishing}

The deduction needs, for every rational $x_0$ and every large scale, one member of the family that does not vanish at
$x_0$. We prove this in characteristic zero for every frame of \S\ref{ssec:frames}. The number of deleted common nodes
is at most half the number of poles of $F$ in a period, counted with order. The proof has three steps. The kernel of
the matrix is a space of solutions of a Hermite interpolation problem. A discrete Wronskian bounds its dimension. The
values of the solutions at the common nodes then tell us which nodes to delete.

Throughout this section $F$ is a frame with $d$ poles per period, counted with order, and $\mathcal G(r,m)$ is a
geometry with $r\ge1$. Shifts of order $0$ play no role here, since removing them from $S$ does not change $F$. By the
balance $d\ge2$. Fix $x_0\in\Q$ and put $\vartheta_t=x_0-e_t$. Then $\Phi_{\mathcal G}(x_0)$ has the entries
$(f_{t'}-f_t)/(t'-t)$ for $t\ne t'$ and $\vartheta_t$ on the diagonal of the common block. Put $n=r+m$, and let
$Z=[0,\bar N)$ be the set of all nodes and $C=[0,r)$ the set of common nodes. So $|R|=|Q|=n$ and $|Z|+|C|=2n$.

\subsection{The kernel as an interpolation space}\label{ssec:selection}
Let $\mathcal V$ be the space of pairs $(\mathsf u,\mathsf w)$ of polynomials over $\Q$ of degree less than $n$ that
satisfy the \emph{value conditions} and the \emph{jet conditions}
\[
 \mathsf w(t)=f_t\,\mathsf u(t)\quad(t\in Z),\qquad
 \mathsf w'(t)=f_t\,\mathsf u'(t)+\vartheta_t\,\mathsf u(t)\quad(t\in C).
\]
For a finite set $T$ of integers write $\omega_T(z)=\prod_{t\in T}(z-t)$.

\begin{lemma}[kernel]\label{lem:kernelpairs}
The map that sends $(\mathsf u,\mathsf w)\in\mathcal V$ to the vector $\mathbf y$ with
$y_t=\mathsf u(t)/\omega_Q'(t)$ for $t\in Q$ is an isomorphism of $\mathcal V$ onto the right kernel of
$\Phi_{\mathcal G}(x_0)$. The same holds for the left kernel, with $R$ in place of $Q$. In any pair of $\mathcal V$,
$\mathsf u=0$ implies $\mathsf w=0$.
\end{lemma}
\begin{proof}
Let $\mathbf y$ be a vector indexed by $Q$, and put
\[
 \mathsf u(z)=\sum_{t'\in Q}y_{t'}\frac{\omega_Q(z)}{z-t'},\qquad
 \mathsf w(z)=\sum_{t'\in Q}y_{t'}f_{t'}\frac{\omega_Q(z)}{z-t'} .
\]
By Lagrange interpolation at the $n$ nodes of $Q$, this is a bijection between vectors $\mathbf y$ and pairs of
degree less than $n$ that satisfy the value conditions on $Q$, with inverse $y_t=\mathsf u(t)/\omega_Q'(t)$. Let
$t\in X'$ be a border row. Since $\mathsf u(t)/\omega_Q(t)=\sum_{t'}y_{t'}/(t-t')$, and likewise for $\mathsf w$,
\[
 \bigl(\Phi_{\mathcal G}(x_0)\mathbf y\bigr)_t=\sum_{t'\in Q}y_{t'}\frac{f_{t'}-f_t}{t'-t}
 =-\frac{\mathsf w(t)-f_t\mathsf u(t)}{\omega_Q(t)} .
\]
Let $t\in C$ be a common row, and put $\mathsf v=\mathsf w-f_t\mathsf u$. The term $t'=t$ of $\mathsf v$ vanishes. Each
other term vanishes at $t$ with derivative $y_{t'}(f_{t'}-f_t)\,\omega_Q'(t)/(t-t')$. Since
$\mathsf u(t)=y_t\,\omega_Q'(t)$,
\[
 \mathsf w'(t)-f_t\mathsf u'(t)-\vartheta_t\mathsf u(t)=\mathsf v'(t)-\vartheta_t\mathsf u(t)
 =-\omega_Q'(t)\bigl(\Phi_{\mathcal G}(x_0)\mathbf y\bigr)_t .
\]
So $\mathbf y$ lies in the right kernel exactly when the value conditions on $X'$ and the jet conditions on $C$ hold.
As $Z=Q\sqcup X'$, these are the conditions defining $\mathcal V$. The matrix $M$ is symmetric. So the transpose of
$\Phi_{\mathcal G}(x_0)$ is the matrix of the same nodes with the roles of $R$ and $Q$ exchanged. This exchange leaves
$Z$ and $C$ unchanged, and with them $\mathcal V$. This gives the left kernel. Finally, if $\mathsf u=0$ then
$\mathsf w$ vanishes on $Z$. As $|Z|\ge n>\deg\mathsf w$, this gives $\mathsf w=0$.
\end{proof}

\subsection{The corank}
Put
\[
 \mathsf D(z)=\prod_{s\in S}(z+s)^{w(s)},\qquad \mathsf A=\mathsf D\,\delta ,
\]
where $w(s)$ is the order of $s$. Then $\mathsf D$ has degree $d$, and $\mathsf A$ is a polynomial of degree at most
$d-2$, since $\delta(z)=O(|z|^{-2})$ by the balance $\sum_sa_s=0$. Differentiating $\mathsf D\delta=\mathsf A$ gives
$\mathsf D'\delta+\mathsf D\delta'=\mathsf A'$. For integers $k\ge0$ we have $\mathsf D(k)\ne0$, $f_{k+1}-f_k=\delta(k)$
and $\vartheta_{k+1}-\vartheta_k=\delta'(k)$.

\begin{lemma}[cross product]\label{lem:cross}
For two pairs $(\mathsf u_1,\mathsf w_1)$ and $(\mathsf u_2,\mathsf w_2)$ of $\mathcal V$ one has
$\mathsf w_1\mathsf u_2=\mathsf w_2\mathsf u_1$.
\end{lemma}
\begin{proof}
The polynomial $\Pi=\mathsf w_1\mathsf u_2-\mathsf w_2\mathsf u_1$ has degree at most $2n-2$. The value conditions give
$\Pi(t)=0$ for $t\in Z$. At $t\in C$ the value and jet conditions give
\[
 \Pi'(t)=(f_t\mathsf u_1'+\vartheta_t\mathsf u_1)\mathsf u_2+f_t\mathsf u_1\mathsf u_2'
 -(f_t\mathsf u_2'+\vartheta_t\mathsf u_2)\mathsf u_1-f_t\mathsf u_2\mathsf u_1'=0,
\]
with all polynomials evaluated at $t$. So $\Pi$ has at least $|Z|+|C|=2n$ zeros, counted with multiplicity, and
$\Pi=0$.
\end{proof}

The next lemma is the heart of the argument. It is used twice, with all the conditions and with the value conditions
only.

\begin{lemma}[no short pairs]\label{lem:noshort}
Let $N_0\ge1$ and $0\le c_0\le N_0$ be integers, and let $\mathsf u$, $\mathsf w$ be polynomials of degree at most
$\ell$ that satisfy the value conditions on $[0,N_0)$ and the jet conditions on $[0,c_0)$. If
\[
 2\ell+d-2<(N_0-1)+\max(c_0-1,0),
\]
then $\mathsf u=0$.
\end{lemma}
\begin{proof}
Suppose that $\mathsf u\ne0$. Put $\mathsf W(z)=\mathsf w(z+1)\mathsf u(z)-\mathsf w(z)\mathsf u(z+1)$ and
\[
 \mathsf E(z)=\mathsf D(z)\,\mathsf W(z)-\mathsf A(z)\,\mathsf u(z)\,\mathsf u(z+1).
\]
By Taylor's formula $\mathsf W=\sum_{j\ge1}(\mathsf w^{(j)}\mathsf u-\mathsf w\mathsf u^{(j)})/j!$. The term $j=1$ has
degree at most $2\ell-2$, because the coefficients of $z^{2\ell-1}$ in $\mathsf w'\mathsf u$ and in $\mathsf w\mathsf u'$
are both $\ell$ times the product of the leading coefficients. The terms with $j\ge2$ have degree at most $2\ell-2$ too.
Hence $\deg\mathsf E\le2\ell+d-2$.

Now we count zeros. Let $k\ge0$ with $k+1<N_0$, and put $U(z)=\mathsf u(z)\mathsf u(z+1)$. The value conditions at $k$
and $k+1$ give $\mathsf W(k)=\delta(k)U(k)$, so $\mathsf E(k)=(\mathsf D\delta-\mathsf A)(k)\,U(k)=0$. Let moreover
$k+1<c_0$. The jet conditions at $k$ and $k+1$ give $\mathsf W'(k)=\delta(k)U'(k)+\delta'(k)U(k)$. Hence
\[
 \mathsf E'(k)=(\mathsf D'\delta+\mathsf D\delta'-\mathsf A')(k)\,U(k)+(\mathsf D\delta-\mathsf A)(k)\,U'(k)=0 .
\]
So $\mathsf E$ has at least $(N_0-1)+\max(c_0-1,0)$ zeros, counted with multiplicity. This exceeds its degree, and
$\mathsf E=0$.

Let $\mathsf g$ be the greatest common divisor of $\mathsf u$ and $\mathsf w$, and write $\mathsf u=\mathsf g\mathsf u_1$
and $\mathsf w=\mathsf g\mathsf w_1$. Then $\mathsf E=\mathsf g(z)\mathsf g(z+1)\mathsf E_1$, where $\mathsf E_1$ is the
same expression for $(\mathsf u_1,\mathsf w_1)$. So $\mathsf E_1=0$, with $\mathsf u_1\ne0$ and $\mathsf u_1$,
$\mathsf w_1$ coprime. Take a shift $s_0$ of positive order and put $a=-s_0$. Then $\mathsf D(a)=0$, and $\mathsf A(a)$
is $a_{s_0}$ or $-b_{s_0}$, according to the order, times $\prod_{s\ne s_0}(s-s_0)^{w(s)}$. So $\mathsf A(a)\ne0$.
The points $a+j$ with $j\in\Z\setminus\{0\}$ are not roots of $\mathsf D$, since distinct shifts in $(0,1]$ differ by
less than $1$. Evaluating $\mathsf E_1=0$ at $a$ gives $\mathsf A(a)\mathsf u_1(a)\mathsf u_1(a+1)=0$.

Suppose first that $\mathsf u_1(a+j)=0$ for some $j\ge1$. Evaluating $\mathsf E_1=0$ at $a+j$ gives
$\mathsf D(a+j)\mathsf w_1(a+j)\mathsf u_1(a+j+1)=0$. Here $\mathsf w_1(a+j)\ne0$ by coprimality, so
$\mathsf u_1(a+j+1)=0$. If $\mathsf u_1(a+1)=0$, induction gives infinitely many roots of $\mathsf u_1$. Suppose next
that $\mathsf u_1(a-j)=0$ for some $j\ge0$. Evaluating $\mathsf E_1=0$ at $a-j-1$ gives
$\mathsf D(a-j-1)\mathsf w_1(a-j)\mathsf u_1(a-j-1)=0$, so $\mathsf u_1(a-j-1)=0$. If $\mathsf u_1(a)=0$, induction
again gives infinitely many roots. Both cases contradict $\mathsf u_1\ne0$.
\end{proof}

\begin{proposition}[corank]\label{prop:corank}
$\dim\mathcal V\le\lfloor d/2\rfloor$. Hence $\corank\Phi_{\mathcal G}(x_0)\le\lfloor d/2\rfloor$.
\end{proposition}
\begin{proof}
Put $\varsigma=\lfloor d/2\rfloor$. If $\varsigma\ge n$ there is nothing to prove, since the right kernel has dimension
at most $n$. Let $\varsigma<n$. By Lemma~\ref{lem:kernelpairs} every nonzero pair has $\mathsf u\ne0$. By
Lemma~\ref{lem:cross} the difference $\deg\mathsf w-\deg\mathsf u$, with $\deg0=-\infty$, is the same for all nonzero
pairs. So one of the two coordinates, call it $\ell$, has the larger degree in every pair. Map each pair to the
coefficients of $z^{n-\varsigma},\dots,z^{n-1}$ in $\ell$. This is a linear map to $\Q^\varsigma$. A pair in its
kernel has both coordinates of degree at most $n-\varsigma-1$. Apply Lemma~\ref{lem:noshort} with $N_0=\bar N$,
$c_0=r$ and $\ell=n-\varsigma-1$. Its condition reads $2n-2\varsigma+d-4<2n-2$, that is $d-2<2\varsigma$, which holds
since $2\lfloor d/2\rfloor\ge d-1$. So $\mathsf u=0$, and the pair is zero. The map is injective, and
$\dim\mathcal V\le\varsigma$. The second statement follows from Lemma~\ref{lem:kernelpairs}.
\end{proof}

\subsection{The choice of the deleted nodes}\label{ssec:nvgeneral}
It remains to see the kernel on the common nodes. At a common node $t$ the coordinates of a right and of a left kernel
vector are $\mathsf u(t)/\omega_Q'(t)$ and $\mathsf u(t)/\omega_R'(t)$, by Lemma~\ref{lem:kernelpairs}.

\begin{lemma}[visibility]\label{lem:visible}
Let $r\ge d-2$. If a pair of $\mathcal V$ has $\mathsf u(t)=0$ for every $t\in C$, then it is zero.
\end{lemma}
\begin{proof}
The value conditions give $\mathsf w(t)=0$ for $t\in C$. So $\omega_C$ divides $\mathsf u$ and $\mathsf w$. Write
$\mathsf u=\omega_C\mathsf u_0$ and $\mathsf w=\omega_C\mathsf w_0$, with degrees less than $n-r=m$. If $m=0$ then
$\mathsf u_0=0$. Let $m\ge1$. For $t\in Z\setminus C$ we have $\omega_C(t)\ne0$, so $\mathsf w_0(t)=f_t\mathsf u_0(t)$.
For $t\in C$ we have $\mathsf u(t)=0$, $\mathsf u'(t)=\omega_C'(t)\mathsf u_0(t)$ and $\mathsf w'(t)=\omega_C'(t)\mathsf
w_0(t)$, with $\omega_C'(t)\ne0$. So the jet condition at $t$ reads $\mathsf w_0(t)=f_t\mathsf u_0(t)$. The pair
$(\mathsf u_0,\mathsf w_0)$ therefore satisfies the value conditions on all of $Z$. Apply Lemma~\ref{lem:noshort} with
$N_0=\bar N=r+2m$, $c_0=0$ and $\ell=m-1$. Its condition reads $2m+d-4<r+2m-1$, that is $r\ge d-2$. So
$\mathsf u_0=0$. Then $\mathsf u=0$, and the pair is zero by Lemma~\ref{lem:kernelpairs}.
\end{proof}

\begin{theorem}[nonvanishing for frames]\label{thm:nvframe}
Let $F$ be a frame, and let $r\ge\max(1,d-2)$ and $m\ge0$. For every $x_0\in\Q$ some member $P_{r,I}$ with
$|I|\le\lfloor d/2\rfloor$ satisfies $P_{r,I}(x_0)\ne0$.
\end{theorem}

\begin{proof}
Let $k$ be the corank of $\Phi_{\mathcal G}(x_0)$. Then $k=\dim\mathcal V\le\lfloor d/2\rfloor$ by
Proposition~\ref{prop:corank}. The map $(\mathsf u,\mathsf w)\mapsto(\mathsf u(t))_{t\in C}$ is injective on
$\mathcal V$, by Lemma~\ref{lem:visible}. Its image is a subspace of $\Q^C$ of dimension $k$. So there is a set
$I\subseteq C$ with $|I|=k$ such that the coordinates in $I$ already determine the pair. By the remark before
Lemma~\ref{lem:visible}, no nonzero right kernel vector and no nonzero left kernel vector vanishes on $I$.

Write $\Phi=\Phi_{\mathcal G}(x_0)$, of size $n$ and rank $n-k$. A linear relation among the columns $Q\setminus I$
is a right kernel vector that vanishes on $I$. So these $n-k$ columns are independent, and they span the column space.
In the same way the rows $R\setminus I$ span the row space. Let $\mathbf z$ be a vector on $Q\setminus I$ with
$\Phi[R\setminus I,Q\setminus I]\,\mathbf z=0$. Every row of $\Phi$ is a combination of the rows $R\setminus I$, so
$\Phi[R,Q\setminus I]\,\mathbf z=0$, and $\mathbf z=0$ by the independence of the columns. Hence
$P_{r,I}(x_0)=\det\Phi[R\setminus I,Q\setminus I]\ne0$.
\end{proof}

For the frames of the first plane $d\le3$, and one deleted node suffices. For the frames $F^{(2)}_{v,u}$ with $uv\ne0$
we have $\mathsf D=(z+1)^2(z+\frac12)$ and $\mathsf A=-(v+\frac u2)z-\frac{v+u}2$, with $\mathsf A(-1)=\frac v2\ne0$.
For $v=0$ we have $\mathsf D=(z+1)(z+\frac12)$ and $\mathsf A=-\frac u2$. For $u=0$ we have $\mathsf D=(z+1)^2$ and
$\mathsf A=-v$. So again one deleted node suffices. The formal proof uses the same interpolation argument for every
frame.

\subsection{The functions of \eqref{eq:frames}}
\begin{theorem}[selection]\label{thm:selection}
Let $F$ be a function of \eqref{eq:frames}, and consider the alternating geometry of parameter $c$. There is $K$, depending
only on $|\mathcal T|$, with the following property. For every large even $r$ and
every $x_0\in\Q$ some member $P_{r,I}$ with $|I|\le K$ satisfies $P_{r,I}(x_0)\ne0$.
\end{theorem}
\begin{proof}
The function $F_s$ is $\omega$ times the frame on $\{s,1\}$ with $a_1=1$ and $a_s=-1$, so $d=2$. The function $F_\chi$
is $-\frac16F^+_{2,-3,1}$, with $d=3$. Multiplying a frame by $\omega\ne0$ multiplies $M$ by $\omega$. So the matrix
$\Phi_{\mathcal G,I}(x_0)$ of $\omega F$ is $\omega$ times that of $F$ at $x_0/\omega$. Theorem~\ref{thm:nvframe} gives the
claim with $K=\lfloor|\mathcal T|/2\rfloor=1$ for $r\ge2$.
\end{proof}

The formal proof of Theorem~\ref{thm:selection} uses a coarser count of the same kind, with $K\le64$. The deduction
needs only that the depth is bounded.

%% =====================================================================
\section{Analytic decay}\label{sec:decay}

In this section we bound $|P_{r,I}(x)|$ for $x$ near $\xi$. The argument is that of OpenAI \cite{OpenAI2026} for square
determinants, with three changes. The determinant has two families of nodes, one for the rows and one for the columns.
The comparison measure is a density on finitely many segments rather than an explicit equilibrium measure. The
inequalities that remain are verified on a finite cover by cells. The result is one theorem for every geometry parameter
$c\le1$ and every function of \eqref{eq:frames}.

\subsection{A contour representation}
Hermite's formula gives $\Lw_F[t,t']=\frac1{2\pi i}\oint F(z)\upsilon_t(z)\upsilon_{t'}(z)\,dz$ with
$\upsilon_t(z)=1/(z-t)$. The contour encloses the nodes and no pole of $F$. Write
$F=\sum_{\varsigma\in\mathcal T}w_\varsigma\psi(z+\varsigma)$ as in \S\ref{sec:construction}. The reflection formula
$\psi(w)=\psi(1-w)-\pi\cot\pi w$ splits $F=\tilde F+E$. Here $\tilde F$ is analytic in the left half-plane and
$O(1/|z|)$ there, and
\[
 E(z)=-\pi\sum_{\varsigma\in\mathcal T}w_\varsigma\cot\pi(z+\varsigma).
\]
The function $E$ is $1$-periodic. Since the weights sum to zero, $E(z)=O(e^{-2\pi|\Im z|})$ for $|\Im z|\ge1$. For
$F_s$ one has $E(z)=\pi[\cot\pi(z+s)-\cot\pi z]$.

Let $\Gamma$ be a polygonal path, symmetric under complex conjugation. It crosses the real axis at a point $x^*<0$ that
is not a pole of $F$, runs through the left half-plane and ends in two rays. The poles of $F$ are the points
$-(i+\varsigma)$ with $i\ge0$ and $\varsigma\in\mathcal T$. Moving the contour to $\Gamma$ picks up the residues of $F$
at the poles in $(x^*,0)$. The part of the integrand with $\tilde F$ integrates to $0$ over $\Gamma$. Hence, for all
integers $t,t'\ge0$,
\begin{equation}\label{eq:gram}
 \Lw_F[t,t']=\int\upsilon_t\upsilon_{t'}\,d\Omega .
\end{equation}
Here $\Omega$ is the complex measure with the mass $-\operatorname{res}_y(F)$ at each pole $y\in(x^*,0)$ of $F$ and the
density $(2\pi i)^{-1}E(z)\,dz$ on $\Gamma$. This is a Gram representation with respect to a measure with finitely many
atoms and an exponentially small density on $\Gamma$. After scaling by $r$, the atoms lie on a segment $[-r_0,0]$ with
density $|\mathcal T|$. The separation condition on $x^*$ is a finite check on the fractional parts of $\varsigma-x^*$
for $\varsigma\in\mathcal T$.

\subsection{Two-family Andr\'eief bound}
For a set $C''\subseteq C$ of common nodes put $\mathcal M_{C''}=\omega\,\Lw_F[C''\sqcup X',C''\sqcup Y']$. Consider a
member with common nodes $C'=C\setminus I$. Put $R'=C'\sqcup X'$, $Q'=C'\sqcup Y'$ and $n_*=|C'|+m$. By multilinearity in
the diagonal,
\begin{equation}\label{eq:expand}
 P_{r,I}(\xi+h)=\sum_{C''\subseteq C'}h^{|C'|-|C''|}\det\mathcal M_{C''} .
\end{equation}
Andr\'eief's identity \cite{Forrester2018Andreief} for the two families $(\upsilon_t)_{t\in R'}$ and
$(\upsilon_{t'})_{t'\in Q'}$ gives
\[
 \det\Lw_F[R',Q']=\frac1{n_*!}\int\det[\upsilon_t(z_i)]\,\det[\upsilon_{t'}(z_i)]\prod_{i=1}^{n_*}d\Omega(z_i).
\]
Cauchy's determinant formula evaluates both factors. We take absolute values. The logarithmic potential theory behind
what follows is in \cite{SaffTotik}. We obtain
\begin{multline*}
 |\det\Lw_F[R',Q']|\le\frac1{n_*!}\int\prod_{i<i'}|z_i-z_{i'}|^2\,\\
 \prod_{\mathcal F\in\{R',Q'\}}\frac{\prod_{t<t'\in\mathcal F}|t'-t|}{\prod_{i,\,t\in\mathcal F}|z_i-t|}\
 \prod_id|\Omega|(z_i).
\end{multline*}
The integrand is the exponential of minus a discrete logarithmic energy. The free points $z_i$ repel each other with
weight $1$. Each node family repels itself with weight $\frac12$ and attracts the free points with weight $1$. The
density $|E|\approx e^{-2\pi|\Im z|}$ acts as an external field. The rows and columns differ only in the interleaved
border. So the two node families have the same macroscopic density. After scaling by $r$, both are
\[
 \nu_c=\mathrm{Leb}|_{[0,1]}+\tfrac12\mathrm{Leb}|_{[1,1+c]},
\]
of mass $|\nu_c|=1+c/2$.

\subsection{The energy certificate}
For a regularisation radius $\epsilon>0$ let $k_\epsilon(z,w)=-\frac12\log(|z-w|^2+\epsilon^2)$. Then $k_\epsilon(z,w)\le-\log|z-w|$.

\begin{lemma}[conditional positive definiteness]\label{lem:cpd}
Let $z_i\in\C$ be finitely many points and $y_i$ real numbers with $\sum_iy_i=0$. Then
$\sum_{i,j}y_iy_jk_\epsilon(z_i,z_j)\ge0$.
\end{lemma}
\begin{proof}
We use Frullani's integral $\log T=\int_0^\infty(e^{-u}-e^{-Tu})\,du/u$ for $T>0$. Together with $\sum_iy_i=0$ it gives
$\sum_{i,j}y_iy_j(-\log(|z_i-z_j|^2+\epsilon^2))=\int_0^\infty\frac{e^{-u\epsilon^2}}u\sum_{i,j}y_iy_je^{-u|z_i-z_j|^2}du$.
Gaussian kernels are positive definite. Indeed, expanding $e^{2u\langle z,z'\rangle}$ in powers and using the binomial
theorem writes the inner sum as a convergent sum of squares.
\end{proof}

A \emph{trial measure} is a probability measure $\sigma$, symmetric under conjugation, carried by finitely many
segments $\mathcal S_i\subset\Gamma\cup[-r_0,0]\cup[0,1]$ with rational endpoints. On each segment it has a density
$\sigma_i(y)\,dy$ in the parameter $y\in[0,1]$ of the segment. Here $\sigma_i$ is a polynomial with rational coefficients,
and $0\le\sigma_i\le\bar\sigma_i$ for a rational $\bar\sigma_i$. We write $\Theta_\sigma=\sum_i\bar\sigma_i/|\mathcal S_i|$
for the sum of the bounds of its densities with respect to arclength.

We compare the free points and the deleted nodes with $|\nu_c|\sigma$ through Lemma~\ref{lem:cpd}. This bounds the
free--free energy from below by a sum of one-point terms and the constant $|\nu_c|^2\cI(\sigma)$. Riemann-sum comparisons
replace the node sums by the potential $U^{\nu_c}$ and the node self-energy by $\cI(\nu_c)$, up to $O(r\log r)$. With the
external field one arrives at the function
\[
 \Upsilon(z)=|\nu_c|\,U^\sigma(z)-U^{\nu_c}(z)+(1-\varepsilon_0)\pi|\Im z|
\]
for a small fixed $\varepsilon_0>0$. Here $U^\sigma(z)=\int\log\frac1{|z-w|}\,d\sigma(w)$ is the logarithmic potential of
$\sigma$, and $U^{\nu_c}$ is that of $\nu_c$. A \emph{certificate} consists of $\sigma$ and rational numbers
$\Upsilon_\Gamma$, $\Upsilon_0$, $J$, $\cI_\sigma$, $\cI_\nu$ with the following properties.
\begin{enumerate}[label=(\alph*),leftmargin=*]
\item $\Upsilon\ge\Upsilon_\Gamma$ on $\Gamma$.
\item $\Upsilon\ge\Upsilon_0$ on $[0,1]$, the scaled position of the deleted common nodes.
\item $J\le\int_{-r_0}^0\min(\Upsilon(x)-\Upsilon_\Gamma,0)\,dx$.
\item $\cI(\sigma)\le\cI_\sigma$ and $\cI(\nu_c)\ge\cI_\nu$.
\end{enumerate}

\begin{theorem}[decay on a neighbourhood]\label{thm:decay}
Let $0<c\le1$. Let $F$ be a function of \eqref{eq:frames} whose poles are separated from $x^*$. Let
$(\sigma,\Upsilon_\Gamma,\Upsilon_0,J,\cI_\sigma,\cI_\nu)$ be a certificate for $\Gamma$ and $\nu_c$. Let $\epsilon>0$,
$\delta_\epsilon\ge\pi\epsilon\,\Theta_\sigma$, $\delta_{\rm m}>0$ and $\delta_J>0$ be rational, and put
\[
 \delta=3|\nu_c|^2\delta_\epsilon+4|\nu_c|^2\delta_{\rm m}+2|\mathcal T|\,\delta_J .
\]
Let $D$ and $\rho$ be rational numbers with
\[
 D\le2(|\nu_c|\Upsilon_\Gamma+|\mathcal T|J)-|\nu_c|^2\cI_\sigma+\cI_\nu-\delta,\qquad
 \rho\ge2(\Upsilon_\Gamma-\Upsilon_0).
\]
Then for every depth $K$, every member of the alternating geometry of parameter $c$ and every real $h$ with
$|h|\le e^{-\rho r}$,
\[
 |P_{r,I}(\xi+h)|\le\exp\bigl(-(D-o(1))r^2\bigr).
\]
\end{theorem}

The loss $\delta$ comes from replacing $\sigma$ by finitely many atoms in the comparison with the free points. The
regularisation radius $\epsilon$ of Lemma~\ref{lem:cpd} and the mesh of the atoms lower $\Upsilon_\Gamma$ by
$|\nu_c|(\delta_\epsilon+\delta_{\rm m})$. They raise $\cI_\sigma$ by $\delta_\epsilon+2\delta_{\rm m}$. Darboux sums for the
integral in (c) lower $J$ by $\delta_J$. Inserted in the expression for $D$, these changes give $\delta$. The tolerances
$\delta_{\rm m}$ and $\delta_J$ may be any positive rational numbers, since the atoms can be taken as fine as needed. In
the certificates of this paper $\epsilon=10^{-12}$ and $\delta<10^{-9}$.

The node energy $\cI(\nu_c)$ is a finite rational combination of the values $\mathcal H(y_k)$, where
$\mathcal H(y)=\frac12y^2\log|y|-\frac34y^2$ and the $y_k$ are the rational differences of the endpoints of the two pieces
of $\nu_c$. The lower bound $\cI_\nu$ is checked with rational enclosures of the logarithms. These come from the series of
$\operatorname{artanh}$ and a reduction to $[1,2)$. The energy $\cI(\sigma)$ is a finite sum of pair energies of
polynomial densities on segments. Each is a closed form of the type polynomial times logarithm plus arctangent, and is
bounded in the same way. The conditions (a)--(c) are statements about $\Upsilon$ on sets of positive measure. They are
proved on a cover by cells.

\subsection{The cover by cells}\label{ssec:slope-cert}
Each of the conditions (a)--(c) is checked on a finite cover by cells with exact rational arithmetic. On a cell, the
potential $U^\sigma$ is split into the contribution of the segments of $\sigma$ near the cell and the rest. A near segment
is bounded by a Taylor polynomial in the cell variable with an explicit remainder. For a cell that meets the line of a
segment beyond its end, the same bound is used for the segment with its orientation reversed. At a corner of $\Gamma$ the
two incident segments are treated together. Far segments are bounded by a multipole expansion with an explicit tail. The
potential $U^{\nu_c}$ is an explicit elementary function, and it is bounded above by monotonicity on the cell.

Some cells end at a breakpoint of $\sigma$, that is, at an endpoint of one of its segments or at a corner of $\Gamma$.
Let $s$ be the distance from the breakpoint along the cell, with $0<s\le S$. The potential of the segments that meet at
the breakpoint is then a polynomial in $s$ plus a term $\sum_mb_ms^m\log s$, up to an explicit remainder. In most cases
the logarithmic part is bounded in absolute value by $\sum_m|b_m|S^m|\log S|$. For the cells on $\Gamma$ of the shifts
$\frac k7$ we use a sharper bound. For $0<t\le1$ and an order $K\ge0$ write
\[
 -\log t=T_K(t)+R_K(t),\qquad T_K(t)=\sum_{k=1}^K\frac{(1-t)^k}k .
\]
Then $0\le R_K(t)\le(1-t)^{K+1}/((K+1)t)$, and the weighted inequality of arithmetic and geometric means gives
\[
 0\le t^mR_K(t)\le\frac{(m-1)^{m-1}(K+1)^K}{(m+K)^{m+K}}\qquad(m\ge1),
\]
with $0^0=1$. We put $t=s/S$ and write $\log s=\log S+\log t$. This bounds each term $b_ms^m\log s$ from below by a
polynomial of degree $m+K$ in $s$, minus a constant. Here $\log S$ is replaced by a rational number within $10^{-11}$ of
it, checked with the enclosures of the logarithm described above. The condition on the cell becomes an inequality for
one polynomial.

For (a) the cells are intervals of the arclength on each edge of $\Gamma$. The bound is proved on a horizontal
neighbourhood of the edge of fixed width. Beyond the last vertex the ray is handled by an analytic tail bound. For (b)
the cells are intervals of $[0,1]$. For (c) the integrand is bounded above on each cell of $[-r_0,0]$ by a constant or by
a chord. Then $J$ is the sum of the integrals of these majorants.

\begin{proposition}[certified constants]\label{prop:cert}
For each row of Table~\ref{tab:data} there is a trial measure $\sigma$ on a polygon $\Gamma$, with a cover by cells
satisfying the checks above, for which Theorem~\ref{thm:decay} gives the constants $D$ and $\rho$ of the table. For the
second plane the certificate is that of the reference frame of \S\ref{ssec:decaydouble}.
\end{proposition}

All inequalities are verified by the Lean kernel, as described in Section~\ref{sec:lean}. The certificate data are the
polygon, the density $\sigma$ and the cells. They were computed with floating-point optimisation and interval arithmetic
in Arb \cite{Arb}. Lean checks them with rational arithmetic only. The data play no role in the proof beyond satisfying these
checks.

\begin{remark}[how good a certificate can be]\label{rem:certlimit}
For an exponent $\mu_0$ the certificate meets $\mu_0$ exactly when $\mu_0(D-\Lambda)-\rho\ge0$. This quantity is a concave function
of $\sigma$. Its maximum over all probability measures on $\Gamma\cup[-r_0,0]\cup[0,1]$ is attained at a measure
$\sigma^*$. For this measure $|\nu_c|\sigma^*$ is the weighted equilibrium measure of the constrained energy problem on
this set, with the external field $(1-\varepsilon_0)\pi|\Im z|-U^{\nu_c}$ and the upper constraint $|\mathcal T|\,dx$ on
the segment \cite{SaffTotik}. Weak duality expresses the deficit of any certificate as a sum of nonnegative terms. The
first is the energy of $2\mu_0|\nu_c|(\sigma-\sigma^*)$. This is the difference between $2\mu_0|\nu_c|\sigma$ and the optimal
dual measure $2\mu_0|\nu_c|\sigma^*$. So a certificate close to the best possible one must be close to $\sigma^*$ in energy.
Then $\Upsilon-\Upsilon_\Gamma$ is a small ripple on the support of $\sigma^*$. This is why a pointwise verification on a
cover is needed. This remark is not used in any proof and is not part of the formalization.
\end{remark}

\subsection{Node removal and the neighbourhood}
The argument applies to every subset $C''\subseteq C'$ of the common nodes. Removing common nodes turns them into points
of the configuration that are not integrated over. Only common nodes can be removed. A node removed from one family only
would leave an uncancelled term of weight $\frac12$. By condition (b) one obtains
\[
 |\det\mathcal M_{C''}|\le\exp\bigl(-(D-o(1))r^2+\rho r\,|C'\setminus C''|\bigr),
\]
uniformly in $C''$. Insert this into \eqref{eq:expand} with $|h|e^{\rho r}\le1$. The sum has at most
$2^{|C'|}=e^{O(r)}$ terms, each at most $e^{-(D-o(1))r^2}$. This gives Theorem~\ref{thm:decay}. The deletion of at most
$K$ common nodes in a member changes the energy by $O(r)$.

\subsection{Frames with several terms and double poles}\label{ssec:decaydouble}
For a frame of \S\ref{ssec:frames} with all orders at most $1$ the argument above applies verbatim, with $|\mathcal T|=|S|$.
It reads the frame only through the source bound $|E(z)|\le C_Ee^{-2\pi|\Im z|}$ on $\Gamma$, through the positions of
the poles and through the bound on the atom masses. The weights enter only through $C_E$ and the atom masses. A change of
these constants multiplies each term of Andr\'eief's integral by at most $e^{O(r)}$, which is $e^{o(r^2)}$. So a
certificate for a set of shifts serves every frame on those shifts.

This includes the frames with a shift of order $0$. The heads of such a shift carry atoms of mass zero. The integrand
of the bound for $|\det\Lw_F[R',Q']|$ is nonnegative, and the integral over $|\Omega|$ only loses the terms in which a
free point sits at one of these heads. So the bound obtained with all $|S|$ shifts still holds. The construction itself
is unchanged. Only the number $\xi$ loses the term of that shift. We use this for the first plane in
Section~\ref{sec:plane1}, and for the second plane at the points $a\zeta(3)$ and $b\pi^2$.

At a double pole the Gram representation carries a derivative. Write
\[
 E(z)=-\pi\sum_sa_s\cot\pi(z+s)+\pi^2\sum_s\frac{b_s}{\sin^2\pi(z+s)},
\]
and let $\Gamma$ cross the negative axis at $x^*$, with heads $x_{s,n}=-s-n>x^*$.

\begin{lemma}[Gram formula]\label{lem:gramjet}
Let $\kappa_{t,t'}(z)=1/((z-t)(z-t'))$ for nonnegative integers $t,t'$. Then
\[
 \Lw_F[t,t']=\frac1{2\pi i}\int_\Gamma\kappa_{t,t'}E\,dz-\sum_{s,n}\bigl(b_s\kappa_{t,t'}'(x_{s,n})-a_s\kappa_{t,t'}(x_{s,n})\bigr).
\]
\end{lemma}
\begin{proof}
Put $\tilde F(z)=\sum_s(a_s\psi(1-s-z)-b_s\psi'(1-s-z))$. The reflection formulas $\psi(1-w)=\psi(w)+\pi\cot\pi w$ and
$\psi'(1-w)=\pi^2/\sin^2\pi w-\psi'(w)$ give $F=\tilde F+E$. The residues of $F\kappa_{t,t'}$ at $t$ and $t'$ give the
L\"owner entry, also for $t=t'$. At a head, $\psi$ has residue $-1$ and $\psi'$ has principal part $(z-x)^{-2}$, which
gives the atom terms. The function $\tilde F$ is holomorphic to the left of $\Gamma$ and $O(1/|z|)$ there by the balance,
so its integral against $\kappa_{t,t'}$ vanishes. The function $E$ decays like $e^{-2\pi|\Im z|}$, again by the
balance.
\end{proof}

The atom term at a double head is a symmetric bilinear form of rank two in the row and column functions, and it splits
into two rank-one terms. For functions $g$ and $h$ put $\hat g=g'-\frac{a_s}{2b_s}g$ and define $\hat h$ likewise. Then
\begin{equation}\label{eq:rankone}
 -b_s(gh)'+a_sgh=-\frac{b_s}2\bigl[(g+\hat g)(h+\hat h)-(g-\hat g)(h-\hat h)\bigr].
\end{equation}
So the measure $\Omega$ of \eqref{eq:gram} is replaced by a real measure on a parameter space with two labels at each
double head, carrying the weights $-b_s/2$ and $b_s/2$, and one label at each simple head. Andr\'eief's identity holds on any
measure space, so it writes every minor of $\Lw_F$ as an integral of products of two determinants. A repeated label
gives two equal rows, so its term vanishes.

The determinants now contain derivatives of Cauchy kernels at the double heads. We compare them with a frame that has
only simple poles. A \emph{reference frame} for $F$ is a frame $\mathcal R$ with simple poles on a set of shifts
$S_{\mathcal R}$ with the following properties. It contains every shift of $F$ of positive order. For each shift $s$ of
order $2$ it contains a further shift $s^*$, whose heads lie on the same side of $x^*$ as those of $s$, at a distance
$\delta_s$ from them with $\frac14\le\delta_s\le\frac12$. It may contain further shifts. The poles of $S_{\mathcal R}$ are
separated from $x^*$. For every frame $F^{(2)}_{v,u}$ we take the shifts $1$, $\frac34$, $\frac12$, with heads $-n$,
$-n+\frac14$, $\frac12-n$ for $n\ge1$, and $\delta_1=\frac14$. This also serves $v=0$, where all orders are at most
$1$, and $u=0$, where the shift $\frac12$ has order $0$. In the separate proof for $\zeta(3)$ in
\S\ref{ssec:zeta3} we take for $F^{(2)}_{1,0}$ the shifts $1$, $\frac12$ of the function $F_{1/2}$, with
$\delta_1=\frac12$.

Each label of Andr\'eief's integral now has a \emph{reference point}. For a point of $\Gamma$ it is the point itself.
For a simple head it is the head, which is a head of $\mathcal R$. For the two labels at a head $x_{s,n}$ of a double
shift they are $x_{s,n}$ and the head $x_{s^*,n}$ of $s^*$. Distinct labels have distinct reference points. In our
frames the only shift of order $2$ is $1$, so the double heads are the points $-n$ with $n\ge1$.

\begin{lemma}[jets]\label{lem:jet}
Let $k=O(r)$. Let $t_1,\dots,t_k$ be distinct integers in $[0,\bar N)$, and let $z_1,\dots,z_k$ be an injective list of
labels. Each label is a point of $\Gamma$ or a head of $F$ in $[x^*,0]$, and each double head may occur with either of
its two labels. Let $\phi_l$ be the function of \eqref{eq:rankone} attached to $z_l$, that is $\upsilon_t(z_l)$ for a
point of $\Gamma$ or a simple head, and $(g\pm\hat g)(z_l)$ with $g=\upsilon_t$ at a double head. Let $x^0_l$ be the
reference point of $z_l$. Then
\[
 \bigl|\det[\phi_l(t_i)]\bigr|\le e^{O(r\log r)}\,\bigl|\det[(x^0_l-t_i)^{-1}]\bigr| ,
\]
where the implied constant depends only on $\Gamma$, on $\mathcal R$ and on the weights of $F$.
\end{lemma}
\begin{proof}
A row at a double head is $\pm(\kappa_l\upsilon_t+\upsilon_t')$ at the head, with a constant $\kappa_l$. Give each
double-head label its own variable $y_l$, and keep the other labels fixed at their points. By multilinearity the
determinant is $\prod_l(\pm(\kappa_l+\partial_{y_l}))$, over the double-head labels, applied to the Cauchy determinant
$\det[(y_l-t_i)^{-1}]$ and evaluated at the heads. The double heads are at distance at least $1$ from the nodes. So
Cauchy's estimate on circles of radius $\frac18$, one variable at a time, bounds the determinant by
$\prod_l(|\kappa_l|+8)=e^{O(r)}$ times the supremum of the Cauchy determinant over the polydisc. On the polydisc a
double-head label is within $e_l$ of its reference point, where $e_l=\frac18$ for the label at $x_{s,n}$ and
$e_l=\delta_s+\frac18$ for the label at $x_{s^*,n}$. So $e_l\le\frac38$ for $\delta_s=\frac14$ and $e_l\le\frac58$ for
$\delta_s=\frac12$. All other labels have $e_l=0$.

Compare each factor of Cauchy's product formula
\[
 \bigl|\det[(y_l-t_i)^{-1}]\bigr|=\frac{\prod_{l<l'}|y_l-y_{l'}|\prod_{i<i'}|t_i-t_{i'}|}{\prod_{l,i}|y_l-t_i|}
\]
with its value at the reference points. The two labels of one double head stay at distance at most $\frac14$, and
their reference points are at distance $\delta_s\ge\frac14$. So their factor does not increase. For any other pair of
labels put $d_{ll'}=|x^0_l-x^0_{l'}|>0$. Then $|y_l-y_{l'}|\le d_{ll'}(1+e_l/d_{ll'})(1+e_{l'}/d_{ll'})$. For a node
$t_i$ we have $|x^0_l-t_i|\le|y_l-t_i|(1+e_l/|y_l-t_i|)$, and $|y_l-t_i|\ge t_i+\frac78$. Hence the logarithm of the
ratio is at most the sum of $e_l/d_{ll'}$ over pairs, plus the sum of $e_l/(t_i+\frac78)$ over the moved labels $l$ and
the nodes. Fix a moved label $l$. The reference heads are distinct points of a lattice $\frac1{q_{\mathcal R}}\Z$ in an
interval of length $O(r)$. So the sum of $1/d_{ll'}$ over head labels $l'$ is $O(\log r)$, and so is the sum over the
nodes. Now fix a label at a point $z$ of $\Gamma$. The heads lie at distance at least a constant $\delta_0>0$ from
$\Gamma$, by the separation of the poles from $x^*$. So the sum of $e_l/|x^0_l-z|$ over the moved labels $l$ is again
$O(\log r)$. This is the case of mixed configurations, of heads with free points on $\Gamma$. A bound by a constant
factor for each such pair would give $e^{O(r^2)}$, so the logarithmic sum matters here. In total the exponent is
$O(k\log r)=O(r\log r)$.
\end{proof}

\begin{theorem}[decay for frames]\label{thm:decayframe}
Let $F$ be a frame and $\mathcal R$ a reference frame for it. If all orders of $F$ are at most $1$, $\mathcal R$ may be
$F$ itself. Suppose that $\mathcal R$ has a certificate for $\Gamma$ and $\nu_c$ as in Theorem~\ref{thm:decay}, with
$|\mathcal T|$ the number of shifts of $\mathcal R$. Then the conclusion of Theorem~\ref{thm:decay} holds for every
member of the alternating geometry of parameter $c$ for $F$, with the same $D$ and $\rho$.
\end{theorem}
\begin{proof}
By Lemma~\ref{lem:gramjet}, \eqref{eq:rankone} and Andr\'eief's identity, each minor is an integral over the labels of
products of two determinants, one for the rows and one for the columns. Lemma~\ref{lem:jet} bounds each of them by the
Cauchy determinant at the reference points, at a cost $e^{O(r\log r)}=e^{o(r^2)}$. The label weights are bounded and
$|E|\le C_Ee^{-2\pi|\Im z|}$. Distinct labels have distinct reference points, all of them heads of $\mathcal R$ or points
of $\Gamma$. So the integral is at most $e^{o(r^2)}$ times the integral of the energy argument above for $\mathcal R$, in
which every head of $\mathcal R$ carries an atom of bounded mass. Those are the only inputs of that argument.
\end{proof}

The certificate of the second plane is that of the reference frame $(1,\frac34,\frac12)$, with $|\mathcal T|=3$. It serves
every frame $F^{(2)}_{v,u}$, including $v=0$ and $u=0$. The certificate of the separate proof for $\zeta(3)$ is that of
the function $F_{1/2}$, with $|\mathcal T|=2$. It is the certificate of the row for $\pi^2$, unchanged.

%% =====================================================================
\section{The first plane}\label{sec:plane1}

In this section we prove Theorem~\ref{thm:plane1} and parts (iii)--(v) of Theorem~\ref{thm:points1}. We use the frames
$F^\pm_{a,b,c}$ of \S\ref{ssec:frames} with the alternating geometry of parameter $c=\frac16$, so $\alpha=\frac1{12}$ and
$m=\lceil r/12-\frac12\rceil$. The point of the section is that the denominator exponent and the decay certificate do not
depend on the weights $(a,b,c)$, and that a weight zero is allowed. Nonvanishing is Theorem~\ref{thm:nvframe} with
$d\le3$, so one deleted node suffices.

\subsection{Polar data}
Fix a frame $F^\pm_{a,b,c}$ with integer weights. Let $\varrho=2$ on $S^+$ and $\varrho=1$ on $S^-$. For $j\ge1$ the increment
is
\[
 \delta(j-1)=\frac{3a}{3j-\varrho}+\frac{2b}{2j-1}+\frac cj ,
\]
and $-\delta'(j-1)=\frac{9a}{(3j-\varrho)^2}+\frac{4b}{(2j-1)^2}+\frac c{j^2}$. Let $p\ge5$ be a prime with $3\bar N<p^2$. Call an
integer $t$ \emph{singular} at $p$ if $p$ divides $t(3t-\varrho)(2t-1)$. The singular integers form the three lattices
$p\Z$, $k_A+p\Z$ and $k_B+p\Z$, where $k_A\in[1,p)$ solves $3k_A\equiv\varrho\pmod p$ and $k_B=\frac{p+1}2$. We write
$k_p=\min(k_A,k_B)$ for the first singular index. Then $k_A/p\to\kappa_A$ with
\[
 \kappa_A=\tfrac13\ \text{on } S^+,\ p\equiv1\ (6)\ \text{and on } S^-,\ p\equiv5\ (6),\qquad
 \kappa_A=\tfrac23\ \text{otherwise},
\]
and $\kappa_B=\frac12$. Passing from $S^+$ to $S^-$ exchanges the two classes of primes modulo $6$.

\begin{lemma}[polar data]\label{lem:poledata}
For every frame on $S^\pm$ with integer weights and every $1\le j\le\bar N$, the increment above has $p$-adic valuation at
least $-1$, and its derivative at least $-2$. If $j$ is not singular, both are $p$-integral.
\end{lemma}
\begin{proof}
Each of $j$, $3j-\varrho$ and $2j-1$ is smaller than $3\bar N<p^2$. So its valuation is at most $1$, and it is $0$ if $j$ is
not singular. The weights are integers.
\end{proof}

The lemma uses nothing about the weights but their integrality. If $p$ divides a weight, the corresponding pole may
disappear, and the lemma still holds since it gives only lower bounds. A weight zero, as for the frames with $a=0$, is
the case where $p$ divides the weight for every $p$. The singular indices of that shift are still counted, which only
weakens the bounds below. Rational weights are cleared first. Multiplying a
frame by a nonzero integer multiplies $\xi$ by that integer, which does not change $\mu$.

\subsection{Local bounds for every frame}
For $x\ne y$ put
\[
 w(x,y)=\mathbf 1[\text{a singular index lies in }(\min(x,y),\max(x,y)]]+\mathbf 1[p\mid x-y],
\]
and put $w(t,t)=2\cdot\mathbf 1[t\ge k_p]$. By Lemma~\ref{lem:poledata}, the entry of $\Phi_{r,I}$ at $(x,y)$ has Gauss valuation at
least $-w(x,y)$. Lemma~\ref{lem:trop} then applies with $W=w$. We use four facts about $w$: $w\le2$, $w(x,y)\le1$ if
$x\ne y$ and $p\nmid x-y$, $w(x,y)=0$ if $x\ne y$ are both below $k_p$, and $w(t,t)=0$ if $t<k_p$.

\begin{proposition}\label{prop:planelocal}
Let $p\ge5$ be a prime with $3\bar N<p^2$. For every frame on $S^\pm$ and every member, $e_p(r,I)$ is at most each of
the following, under its condition.
\begin{center}
{\renewcommand{\arraystretch}{1.2}
\begin{tabular}{@{}lp{0.55\textwidth}p{0.25\textwidth}@{}}
\toprule
 & bound & condition\\
\midrule
(T1) & $2r+2m$ & none\\
(T1$'$) & $2r+2m-\|m\|_p$ & none\\
(D2) & $2r-2(k_p+p-r)_++2m+(\min(\bar N,p+k_p)-r)_+$ & $2m<p\le r$\\
(D4) & $2(r-\min(k_p,r))+2m+(\bar N-p)_+$ & $p>r$\\
(D6) & $(\bar N-k_p)_+$ & $p\ge\bar N$, $k_p\ge r$\\
(W) & $2r+2m-(t_2-t_1)$ for each window of Theorem~\ref{thm:window} & as there\\
\bottomrule
\end{tabular}}
\end{center}
In (W) the cuts are the three lattices of singular indices and the residue of $r$.
\end{proposition}
\begin{proof}
(T1) and (T1$'$) are proved as in \S\ref{sssec:trop}. For the others we give potentials, with one function $\pi$ used on
rows and on columns. For (D2), $\pi=0$ on the common nodes $t<k_p$ with $t+p\ge r$, $\pi=1$ on the other common nodes,
and $\pi(t)=1+\mathbf 1[t-p<k_p]$ on the border. For (D4), $\pi(t)=\mathbf 1[t\ge k_p]$ on the common nodes and
$\pi(t)=1+\mathbf 1[t\ge p]$ on the border. For (D6), $\pi=0$ on the common nodes and $\pi(t)=\mathbf 1[t\ge k_p]$ on the
border. In each case the four facts give $w\le\pi\oplus\pi$ entrywise, and the sums are the stated values. For (W), the
proof of Theorem~\ref{thm:window} uses the function only through Lemma~\ref{lem:poledata} and the cuts. Between two
consecutive cuts the step function $V$ and the diagonal residues are constant, which is what the translation of
kernel vectors needs. No relation among the weights is used.
\end{proof}

None of these bounds reads the weights. They depend on $(r,m,p)$ and on the class of $p$ modulo $6$ only.

\subsection{The denominator exponent}
Homogenise as in \S\ref{ssec:profiles}, with $k_A\mapsto\kappa_A\theta$, $k_B\mapsto\kappa_B\theta$ and $\theta=p/r$. Let
$\varphi_1$ and $\varphi_2$ be the minima of the homogenised bounds of Proposition~\ref{prop:planelocal} for
$\kappa_A=\frac13$ and $\kappa_A=\frac23$, with the windows for $n\le32$. Put $t_0=\alpha/33$ and
\[
 \Lambda^*=(2+2\alpha)\,t_0+\frac12\int_{t_0}^{9}(\varphi_1+\varphi_2)\,d\theta .
\]
On the strips $(\frac\alpha{n+1},\frac\alpha n]$ with $1\le n\le32$ only (T1), (T1$'$) and the windows with
$n=\lfloor m/p\rfloor$ act. There the profile has the closed form described after Theorem~\ref{thm:planecontent}. On
$(\alpha,9]$ it is piecewise affine with finitely many pieces. So $\Lambda^*$ is a rational number, and
\[
 \Lambda^*=3.22274274961\ldots
\]
Primes of the classes $1$ and $5$ modulo $6$ each have density $\frac12$, and $S^-$ exchanges them. So $\Lambda^*$ serves
the frames on both shift sets $S^+$ and $S^-$. Above $\theta=\frac72$ both profiles vanish.

\begin{theorem}[denominators on the first plane]\label{thm:planecontent}
For every frame on $S^\pm$, every $\epsilon>0$ and every even $r\ge R(\epsilon)$, every member satisfies
$\log\den P_{r,I}\le(\Lambda^*+\epsilon)r^2$. The threshold $R(\epsilon)$ does not depend on the frame.
\end{theorem}
\begin{proof}
Take the product over primes of the bounds of Proposition~\ref{prop:planelocal}. For the primes with $p^2\le3\bar N$, and for
$2$ and $3$, every entry has valuation at least $-2\log(3\bar N)/\log p$. These primes contribute $O(r^{3/2})$. Each
bound is built from affine functions of $r$, $m$, $p$, $k_A$, $k_B$ and $p\lfloor r/p\rfloor$ by $\max$, $\min$ and
truncation, with bounded coefficients. So it differs from $r$ times its homogenisation by a bounded amount. Each
integer condition differs from its homogenised form only for $p$ in an interval of bounded length, where (T1) costs
$O(r\log r)$. The prime number theorem in progressions modulo $6$ then gives the claim, as in
Theorem~\ref{thm:content}. For the primes $p\le t_0r$ we use (T1), which gives $(2+2\alpha)t_0r^2$. On the strips
$\theta\le\alpha$ the content is an instance of Theorem~\ref{thm:periodic}, with $a=1$, $b=12$ and $K_0=2+2\alpha$. The
bounds acting there are closed under translation: (T1) and (T1$'$) grow by $\nu K_0bp$ under the translation of order
$\nu$, and it takes the window with $n=\lfloor m/p\rfloor$ to the window with $n+\nu$. So the table is checked on one
period of $1/\theta$, and its translates give the pieces of the other strips.
\end{proof}

On $\theta\le\alpha$ the profile has a closed form. Only (T1), (T1$'$) and the windows act there. Put $s=1/\theta$. Then
$\varphi(\theta)=2+2\alpha-\theta\,\Gamma_\kappa(s)$. Here $\Gamma_\kappa(s)$ is the larger of $\|\alpha s\|$ and the
longest gap of the lattice $\{0,\kappa_A,\kappa_B\}+\Z$ inside a window. The window is determined by the fractional
parts of $s$ and $\alpha s$. This function has period $12$ in $s$, as Theorem~\ref{thm:periodic} predicts. The table
has $82$ cells per class on one period.

\subsection{Decay}
Let $\Gamma$ be the polygon of the formal proof for the alternating geometry of parameter $\frac16$. It crosses the real
axis at a point $x^*$ with $x^*\equiv\frac14\pmod1$ for the frames on $S^+$ and $x^*\equiv\frac34\pmod1$ for those on
$S^-$. In both cases $x^*$ is at distance at least $\frac14$ from the poles of the shifts. The scaled cover is the same
for both.

\begin{proposition}\label{prop:planecert}
There is a trial measure with $25$ polynomial pieces on $\Gamma$, with a cover by cells satisfying the checks of
\S\ref{ssec:slope-cert} with $\varepsilon_0=10^{-6}$ and $|\mathcal T|=3$, for which Theorem~\ref{thm:decay} gives
\begin{gather*}
 D=\frac{99853629795114224867}{28125000000000000000}=3.55035\ldots,\\
 \rho=\frac{33845462710966783}{5\cdot10^{15}}=6.76909\ldots
\end{gather*}
\end{proposition}

The certificate contains no weight. By \S\ref{ssec:decaydouble} the decay estimate therefore holds for every frame on
$S^+$ and on $S^-$ with these $D$ and $\rho$, also when one weight is zero. The weights change only the threshold in $r$.

\subsection{Proof of Theorem~\ref{thm:plane1}}
Let $a$ and $b$ be rationals, not both zero. By Lemma~\ref{lem:planeconst} there is a frame on $S^+$ or $S^-$ whose
number $\xi$ is a nonzero rational multiple of $aL+b\pi^2$. For $a=0$ it is the frame $F^+_{0,3,-3}$, whose shift
$\frac13$ has weight zero. Theorem~\ref{thm:planecontent} gives the denominators with $\Lambda^*$, and
Proposition~\ref{prop:planecert} gives the decay. Both hold for every frame on $S^\pm$, including those with a weight
zero. Theorem~\ref{thm:nvframe} gives nonvanishing with one deleted node, since $d\le3$. For $a=0$ the frame has two
poles per period, and the argument is that of two digamma terms. The deduction of Theorem~\ref{thm:deduction} at the
even scales gives
\[
 \mu(\xi)\le\frac\rho{D-\Lambda^*}=20.66213\,7\ldots<\MuPlaneA .
\]
The strict inequality gives the bound $|aL+b\pi^2-p/q|>q^{-\MuPlaneA}$ for all large $q$, and in particular the
irrationality.\qed

The thresholds depend on $(a,b)$ through the height of the weights. No uniformity in $(a,b)$ is claimed.

\subsection{Level returns}\label{ssec:levels}
At special weights the partial sums of the increment return to an earlier level, and then more entries are
$p$-integral than Proposition~\ref{prop:planelocal} uses. Fix a frame on $S^\pm$, a class of primes modulo $6$, and a prime
$p\ge7$ with $3\bar N<p^2$ and $2m<p$. Every singular index $t\le\bar N$ has $3t-\varrho$, $2t-1$ or $t$ equal to $p$ times an
integer prime to $p$ and smaller than $p$. Hence
\[
 f_t=\frac{V(t)}p+(\text{a $p$-integer}),\qquad M[t,t]=\frac{U(t)}{p^2}+(\text{a $p$-integer}),
\]
where the \emph{level} $V(t)$ and $U(t)$ are rational numbers. They are sums of the terms $\frac{3a}{3l+n}$,
$\frac{2b}{2l+1}$, $\frac c{l+1}$, respectively of their squared analogues with a sign, over the singular indices up to
$t$. Here $n\in\{1,2\}$ is fixed by the class. So $V$ and $U$ are constant between consecutive singular indices, and they
depend on $p$ only through the class. Along the first period $V$ steps by the residues at the three poles. On $S^+$, for
$p\equiv1\pmod6$, the poles come in the order $k_A$, $k_B$, $p$ with residues $3a$, $2b$, $c$, so the level after two
poles is $3a+2b$, which is $0$ exactly on the line of $L$. For $p\equiv5\pmod6$ the order is $k_B$, $k_A$, $p$ with
residues $2b$, $\frac32a$, $c$. A frame for which a partial sum of these residues vanishes has a \emph{level return}.
The frames $(2,-1,-1)$, $(1,-2,1)$ and $(4,-3,-1)$ have level returns, with $b+\frac a2=0$, $2a+b=0$ and
$2b+\frac32a=0$.

\begin{lemma}[period sums]\label{lem:periodsums}
Let $t'=t+sp$ with $1\le s<p$. Then $f_{t'}-f_t=(V(t')-V(t))/p+\Delta$ with $v_p(\Delta)\ge1$.
\end{lemma}
\begin{proof}
The polar terms give $(V(t')-V(t))/p$ exactly. For each shift, every block of $p$ consecutive indices runs once
through the residues modulo $p$. The nonpolar terms of a block therefore contribute a multiple of
$\sum_{v\in\F_p^\times}v^{-1}\equiv0\pmod p$.
\end{proof}

So two congruent nodes at the same level have a $p$-integral entry, not one of valuation $-1$. A second saving comes
from row operations inside a residue class. Let $x$ be a border node and $a<x$ a common node with $x-a=jp$. Suppose that
$V(x)\ne V(a)$, $U(a)\ne0$, and that
\[
 \lambda_x=\frac{V(x)-V(a)}{j\,U(a)}
\]
satisfies $\lambda_x\ne1$ and has no prime factor larger than $5$ in its denominator. Then $\lambda_x$ is $p$-integral. We
call $(x,a)$ a \emph{pivot pair}, and we put $V^*(x)=(V(x)-\lambda_xV(a))/(1-\lambda_x)$. Pivot pairs whose border node is a
column are defined in the same way. Each border node lies in at most one pivot pair, and $V^*=V$ on the nodes in no
pair. Define weights $w'$ on rows times columns as follows.
\begin{enumerate}[label=(\roman*),leftmargin=*]
\item For $x\not\equiv y\pmod p$, $w'(x,y)=\mathbf 1[V^*(x)\ne V^*(y)]$.
\item For a pivot pair, the entry at the border node and its common node has $w'=0$.
\item For other $x\equiv y$, $w'(x,y)$ is the largest $w_0(x_1,y_1)$ with $x_1$ equal to $x$ or its common node and $y_1$
 equal to $y$ or its common node. Here $w_0(x_1,y_1)=2\cdot\mathbf 1[V(x_1)\ne V(y_1)]$ for $x_1\ne y_1$, and
 $w_0(t,t)=2\cdot\mathbf 1[U(t)\ne0]$.
\end{enumerate}

\begin{lemma}[levels]\label{lem:levels}
If nonnegative numbers $u$ on the rows and $u'$ on the columns satisfy $w'\le u+u'$ entrywise, then
$e_p(r,I)\le\sum u+\sum u'+2|I|$ for every member.
\end{lemma}
\begin{proof}
Subtract $\lambda_x$ times the row of $a$ from the row of $x$ for every pivot pair, and likewise for columns. The
pivots are common nodes, which are never modified themselves, so these operations are unimodular and the determinant is
unchanged as a polynomial in the variable. It suffices to bound the Gauss valuation of each new entry from below by
$-w'$. For congruent $x\ne y$ without pivots, Lemma~\ref{lem:periodsums} gives
$M[x,y]=(V(y)-V(x))/(sp^2)+\Delta/(sp)$, so the valuation is at least $-2\cdot\mathbf 1[V(x)\ne V(y)]$. With pivots the
entry is a combination of at most four such entries with $p$-integral coefficients, which gives (iii). For a pivot
pair the entry is $M[x,a]-\lambda_x(M[a,a]+x_{\rm var})$, where $x_{\rm var}$ is the variable. The terms of order $p^{-2}$
cancel by the choice of $\lambda_x$, the rest is $p$-integral, and the variable has the $p$-integral coefficient
$-\lambda_x$. This is (ii). For $y\not\equiv x$ write $\frac1{y-x}=\frac1{y-a}+\frac{jp}{(y-a)(y-x)}$. Then
\[
 M[x,y]-\lambda_xM[a,y]=\frac{(1-\lambda_x)(V(y)-V^*(x))}{p\,(y-a)}+(\text{a $p$-integer}),
\]
since $y-a$ and $y-x$ are units. If $y$ is pivoted too, the same identity applied to its column gives the leading term
$(1-\lambda_x)(1-\lambda_y)(V^*(y)-V^*(x))/(p(y-a))$. This is (i). Finally, let a common node be deleted. Since $2m<p$, each
residue class contains at most one border node. So the deleted node is the common node of at most one pivot pair. That
pair is undone. This raises the weights of the border line off its fibre by at most $1$, and does not raise its
weights in the fibre. Raising the potential of that line by $1$ restores $w'\le u+u'$. So each deleted node costs at
most $1$, within the allowance $2|I|$.
\end{proof}

The bound of Lemma~\ref{lem:levels} with no pivot pairs is already better than Lemma~\ref{lem:trop} at a level return.
The pivots add more. They are the assignment form of the rank of the reduction modulo $p$ on the fibres where the
level returns. For instance, take the frame $(2,-1,-1)$ on $S^+$ and $p\equiv5\pmod6$. The poles fall at about $\frac
p2$, $\frac{2p}3$, $p$, $\frac{3p}2$, with residues $-2$, $3$, $-1$, so $V$ returns to $0$ at $p$. A border node in
$[\frac{3p}2,\frac{5p}3)$ has level $-\frac23$. Its common node in $[\frac p2,\frac{2p}3)$ has level $-2$ and $U=4$. Then
$\lambda=\frac13$ and $V^*=0$, so after the row operation the border node behaves like a node at the returned level.

\begin{theorem}[the special points]\label{thm:dips}
For the frames $(2,-1,-1)$, $(1,-2,1)$ and $(4,-3,-1)$, on $S^+$ and on $S^-$, the minimum of the bounds of
Proposition~\ref{prop:planelocal} and Lemma~\ref{lem:levels} gives the denominator exponents
\[
 \Lambda=3.06683997\ldots,\qquad 3.10960617\ldots,\qquad 3.17760386\ldots
\]
\end{theorem}
\begin{proof}
On each cell of the common refinement of the conditions, the nodes split into boundedly many classes by residue
interval, parity and period. Pivot pairs, $\lambda$ and $V^*$ are constant on these classes. Potentials constant on the
classes are affine in $\theta=p/r$, with a bounded rounding error per prime, as in Theorem~\ref{thm:planecontent}. The
potentials of the formal proof take half-integer values. They are verified on every pair of classes, and the minimum with
the plane bounds is integrated as there. A cell is fixed by the integer parts of $r/p$ and $\bar N/p$
and by the order of the residues of $k_A$, $k_B$, $r$ and $\bar N$ in a period. At a common endpoint $x$ of two cells
the bounds of the two sides take the same value. A prime with $p/r$ near $x$ may lie in neither cell, through the
rounding of $m$ or a coincidence of two residues. So on the window $|p/r-x|\le10^{-4}x$, in each class of primes
modulo $6$, we use the largest of the bounds of all cells that occur in the window. There are at most three. That one of
them applies to each $(r,p,m)$ in the window is a statement of linear integer arithmetic in $r$, $p$, $m$, $k_A$ and
$k_B$, and it is proved for every prime $p$. This bound is the second bound of (C3) at $x$. All pivots used have $\lambda\in\Z[\frac1{30}]$, so every prime $p\ge7$ is
admissible. The level branches act only for $\theta\ge\frac3{10}$. The frames on $S^-$ give the same profiles with the
classes exchanged.
\end{proof}

\begin{proof}[Proof of Theorem~\ref{thm:points1}(iii)--(v)]
The frame $(2,-1,-1)$ has $\tau=\frac2{27}$. On $S^+$ its number is $9(L+\frac{2\pi^2}{27})=\frac12\Psi(\frac16,\frac13)$, and
on $S^-$ it is $-9(L-\frac{2\pi^2}{27})=-\frac12\Psi(\frac23,\frac56)$. The frame $(1,-2,1)$ has $\tau=-\frac1{27}$ and gives
$\Psi(\frac12,\frac23)$ on $S^+$ and $\Psi(\frac13,\frac12)$ on $S^-$, up to sign. The frame $(4,-3,-1)$ has $\tau=\frac1{18}$ and
gives $L\pm\frac{\pi^2}{18}$. With the denominator exponents of Theorem~\ref{thm:dips}, the decay of
Proposition~\ref{prop:planecert} and Theorem~\ref{thm:nvframe}, the deduction gives
\[
 \frac\rho{D-\Lambda}=13.99986\,39\ldots,\qquad 15.35829\,31\ldots,\qquad 18.15999\,83\ldots,
\]
below $\MuDipA$, $\MuDipB$ and $\MuDipC$.
\end{proof}

The point $t=0$ is the frame $(2,-3,1)$, which has a level return after two poles in the class $p\equiv1\pmod6$. It is
treated with its own geometry in Section~\ref{sec:chi3-quarter}. The points $t=\pm\frac19$ are given by the
differences $F_{1/3}$ and $F_{2/3}$ of two digamma functions, with their own geometry, in Section~\ref{sec:proof}.

%% =====================================================================
\section{The point \texorpdfstring{$\Psi(\frac13,\frac23)$}{Psi(1/3,2/3)} of the first plane}
\label{sec:chi3-quarter}

For $\Psi(\frac13,\frac23)$ we use the function $F_\chi$ of \eqref{eq:frames}, which combines three digamma terms, with
the alternating geometry of parameter $c=\frac12$ at all even scales $r$. Then $\alpha=\frac14$ and
$m=\lceil r/4-\frac12\rceil$. The decay and nonvanishing estimates of Theorems~\ref{thm:decay} and~\ref{thm:selection}
apply unchanged. The denominators need more care. At a prime $p$ the increments $g(k)$ of \S\ref{sec:construction} have
poles in three residue classes instead of two, so the local bounds depend on two first poles. This section describes the
two poles, the bounds we use, and the resulting constant.

\subsection{The first poles}
Let $p>3$ be a prime. The increment $g(k)=\frac1{3k-2}-\frac1{2k-1}+\frac1{6k}$ has a pole modulo $p$ exactly when $p$
divides $k$, $3k-2$ or $2k-1$. Let $k_A$ and $k_B$ be the least $k\ge1$ with $p\mid3k-2$ and with $p\mid2k-1$. Then
\[
 k_B=\frac{p+1}2,\qquad
 k_A=\begin{cases}\frac{p+2}3,&p\equiv1\pmod6,\\[3pt]\frac{2p+2}3,&p\equiv5\pmod6.\end{cases}
\]
The residues of $p\,g(k)$ at $k_A$ and $k_B$ are $1$ and $-1$ if $p\equiv1\pmod6$, and $\frac12$ and $-1$ if
$p\equiv5\pmod6$. In the first case $k_A<k_B$ and the two residues cancel. So the partial sums $\sum_{k\le t}g(k)$ are
$p$-integral for $t<k_A$, and again for $k_B\le t<p$.

Both first poles are affine in one index. Define $k_p$ by $6k_p=jp+5$ with $0<j<6$. Then $j=1$ for $p\equiv1\pmod6$ and
$j=5$ for $p\equiv5\pmod6$, $k_p=\frac j6p+O(1)$, and
\[
 (k_A,k_B)=\begin{cases}(2k_p-1,\ 3k_p-2),&p\equiv1\pmod6,\\[2pt](2k_p-p-1,\ 3k_p-2p-2),&p\equiv5\pmod6.\end{cases}
\]
Thus $j$ plays the role of the root class of \S\ref{sec:content}, with $q$ replaced by $6$.

\subsection{The local bounds}
Fix a prime $p>6$ with $6\bar N<p^2$. Every pole below $\bar N$ is then simple. Call a node $t$ \emph{regular} if
$t<\min(k_A,k_B)$, or if $p\equiv1\pmod6$ and $t$ is a border node with $k_B\le t<p$. At a regular node the partial sum
of the increments up to $t$ is $p$-integral. When $t$ is a common node, its diagonal entry is $p$-integral too.

\begin{proposition}\label{prop:chi3-local}
Every member of the alternating geometry of parameter $\frac12$ for $F_\chi$ satisfies $e_p(r,I)\le\tau$ for each of the
following $\tau$.
\begin{enumerate}[label=(\roman*),leftmargin=*]
\item The bounds (T1) and (T1$'$), namely $2(r+m)$ and $2r+2m-\|m\|_p$.
\item If $\bar N\le2p$, the bound $2r+2m-n_{\rm reg}$. Here $n_{\rm reg}$ is the number of pairs formed by a regular node $t$
 and a side on which $t$ lies, row or column, such that $t+p$ is not a node of the other side.
\item If $\bar N<\min(k_A,k_B)$, the bound $0$.
\item The bound $2r+2m-(t_2-t_1)$ of Theorem~\ref{thm:window}, for every $n$ and every window $[t_1,t_2)$ as there, with
 the cuts $t\equiv0$, $t\equiv k_A$, $t\equiv k_B$ and $t\equiv r\pmod p$.
\end{enumerate}
\end{proposition}

\begin{proof}
The bounds (i) are proved as in \S\ref{sssec:trop}. For (ii) we apply Lemma~\ref{lem:trop} with the following potentials.
At a node that is not regular, the potential is $1$ on its row and on its column. At a regular node $t$, the potential on
its row is $1$ if $t+p$ is a column node, and $0$ otherwise. The potential on its column is $1$ if $t+p$ is a row node,
and $0$ otherwise. Their sum is $2(r+m)-n_{\rm reg}$, and they dominate the weights $W$. Indeed, a diagonal weight is $0$
at a regular common node and at most $2$ in general. Two nodes that are not congruent modulo $p$ have weight at most $1$,
since all poles are simple. Their weight is $0$ if both are regular. Two distinct congruent nodes have weight at most
$2$. Since $\bar N\le2p$ and regular nodes are smaller than $p$, at most one of them is regular. Suppose that the row node $t$
of such an entry is regular. Then its column node is $t+p$, so the potential on the row of $t$ is $1$. The case of a
regular column node is symmetric.

For (iii), the condition $\bar N<\min(k_A,k_B)$ implies $\bar N<p$. Then no index below $\bar N$ is a pole and no two nodes are
congruent modulo $p$, so every entry is $p$-integral. For (iv), the proof of Theorem~\ref{thm:window} uses $F$ only
through its poles. The step function $V$ and the diagonal residues are constant between consecutive cuts. This is what
the translation between the fibres of the window needs.
\end{proof}

The bounds (T2)--(T5), (LP1), (LP2), (R) and (S) of Theorem~\ref{thm:local} are not used for $F_\chi$.

\subsection{The denominators}
The tables for $\Lambda$ of \S\ref{ssec:checker} and Theorem~\ref{thm:content} apply to $F_\chi$ with $q$ replaced by $6$. The root
class is $j\in\{1,5\}$ as above, and the bounds are those of Proposition~\ref{prop:chi3-local}. These bounds are
expressions of the class $\mathcal E$ in $r,p,m,k_A,k_B$. On each class of primes modulo $6$ the substitution above makes
them expressions in $r,p,m,k_p$. Their homogenisations are taken with $\kappa=j/6$. The proof of
Theorem~\ref{thm:content} uses $F$ only through these bounds and the root class. It goes through with the bound
(iii) in place of (T5) beyond the end of the table.

The table of the formal proof starts at $\theta_0=\frac1{68}=\alpha/17$. Below it the trivial bound $\ell_0r$ with
$\ell_0=\frac52=2+2\alpha$ is used. On $(\alpha/17,\alpha]$ the table uses the bounds (i) and the windows (iv) with
$1\le n\le16$. Beyond $\alpha$ it uses (i), (ii), the windows (iv) with $n=0$ between consecutive cuts, and the bound
(iii), which ends the table. The table is admissible for the two classes $j=1,5$. The arithmetic estimate therefore
holds for every even $r$ and every member with $|I|\le64$, with the denominator exponent
\[
 \Lambda_\chi=4.09592776731\ldots
\]
This is an explicit rational number, whose numerator and denominator have $117$ and $116$ digits.

\subsection{Decay and nonvanishing}
The decay estimate is Theorem~\ref{thm:decay} for the function $F_\chi$, with $|\mathcal T|=3$ and $c=\frac12$. The node
measure $\nu_{1/2}$ has mass $\frac54$ and energy
\[
 \cI(\nu_{1/2})=\frac{75}{32}-\frac98\log\frac32-\frac1{16}\log2 .
\]
A trial measure on a polygon, with a cover by cells satisfying the checks of \S\ref{ssec:slope-cert}, gives
\[
\begin{gathered}
 D=\frac{4643511791526209475773}{10^{21}}=4.64351\ldots,\\ \rho=\frac{37521870783566737}{5\cdot10^{15}}=7.50437\ldots
\end{gathered}
\]
Nonvanishing is Theorem~\ref{thm:selection} for $F_\chi$, which has $d=3$. The pole of its increment at $-1$ is
isolated in its class modulo $\Z$, and the members with $|I|\le64$ are used. The deduction of Section~\ref{sec:proof}, at the even scales, gives
\[
 \mu\bigl(\Psi(\tfrac13,\tfrac23)\bigr)\le\frac\rho{D-\Lambda_\chi}=\MuChiExact\ldots<\MuChi,
\]
since $\Psi(\frac13,\frac23)=36\,\xi_\chi$.

%% =====================================================================
\section{The second plane}\label{sec:plane2}

In this section we prove Theorem~\ref{thm:plane2} and the bound for $\zeta(3)$ in Theorem~\ref{thm:ends}. We use the
frames $F^{(2)}_{v,u}$ of \S\ref{ssec:frames} with integers $v$ and $u$, not both zero, and we put
$\varepsilon=\mathbf 1[u\ne0]$. For Theorem~\ref{thm:plane2} we take the alternating geometry of parameter
$c=\frac6{41}$, so $\alpha=\frac3{41}$, for every frame. This includes $v=0$, the points $b\pi^2$, and $u=0$, the
points $a\zeta(3)$. For the separate bound for $\zeta(3)$ in Theorem~\ref{thm:ends} we take $c=\frac3{10}$. In both
cases $2m\le r$. The new feature is the double pole at the shift $1$. The entries of the L\"owner matrix then have
valuation down to $-3$, and most of the work goes into recovering what is lost against a simple pole.

\subsection{Cells and weights}
Let $p\ge5$ be a prime with $p^2>2\bar N+2p$. The increment is
\[
 \delta(k-1)=\frac uk-\frac{2u}{2k-1}-\frac v{k^2}\qquad(k\ge1).
\]
Its poles below $\bar N$ are the double poles $k=jp$ and, if $u\ne0$, the simple poles $k=jp+\frac{p+1}2$. Write an integer
$t\ge0$ as $t=jp+\omega$ with $0\le\omega<p$. We call $j$ the \emph{period} of $t$, $h=\mathbf 1[\omega\ge\frac{p+1}2]$ its
\emph{half}, and $c(t)=(j,h)$ its \emph{cell}. A pole lies at the start of every cell. Let $S(c)$ be the sum of the polar
terms of $\delta$ up to the start of the cell $c$, and $G(c)$ the same for $-\delta'$. Then
\[
 f_t=S(c(t))+\rho_t,\qquad e_t=G(c(t))+\gamma_t ,
\]
with $\rho_t$ and $\gamma_t$ in $\Z_{(p)}$, and $p^2S$ and $p^3G$ in $\Z_{(p)}$. Modulo $p$ the numbers $p^2S$ and $p^3G$
depend only on the period, because the simple poles have lower order than the double ones.

\begin{lemma}[weights]\label{lem:weights2}
Let $W(t,t')=\max(0,-v_p(M[t,t']))$. Then $W(t,t')\le3$ if $t\equiv t'\pmod p$, and for $t\not\equiv t'$
\[
 W(t,t')\le W_3(c(t),c(t')),\qquad W_3(c,c')=\begin{cases}0&c=c',\\ 1&\text{same period, other half},\\ 2&\text{other
 period.}\end{cases}
\]
On the diagonal $W(t,t)$ is $0$ on the cell $(0,0)$, at most $2\varepsilon$ on $(0,1)$, and at most $3$ elsewhere. This holds
also when $p$ divides $v$ or $u$.
\end{lemma}
\begin{proof}
For $t\ne t'$ the difference $t'-t$ is a unit unless $t\equiv t'$, and then $v_p(t'-t)=1$ since $p^2>\bar N$. The
numerator $f_{t'}-f_t$ has valuation at least $-2$, at least $-1$ if $t$ and $t'$ lie in one period, and at least $0$ if
they lie in one cell. The diagonal is $-e_t$. The weights $u$ and $v$ are integers, so a prime dividing them only
raises valuations.
\end{proof}

The primes that divide $v$ need one remark. If $p\mid v$, the double poles
contribute terms $v/k^2$ of valuation at least $-1$ to $f_t$, and terms $2v/k^3$ of valuation at least $-2$ to $e_t$.
So every entry of $M$ has valuation at least $-2$, and Lemma~\ref{lem:trop} with all potentials $1$ gives
$e_p(r,I)\le2(r+m)$. The bounds of Proposition~\ref{prop:rows2} below use Lemma~\ref{lem:weights2} only and hold for
all $p$. Those of Theorem~\ref{thm:rows2} assume $p\nmid v$. Each of them is at least $2(r+m)$, as we check after its
proof. So for $p\mid v$ the trivial bound $2(r+m)$ takes their place, and every bound used below holds for all frames
with $v\ne0$.

\begin{proposition}[assignment bounds]\label{prop:rows2}
Assume $2m\le r$. For every member,
\begin{center}
{\renewcommand{\arraystretch}{1.2}
\begin{tabular}{@{}lll@{}}
\toprule
 & bound & condition\\
\midrule
(P1) & $3(r+m)-\|m\|_p$ & none\\
(P4) & $6(r+m-p)+\varepsilon p$ & $p\le r$\\
(P5) & $\varepsilon r+3(\bar N-p)$ & $r\le p\le\bar N$\\
(P6) & $\varepsilon(2r+2m-p+1)$ & $\bar N\le p\le2r$\\
(P7) & $\varepsilon(\bar N-\frac{p-1}2)_+$ & $p\ge2r$\\
\bottomrule
\end{tabular}}
\end{center}
\end{proposition}
\begin{proof}
All follow from Lemmas~\ref{lem:trop} and~\ref{lem:weights2}. For (P1) take the potentials of (T1$'$) shifted by
$\frac12$: on the rows of a residue class $\frac32-\frac12\sigma$ and on its columns $\frac32+\frac12\sigma$, where
$\sigma=\operatorname{sign}(|X'\cap\text{class}|-|Y'\cap\text{class}|)$. Congruent pairs get $3$ and the others at least
$2$. For the other bounds take one function $\varphi$ on rows and columns, with $\varphi=0$ on the cell $(0,0)$,
$\varphi=\varepsilon$ on $(0,1)$ and $\varphi=3$ on the periods $j\ge1$. Pairs inside the period $0$ have weight at most
$\varphi\oplus\varphi$, and any pair with a node in a later period has weight at most $3$. Counting the common nodes twice
and the border nodes once gives the four bounds. The hypothesis $2m\le r$ enters twice. For (P5) it gives
$p\le\bar N\le2r$, so $(p-1)/2<r$. For (P7) it gives $\bar N\le2r\le p$, so all nodes lie in the period $0$ and the
common nodes lie in the cell $(0,0)$.
\end{proof}

Proposition~\ref{prop:rows2} bounds the contribution of a prime $p\le2m$ by $3(r+m)-\|m\|_p$, and that of a prime
$2m<p\le r$ by no better. The double
pole costs one more unit per node than a simple pole. The next two subsections recover most of this unit on the primes
below $r+2m$.

\subsection{Elimination of complete periods}
Let $\cO$ be the discrete valuation ring of \S\ref{sec:content}, in a variable $X$, and put $x=p^{-3}X$, so that a member
is a polynomial in $x$. Fix $\kappa\ge1$ and $k'\ge0$ with $(\kappa+2k')p\le\bar N+p$, and put $\bar N'=\max(\bar N,
(\kappa+2k')p)$. Then $p^2>2\bar N'$. Let $H=M+p^{-3}X\Id$ on $[0,\bar N')$, so that $H$ restricts to the matrices of the
members. Extend the border alternately, so $\hat X=\{t\ge r:t-r\text{ even}\}$ and $\hat Y=\{t\ge r:t-r\text{ odd}\}$. The
\emph{pivot} consists of the complete periods below $\kappa p$ and of $2k'$ further complete periods:
\[
 \mathcal P_r=[0,\kappa p)\cup(\hat X\cap[\kappa p,(\kappa+2k')p)),\qquad
 \mathcal P_c=[0,\kappa p)\cup(\hat Y\cap[\kappa p,(\kappa+2k')p)).
\]

\begin{lemma}[displacement and bordering]\label{lem:displace}
For all nodes $s,t$ one has $(s-t)H(s,t)=f_s-f_t$, whatever the diagonal. Suppose $H[\mathcal P_r,\mathcal P_c]$ is
invertible, and let $Z$ be its inverse. Let $\mathrm{Sw}$ be the Schur complement of $H[\mathcal P_r,\mathcal P_c]$ in the
matrix bordered by unit vectors, as in Lemma~\ref{lem:bordered}. Then
\begin{enumerate}[label=(\alph*),leftmargin=*]
\item every minor of $H$ that contains $\mathcal P_r\setminus A_0$ and $\mathcal P_c\setminus A_0'$ equals
 $\pm\det H[\mathcal P_r,\mathcal P_c]$ times a minor of $\mathrm{Sw}$ with rows $B_0\cup A_0'$ and columns
 $B_0'\cup A_0$, where $B_0$ and $B_0'$ are the remaining rows and columns.
\item there are vectors $g_u\in\Q(X)^2$, one for each label $u$, such that $(v-u)\,\mathrm{Sw}(u,v)=g_u\wedge g_v$ for all
 labels $u\ne v$.
\end{enumerate}
\end{lemma}
\begin{proof}
The first identity is the definition of the L\"owner matrix. With $\mathsf T=\operatorname{diag}(t)$ it says that
$\mathsf TH-H\mathsf T=f\mathbf 1^{\mathsf T}-\mathbf 1f^{\mathsf T}$ has rank two. The bordered matrix satisfies the same
equation with generators $(1,f_t)$ on actual nodes and $0$ on the border, since $(s-s)\cdot1=0$. A Schur complement of a
block preserves this displacement structure, with generators $g=G-VZG_{\mathcal P}$. Laplace expansion along the unit
border gives (a).
\end{proof}

The generators are explicit. On a label $z$ outside the pivot $g_z=(1-H(z,\mathcal P_c)Z\mathbf 1,\
f_z-H(z,\mathcal P_c)Zf)$, and on a pivot label $t$ it is $((Z\mathbf 1)_t,(Zf)_t)$ up to sign. We rescale them as
$\tilde g_u=(p^{-3}g_{u,1},p^{-1}g_{u,2})$ on pivot labels and $\tilde g_u=(g_{u,1},p^2g_{u,2})$ on the others, and the pivot
rows and columns of $\mathrm{Sw}$ by $p^{-3}$. The rescaled matrix $\widetilde{\mathrm{Sw}}$ then has entries
$p^{-2}(\tilde g_u\wedge\tilde g_v)/(v-u)$.

The reduction of the pivot is governed by matrices on period indices. For $h\in\{0,1\}$ put $s_h(j)=p^2S(j,h)$,
$g_h(j)=p^3G(j,h)$, $\Gamma_h(j)=(1,s_h(j))$, and
\[
 B^h(l,j)=\frac{s_h(j)-s_h(l)}{j-l}\ (l\ne j),\qquad B^h(j,j)=X-g_h(j).
\]
Modulo $p$ one has $s_h(j)\equiv-v\,h_j$, where $h_j=\sum_{i\le j}i^{-2}$. So the off-diagonal entries of $B^h$ vanish
modulo $p$ when $p\mid v$. This is why $p\nmid v$ is assumed from now on.
Let $J_X=\{\kappa,\kappa+2,\dots,\kappa+2k'-2\}$, $J_Y=J_X+1$, and $B_X^h=B^h[[0,\kappa)\cup J_X,[0,\kappa)\cup J_Y]$. Write
$\bar B_X$ for its reduction modulo $p$, which does not depend on $h$. If $t\equiv t'\pmod p$ then
$p^3H(t,t')\equiv B^h(j_t,j_{t'})\pmod{p^3}$. If $t\not\equiv t'$ then $p^3H(t,t')\in p\cO$. Hence
$p^3H[\mathcal P_r,\mathcal P_c]$ is congruent modulo $p$ to a direct sum of copies of $\bar B_X$ and its transpose, one for
each residue class.

\begin{theorem}[nondegeneracy]\label{thm:nd}
Let $\kappa\ge1$, $k'\ge0$, $p>\kappa+2k'$ and $p\nmid v$. Then $\det\bar B_X\ne0$ in $\F_p(X)$.
\end{theorem}
\begin{proof}
For $k'=0$ the determinant is monic in $X$. Let $k'\ge1$. The variable occurs only on the diagonal of the block
$[0,\kappa)$, so $\det\bar B_X$ is $\pm\det\bar B[J_X,J_Y]$ times a monic polynomial. Since $\bar B[J_X,J_Y]$ is $-v$
times the matrix $[(h_j-h_l)/(j-l)]_{l\in J_X,j\in J_Y}$ and $p\nmid v$, it suffices that the determinant $D$ of the
latter is nonzero in $\F_p$. Suppose that $c\ne0$
lies in its kernel. Let $P=\prod_{j\in J_Y}(z-j)$, $\mathsf u=\sum_jc_jP/(z-j)$ and $\mathsf w=\sum_jc_jh_jP/(z-j)$. These
polynomials over $\F_p$ have degree at most $k'-1$, $\mathsf u\ne0$, and $\mathsf w(t)=h_t\mathsf u(t)$ on the $2k'$
consecutive nodes $\kappa,\dots,\kappa+2k'-1$. Since $h_{t+1}-h_t=(t+1)^{-2}$, the polynomial
\[
 \mathsf E(z)=(z+1)^2\bigl(\mathsf w(z+1)\mathsf u(z)-\mathsf w(z)\mathsf u(z+1)\bigr)-\mathsf u(z)\mathsf u(z+1)
\]
vanishes at the $2k'-1$ nodes other than the last. The top terms cancel, so $\deg\mathsf E\le2k'-2$ and $\mathsf E=0$. Divide
$\mathsf u$ and $\mathsf w$ by their greatest common divisor. At $z=-1$ the identity gives $\mathsf u(-1)\mathsf u(0)=0$. If
$\mathsf u(x+1)=0$ with $x\ne-1$ then $\mathsf u(x)=0$, and conversely, by coprimality. So every element of $\F_p$ is a root
of $\mathsf u$, whose degree is less than $p$. This is a contradiction.
\end{proof}

The argument is that of Theorem~\ref{thm:nvframe}, in characteristic $p$. The double pole at $-1$ is isolated in its
orbit, and the orbit is now finite, which replaces the infinitely many roots by $p$ roots.

\begin{proposition}[generators]\label{prop:gen}
Let $p>\kappa+2k'$ and $p\nmid v$. Then all $\tilde g_u$ lie in $\cO^2$. There are vectors $\mathcal V(c)\in\cO^2$ on the cells of the
periods outside the pivot, of the period $\kappa-1$ and of the period $\kappa+2k'-1$, such that every label $u$ in these
periods satisfies $\tilde g_u\equiv\epsilon_u\mathcal V(c(u))\pmod{p^2}$ with a unit $\epsilon_u\equiv1\pmod p$, and
$\mathcal V(j,0)\equiv\mathcal V(j,1)\pmod p$. Consequently, for two such labels $u\ne v$,
\begin{gather*}
 v\bigl(\widetilde{\mathrm{Sw}}(u,v)\bigr)\ge-3\ \text{if }u\equiv v\pmod p,\\
 v\bigl(\widetilde{\mathrm{Sw}}(u,v)\bigr)\ge-W_3(c(u),c(v))\ \text{otherwise},
\end{gather*}
and $v(\widetilde{\mathrm{Sw}}(t,t))\ge-3$ on deleted nodes.
\end{proposition}
\begin{proof}
Let $Q=p^3H[\mathcal P_r,\mathcal P_c]$. By Theorem~\ref{thm:nd} and the direct sum above, $Q$ is invertible over $\cO$. We
look for the solution $\Theta$ of $Q\Theta=(\mathbf 1,p^2f)$ modulo $p^2$. Modulo $p^2$ the right side is a function of the
cell, and modulo $p$ of the period. Put $\Theta(t)=\epsilon(t)Y^{h}(j_t)$ with $Y^h=(B^h_X)^{-1}\Gamma_h$ on a class
oriented like $\hat X$, and the transposed matrix on the others. The congruent part of $Q(u,\cdot)\Theta$ is then
$\epsilon B^h(j_u,\cdot)Y^h$, which is $\epsilon\Gamma_h(j_u)$ on the pivot. The non-congruent part lies in $p\cO$, and
modulo $p^2$ it equals
\[
 p\sum_{\omega'\ne\omega_u}\frac1{\omega'-\omega_u}\bigl(\Psi_{o(\omega')}[1,\cdot]-\bar s(j_u)\Psi_{o(\omega')}[0,\cdot]\bigr),
 \qquad \Psi_X=\bar\Gamma^{\mathsf T}\bar B_X^{-1}\bar\Gamma ,
\]
with $\Psi_Y=\Psi_X^{\mathsf T}$. Here we use that the pivot consists of complete periods, so every residue carries
$\Psi_X$ or $\Psi_Y$ and nothing else. Since $\sum_{\omega'\ne\omega}(\omega'-\omega)^{-1}\equiv0\pmod p$, only the
antisymmetric part $\Psi_X-\Psi_X^{\mathsf T}$ survives. It acts on $2$-vectors as a scalar $\delta$ times the rotation, so
the sum is $p\delta E(\omega_u)\bar\Gamma(j_u)$, where $E(\omega)$ is the sum of $(\omega'-\omega)^{-1}$ over the residues
of one orientation. Taking $\epsilon=1-p\delta E$ gives $Q\Theta\equiv(\mathbf 1,p^2f)\pmod{p^2}$ on the pivot, and the
same computation outside the pivot gives $\tilde g_z\equiv\epsilon(\Gamma(z)-B^h(j_z,\cdot)Y^h)\pmod{p^2}$. These are the
vectors $\mathcal V$. They depend on the cell modulo $p^2$, and on the period modulo $p$, since $B^h$ and $\Gamma_h$ do
not depend on $h$ modulo $p$. For $k'=0$ we have $\delta=0$.

By Lemma~\ref{lem:displace}, $\widetilde{\mathrm{Sw}}(u,v)=p^{-2}(\tilde g_u\wedge\tilde g_v)/(v-u)$. The wedge lies in
$p^2\cO$ for two labels in one cell and in $p\cO$ for two labels in one period. The difference $v-u$ has valuation $1$ if
$u\equiv v$ and $0$ otherwise. The diagonal entry on a deleted node is $-p^{-3}Q^{-1}(t,t)$.
\end{proof}

\subsection{The local bounds from elimination}
We combine Lemma~\ref{lem:displace}(a) with Proposition~\ref{prop:gen}. A potential on the labels bounds the Leibniz
terms of the minor of $\widetilde{\mathrm{Sw}}$, and the factor $\det H[\mathcal P_r,\mathcal P_c]$ has valuation exactly
$-3|\mathcal P_r|$.

\begin{theorem}[local bounds of the second plane]\label{thm:rows2}
Let $p\ge5$ with $p^2>2\bar N+2p$ and $p\nmid v$, and let $2m\le r$. For every member the following hold.
\begin{enumerate}[label=(\roman*),leftmargin=*]
\item (E$_1$) If $p>2m$, then $e_p(r,I)\le3r+\sum u+\sum u'$ for all nonnegative potentials $u$ on $X'$ and $u'$ on $Y'$ with
 $W_3\le u\oplus u'$ on the border.
\item (E$_2$) If $p\le r$ and $\bar N\le2p$, then (E$_1$) holds with $3r$ replaced by $3r-(2p-\bar N)$.
\item (R) If $p\le2m$, write $2m=kp+s$ with $0\le s<p$, and let $E=[r,r+s)$ for even $k$ and $E=[r+s,r+p)$ for odd $k$. For
 $t\in E$ let $c_b(t)\in\{0,1,2,3\}$ be twice the period of $t$ minus that of $r$, plus the half of $t$. Then, for each
 $w\in\{(2,2,0,1),(2,2,1,0),(0,1,2,2),(1,0,2,2),(1,1,1,1)\}$,
 \[
 e_p(r,I)\le3(r+m)-3\|m\|_p+\sum_{t\in E}w(c_b(t)).
 \]
\end{enumerate}
\end{theorem}
\begin{proof}
In each case $\kappa+2k'\le\bar N/p+1<p$, so Theorem~\ref{thm:nd} and Proposition~\ref{prop:gen} apply.

(i) Take $\kappa=\lfloor r/p\rfloor+1$ and $k'=0$. Since $2m<p$, no two border nodes are congruent. The nodes of
$[\bar N,\kappa p)$ lie in one period with the border nodes beyond $\kappa p-p$. Give them the potentials $\frac32$ on both
sides, and the deleted nodes the potentials $3$. Proposition~\ref{prop:gen} bounds every Leibniz term of the minor of
$\widetilde{\mathrm{Sw}}$, and Lemma~\ref{lem:displace}(a) with the count of the pivot gives the bound for the Gauss
valuation. The coefficient of $x^i$ carries an extra $3i\ge0$.

(ii) Take $\kappa=2$, $k'=0$ and $X$ replaced by $pX$. Then $\bar B_X=v\bigl(\begin{smallmatrix}0&-1\\-1&2\end{smallmatrix}
\bigr)$, because there is no pole at $k=0$. It is invertible since $p\nmid v$, and its inverse has entry $0$ at $(1,1)$.
So $Q^{-1}(t,t)\in p\cO$ for $t$ in the period $1$, and the virtual nodes $[\bar N,2p)$ need the potentials $1$ only.
This saves $2p-\bar N$.

(iii) Take $\kappa=\lfloor r/p\rfloor+1$ and $k'=\lfloor k/2\rfloor$, and put $h^*=(\kappa+2k')p$. Then
$h^*\le\bar N+p$. The \emph{virtual nodes} $E_v=[\bar N,h^*)$ are empty if $h^*\le\bar N$. The labels of
Lemma~\ref{lem:displace}(a) outside $I$ are the nodes of
\[
 \mathcal L=[r,\kappa p)\cup E_v\cup[h^*,\bar N).
\]
On $[r,\kappa p)$ and on $[h^*,\bar N)$ a node of $\hat X$ is a row label and a node of $\hat Y$ a column label. On $E_v$
the two types are exchanged, since there the pivot rows and columns that are not actual rows and columns become labels.
Put $\mathrm{sh}(t)=t-2k'p$ for $t\ge h^*$ and $\mathrm{sh}(t)=t$ otherwise. This map keeps parity, residue and half,
and it preserves the weights of Proposition~\ref{prop:gen}. The set $E$ has $2\|m\|_p$ nodes and lies in the periods
$\lfloor r/p\rfloor$ and $\lfloor r/p\rfloor+1$, so $c_b$ maps it into four fixed cells. Each $w$ in the list dominates
$W_3$ on these four cells.

Choose a period $j_0$ as follows. If $E$ meets one period, $j_0$ is that period. If $E$ is empty, which happens when
$\|m\|_p=0$, put $j_0=\kappa-1$. If $E$ meets two periods $j$ and $j+1$, put $j_0=j$ when $w\ge1$ on the cells of $E$
in $j+1$, and $j_0=j+1$ otherwise. This is always possible, since for each $w$ in the list a zero in one period forces
$w=2$ on the other. Define the potential of a label $t$, the same on both sides, by
\[
 \phi(t)=\begin{cases}w(c_b(\mathrm{sh}\,t))&\text{if }\mathrm{sh}\,t\in E,\\ 1&\text{if }\mathrm{sh}\,t\notin E
 \text{ lies in the period }j_0,\\ 2&\text{otherwise.}\end{cases}
\]
Then $W\le\phi(t)+\phi(t')$ for every row label $t$ and column label $t'$, where $W$ is the bound of
Proposition~\ref{prop:gen}. Two labels of $E$ are covered by the dominance of $w$. A label outside $E$ in the period
$j_0$ against a label of $E$ has $W\le1$ inside $j_0$, and $W=2\le1+w$ outside $j_0$, by the choice of $j_0$. A label
outside $E$ and outside $j_0$ has potential $2$, and every pair with it that is not congruent has $W\le2$. Every
congruent pair $(t,t+p)$ off $E$ has exactly one node in $j_0$. So does every pair of $U_2=[r+s,\kappa p)$ and $E_v$
when $k$ is even and $h^*>\bar N$, since $t\mapsto t+kp$ maps $U_2$ onto $E_v$. Such pairs have $W\le3=1+2$. All
other pairs off $E$ have $W\le2$. The sum of the potentials, with the potentials $3$ on deleted nodes and the count of
the pivot as in (i), is $3(r+m)-3\|m\|_p+\sum_{t\in E}w(c_b(t))$.
\end{proof}

Each bound of Theorem~\ref{thm:rows2} is at least $2(r+m)$, which justifies the remark before
Proposition~\ref{prop:rows2} on the primes $p\mid v$. Indeed (E$_1$) is at least $3r\ge2r+2m$, (E$_2$) is at least
$4r+2m-2p\ge2r+2m$ since $p\le r$, and (R) is at least $3r+3m-3\|m\|_p\ge3r\ge2r+2m$.

For $k=1$ part (iii) needs only the pivot of (i), with the weight $3$ of congruent pairs kept. For $k\ge2$ it needs the
complete double periods. The pivot contracts the window of length $2m$ to one of length $p+s$ or $s$, and the cost of
each eliminated pair of periods is exactly $3$ per node.

\subsection{The denominator exponent}
Homogenise with $\theta=p/r$ as before. The profile $\mathcal P(\theta)$ is the minimum of the homogenised bounds of
Proposition~\ref{prop:rows2} and Theorem~\ref{thm:rows2}. For (E$_1$) and (E$_2$) the potentials are constant on the at most
three cells cut from the border at the multiples of $\theta/2$. On the band $\frac{2\alpha}{k+1}<\theta\le\frac{2\alpha}k$, the
bound (R) is
\[
 3+\tfrac32k\theta+\tau(E)\ (k\text{ even}),\qquad 3+6\alpha-\tfrac32(k+1)\theta+\tau(E)\ (k\text{ odd}),
\]
where $\tau(E)$ is the least over the five vectors $w$ of the sums $\sum_cw(c)\,|c\cap E|$. Here $c$ runs over the four
cells, scaled by $r$, $|c\cap E|$ is the length of the intersection, and $E=[1,1+2\alpha-k\theta]$, respectively
$E=[1+2\alpha-k\theta,1+\theta]$. Each node of $E$ carries one potential, so there is no factor $\frac12$. Between the points $\theta=2/j$ and $\theta=2(1+2\alpha)/j$ the cell type of $E$ is fixed. So
$\mathcal P$ is piecewise affine, with finitely many pieces on every interval bounded away from $0$.

On $\theta<2\alpha$ the bounds (P1) and (R) are closed under translation in the sense of Theorem~\ref{thm:periodic},
with $a=3$, $b=41$ and $K_0=3+3\alpha$. The translation of order $\nu$ takes the band $k$ to the band $k+6\nu$. So the
table for $\Lambda$ is checked on one period of $1/\theta$, the bands $k=1,\dots,6$, and on the head $\theta\ge2\alpha$.
Its translates cover the bands down to $t_1=3\cdot10^6/(1025\cdot10^6+3)$, just below $\alpha/25$. On $(0,t_1]$ we use
(P1) without its last term.

\begin{theorem}[denominators on the second plane]\label{thm:content2}
Let $\alpha=\frac3{41}$ and $v\ne0$, with $u$ any integer. For every $\epsilon>0$ and every even $r\ge R(\epsilon)$, every member with $|I|\le1$
satisfies $\log\den P_{r,I}\le(\Lambda_2+\epsilon)r^2$, where
\[
 \Lambda_2=(3+3\alpha)t_1+\int_{t_1}^\infty\mathcal P(\theta)\,d\theta\le\frac{3476837531}{10^9}.
\]
The threshold does not depend on $(u,v)$.
\end{theorem}
\begin{proof}
For each prime with $p^2>2\bar N+2p$ the local bounds give $e_p(r,I)\le r\,\mathcal P(p/r)+O(1)$. For $p\mid v$ the bounds
of Theorem~\ref{thm:rows2} are replaced by the trivial bound $2(r+m)$, which is at most each of them. For the bounds
(E$_1$), (E$_2$) and (R) the cells of the integer window and of its homogenisation differ by a bounded number of cells of bounded length,
which take the potential $2$. The hypotheses $p>\kappa+2k'$ of Proposition~\ref{prop:gen} hold for $r$ large, since
$p>t_1r$ on the bands used. For the other primes every entry has valuation at least $-3\log(2\bar N)/\log p$, which
costs $O(r^{3/2})$ in total. The prime number theorem then gives the integral, as in Theorem~\ref{thm:content}, and
Theorem~\ref{thm:periodic} gives the pieces of the translated bands from those of the period. All pieces are affine, so
the integral is a rational number. It is at most the stated one.
\end{proof}

For $v=0$ we take the frame $F^{(2)}_{0,1}=\psi(z+1)-\psi(z+\frac12)$. Its matrix $M$ is $4$ times that of $F_{1/2}$
in \eqref{eq:frames}, so the coefficient of $x^i$ in a member is $4^{n-i}$ times the corresponding coefficient for
$F_{1/2}$, where $n$ is the size of the member. Its denominators are therefore at most those of $F_{1/2}$. Theorem~\ref{thm:content}, with an admissible
table for $q=2$ and $c=\frac6{41}$, bounds them with the exponent $\Lambda_\pi=2.60981\,51766\ldots$, which is less
than $\Lambda_2$.

\subsection{Decay and proof of Theorem~\ref{thm:plane2}}
By Theorem~\ref{thm:nvframe} with $d\le3$, one deleted node suffices for nonvanishing. The decay is
Theorem~\ref{thm:decayframe} with the reference frame of shifts $1$, $\frac34$, $\frac12$.

\begin{proposition}\label{prop:cert2}
Take the reference frame of shifts $1$, $\frac34$, $\frac12$ and $c=\frac6{41}$. There is a trial measure with $49$
polynomial pieces and a cover by $508$ cells satisfying the checks of \S\ref{ssec:slope-cert} with $\varepsilon_0=10^{-6}$
and $|\mathcal T|=3$. For it Theorem~\ref{thm:decay} gives $D\ge\DPlB$ and $\rho\le\gPlB$.
\end{proposition}

\begin{proof}[Proof of Theorem~\ref{thm:plane2}]
Let $a$ and $b$ be rationals, not both zero. Choose a rational $\lambda\ne0$ for which $v=\lambda a$ and $u=6\lambda b$
are integers. By \S\ref{ssec:frames} the number $\xi$ of $F^{(2)}_{v,u}$ is $-2\lambda(a\zeta(3)+b\pi^2)$. For $a=0$ we
take $\lambda$ with $u=1$, so the frame is $\psi(z+1)-\psi(z+\frac12)$. Its trigamma term has weight zero, and it has
two simple poles per period. For $b=0$ we take $\lambda$ with $v=1$, so the frame is $\psi'(z+1)$. Its shift
$\frac12$ has weight zero, and it has one double pole per period. Both are handled exactly as the first plane handles $a=0$, through a weight zero in a
frame on the same shifts. The geometry is the same for all $(a,b)$. For $v\ne0$, Theorem~\ref{thm:content2} gives the
denominators with $\Lambda_2$. For $v=0$, the exponent $\Lambda_\pi$ above gives them, and $\Lambda_\pi<\Lambda_2$.
Proposition~\ref{prop:cert2} with Theorem~\ref{thm:decayframe} gives the decay for every frame, also when one weight is
zero, and Theorem~\ref{thm:nvframe} gives nonvanishing with one deleted node, since $d\le3$. Theorem~\ref{thm:deduction}
gives
\[
 \mu(\xi)\le\frac\rho{D-\Lambda_2}\le\ExpPlB<\MuPlaneB .\qedhere
\]
\end{proof}

Theorem~\ref{thm:ends} gives better bounds at the two ends, by separate proofs.

\subsection{The value $\zeta(3)$}\label{ssec:zeta3}
For the separate bound of Theorem~\ref{thm:ends} we take $u=0$, $v=1$ and $c=\frac3{10}$. The frame is $\psi'(z+1)$,
with $d=2$ and $\xi=-2\zeta(3)$. Everything above holds with the half map set to $0$.
The cells are then the periods, $W_3$ takes the values $0$ and $2$, and $\varepsilon=0$ in
Proposition~\ref{prop:rows2}. The bounds (P4) and (P5) become $6(r+m-p)$ and $3(\bar N-p)_+$, and (P6), (P7) vanish. In (R)
the four cells become two, and the potentials are $(0,0,2,2)$ and $(2,2,0,0)$. With $c=\frac3{10}$, and with (P1)
without its last term on $(0,\frac3{500}]$, the profile has integral
\[
 \Lambda_{\zeta(3)}\le\frac{2912839}{10^6}.
\]
The reference frame is the function $F_{1/2}$, and its certificate is that of the row for $\pi^2$ in
Table~\ref{tab:data}. Theorem~\ref{thm:nvframe} with $d=2$ gives nonvanishing with one deleted node.

\begin{proof}[Proof of Theorem~\ref{thm:ends} for $\zeta(3)$]
The three estimates above and Theorem~\ref{thm:deduction} give
$\mu(\zeta(3))\le\rho/(D-\Lambda_{\zeta(3)})\le\ExpZ<\MuZetaThree$. The bound for $\pi^2$ is proved in
Section~\ref{sec:proof}, together with Theorem~\ref{thm:offplane}.
\end{proof}

%% =====================================================================
\section{The deduction and the remaining proofs}\label{sec:proof}

The three estimates combine as follows.

\begin{theorem}[deduction]\label{thm:deduction}
Let $F$ be a frame with a geometry. Let $K$ be a depth such that, for every large $r$ and every rational $x_0$, some
member $P_{r,I}$ with $|I|\le K$ satisfies $P_{r,I}(x_0)\ne0$, as in Theorems~\ref{thm:selection} and~\ref{thm:nvframe}.
Suppose that the members satisfy $\log\den P_{r,I}\le(\Lambda+o(1))r^2$ and the decay estimate of
Theorem~\ref{thm:decay} with constants $D$ and $\rho$, and that $\Lambda<D$. Then $\xi$ is irrational and
$\mu(\xi)\le\rho/(D-\Lambda)$.
\end{theorem}

\begin{proof}
Let $\eta=D-\Lambda$, and let $\mu_0>\mu_1>\rho/\eta$. Suppose that $|\xi-a/b|<b^{-\mu_0}$ with $b$ large. Let $r$ be
the least admissible scale with $r\ge\mu_1\log b/\rho$. With $x_0=a/b$ there is a member with $P_{r,I}(a/b)\ne0$. By
Lemma~\ref{lem:integrality}, $|\den(P_{r,I})\,P_{r,I}(a/b)|\ge b^{-r}$. Write $a/b=\xi+h$.
Then $|h|<b^{-\mu_0}\le e^{-\rho r}$, because $\rho r\le\mu_1\log b+O(1)<\mu_0\log b$ for large $b$. By the two
estimates,
\[
 |\den(P_{r,I})\,P_{r,I}(a/b)|\le e^{(\Lambda+o(1))r^2}e^{-(D-o(1))r^2}=e^{-(\eta-o(1))r^2}.
\]
Comparing the two bounds gives $(\eta-o(1))r^2\le r\log b$. But $\eta r\ge(\mu_1\eta/\rho)\log b$ with
$\mu_1\eta/\rho>1$. This is a contradiction for large $b$. If $\xi=a/b$ were rational, then for every large $r$ some member would satisfy
$b^{-r}\le|\den(P_{r,I})\,P_{r,I}(\xi)|\le e^{-(\eta-o(1))r^2}$, which is impossible. So $\xi$ is irrational.
\end{proof}

\begin{proof}[Proof of Theorem~\ref{thm:offplane}, Theorem~\ref{thm:points1}(ii) and Theorem~\ref{thm:ends} for $\pi^2$]
For each of these rows of Table~\ref{tab:main} the function is $F_s$ of \eqref{eq:frames}, and the geometry parameter and
the constants are those of Table~\ref{tab:data}. The denominator exponent is given by Theorem~\ref{thm:content} with the
table for $\Lambda$ of the formal proof. The decay rates are given by Proposition~\ref{prop:cert}, and nonvanishing by
Theorem~\ref{thm:selection}. Theorem~\ref{thm:deduction}, applied at the even scales, gives $\mu(\xi)\le\rho/(D-\Lambda)$,
and in every row $\rho/(D-\Lambda)$ is below the stated bound. By \eqref{eq:frames}, $\xi$ is a nonzero rational multiple
of $\Psi(s,1)$. Since the inequality is strict, the definition of $\mu$ gives $|\xi-a/b|>b^{-\mu_0}$ for all integers $a$
and all sufficiently large $b$.
\end{proof}

Theorem~\ref{thm:points1}(i) is proved in Section~\ref{sec:chi3-quarter}, parts (iii)--(v) and Theorem~\ref{thm:plane1}
in Section~\ref{sec:plane1}, and Theorem~\ref{thm:plane2} and the bound for $\zeta(3)$ in Section~\ref{sec:plane2}. In each
case the bound is attained in the same strong form.

\begin{table}[ht]
\centering
{\setlength{\tabcolsep}{1.5pt}
\begin{tabular}{@{}lcccccc@{}}
\toprule
case & $c$ & $\Lambda$ & $D$ & $\rho$ & $\rho/\eta$ & $\mu$\\
\midrule
plane 1 & $\frac16$ & $\LamPl$ & $\DPl$ & $\gPl$ & $\ExpPl$ & $\MuPlaneA$\\
$L\pm\frac{2\pi^2}{27}$ & $\frac16$ & $\LamDA$ & $\DPl$ & $\gPl$ & $\ExpDA$ & $\MuDipA$\\
$L\pm\frac{\pi^2}{27}$ & $\frac16$ & $\LamDB$ & $\DPl$ & $\gPl$ & $\ExpDB$ & $\MuDipB$\\
$L\pm\frac{\pi^2}{18}$ & $\frac16$ & $\LamDC$ & $\DPl$ & $\gPl$ & $\ExpDC$ & $\MuDipC$\\
$(\frac13,\frac23)$ & $\frac12$ & $\LamChi$ & $\DChi$ & $\gChi$ & $\ExpChi$ & $\MuChi$\\
$(\frac k3,1)$ & $\frac3{10}$ & $\LamThree$ & $\DThree$ & $\gThree$ & $\ExpThree$ & $\MuThirdsU$\\
\addlinespace
plane 2 & $\frac6{41}$ & $\LamPlB$ & $\DPlB$ & $\gPlB$ & $\ExpPlB$ & $\MuPlaneB$\\
$\zeta(3)$ & $\frac3{10}$ & $\LamZ$ & $\DHalf$ & $\gHalf$ & $\ExpZ$ & $\MuZetaThree$\\
$(\frac12,1)$ & $\frac3{10}$ & $\LamHalf$ & $\DHalf$ & $\gHalf$ & $\ExpHalf$ & $\MuHalvesU$\\
\addlinespace
$(\frac k4,1)$ & $\frac15$ & $\LamFour$ & $\DFour$ & $\gFour$ & $\ExpFour$ & $\MuLamU$\\
$(\frac k5,1)$ & $\frac29$ & $\LamFive$ & $\DFive$ & $\gFive$ & $\ExpFive$ & $\MuFifthsU$\\
$(\frac k7,1)$ & $\frac16$ & $\LamSeven$ & $\DSeven$ & $\gSeven$ & $\ExpSeven$ & $\MuSeventhsU$\\
\bottomrule
\end{tabular}}
\caption{The data of the proofs. A pair $(s,t)$ denotes the function of \eqref{eq:frames} for $\Psi(s,t)$, with $k$
prime to the denominator. For each row the table gives the geometry parameter $c$, the denominator exponent $\Lambda$, the
decay rate $D$, the neighbourhood exponent $\rho$ and the exponent. The constants are rational numbers, rounded outward to
the digits shown: $\Lambda$ and $\rho$ up, $D$ down. The column $\rho/\eta$, with $\eta=D-\Lambda$, is the exact quotient,
rounded up. The three rows below plane 1 share its decay certificate. The row for $\zeta(3)$ shares the
certificate of the row $(\frac12,1)$.}\label{tab:data}
\end{table}

\begin{remark}[relation to linear independence criteria]
The deduction is a one-dimensional version of the familiar scheme behind the criteria of Nesterenko
\cite{Nesterenko85} and their use by Ball and Rivoal \cite{BallRivoal01}. The nonzero integer
$b^{|C\setminus I|}\den(P_{r,I})\,P_{r,I}(a/b)$ plays the role of a small linear form with integer coefficients. There are
two differences. The lower bound for the form comes from an arithmetic nonvanishing statement rather than from an analytic
asymptotic. The size $r$ of the determinant is chosen in terms of $b$.
\end{remark}

\subsection{Notes on the individual cases}\label{sec:families}
For the rows $(\frac kq,1)$ the function is $F_s$ of \eqref{eq:frames} with $s=k/q$, $\omega=q^{-2}$ and
$\xi_s=-q^{-2}\Psi(\frac kq,1)$. The table for $\Lambda$, the trial measure and the cover are data of the formal proof. All
shifts with the same denominator share one table for $\Lambda$, indexed by the root classes, and one certificate. Indeed the
bounds depend on the shift only through the root class, and the decay depends on it only through the separation of the
poles from the contour. For $q=4$ the root classes $j\equiv kp^{-1}\pmod 4$ of $k=1$ and $k=3$ exchange the two classes of
primes modulo $4$. For $q=3$ the table uses the primes in the classes modulo $6$, and the root classes $j=1,2$ correspond to
first poles near $p/3$ and $2p/3$. For $q=7$ there are six root classes. Here $D-\Lambda$ is only about $4\cdot10^{-3}$,
which explains the large exponent. A small further gain in either constant would lower it substantially
(Section~\ref{sec:limits}). The certificate for $q=7$ therefore takes $\varepsilon_0=10^{-6}$ in the function $\Upsilon$,
where the cases $q=2,3,4,5$ take $\varepsilon_0=10^{-3}$. Its trial measure has $22$ polynomial pieces, and its cells on
$\Gamma$ at the breakpoints of $\sigma$ use the truncated logarithm of \S\ref{ssec:slope-cert}. For $q=2$ there is one root
class, $j=1$, with first pole $k_p=\frac{p+1}2$.

The certificates of the two planes also take $\varepsilon_0=10^{-6}$. On the second plane $D-\Lambda$ is about
$0.021$, which explains the size of its exponent. On the first plane the points of Theorem~\ref{thm:points1}(iii)--(v)
gain only in $\Lambda$, since they use the certificate of the plane.
%% =====================================================================
\section{Limits of the method}\label{sec:limits}

The exponents are determined by three constants through $\rho/(D-\Lambda)$. Since $\partial\mu/\partial D=-\rho/\eta^2$,
the exponent is most sensitive to the decay where $\eta$ is small. For the shifts $\frac14$ and $\frac34$,
$\rho/\eta^2\approx136$. For the shifts $\frac k7$ it is about $4.4\cdot10^5$.

On the arithmetic side, the profiles $\mathcal P_j$ are minima of the bounds of Theorem~\ref{thm:local}. Every further
local bound lowers $\Lambda$ by the integral of the amount by which it lies below $\mathcal P_j$. The denominator exponent of
the method is the integral of the true exponent $e_p$, which the local bounds approximate from above. At third order this
approximation is close to exhaustion for the shifts $\frac k5$ and $\frac k7$. We computed the Smith normal forms exactly
at scales $r\le2400$ and primes $p\le2m$. There the minimum of the local bounds, with (U$_n$) included, equals the true
exponent at all but a few isolated primes. For the shifts $\frac k7$ the invariants of valuation at least $4$ contribute
about $5\cdot10^{-6}$ to $\Lambda$. The bounds (U$_n$) gain about $7\cdot10^{-4}$. So at third order the local bounds are
essentially exhausted. These computations are not part of the proof.

On the analytic side, Remark~\ref{rem:certlimit} explains why a certificate of the present form cannot do better than the
constrained equilibrium measure on the chosen contour. A larger decay rate needs a different contour or a different
treatment of the energy problem. The geometry parameter $c$ trades the two sides against each other. A longer border
increases $D$, through the larger node measure. It also increases $\Lambda$, through the border rows and columns of the
assignment bound.

%% =====================================================================
\section{Formalization}\label{sec:lean}
\raggedbottom

The theorems of Section~\ref{sec:intro} are formalized in Lean~4 \cite{Lean4} with Mathlib \cite{Mathlib}, together with
every theorem, proposition and lemma their proofs use. The remarks and Section~\ref{sec:limits} are not part of the
formalization. The development is available at
\begin{center}
\FormalizationLink
\end{center}
It combines analytic and arithmetic proofs with generated rational data. Kernel checks verify the tables for $\Lambda$ and
the decay certificates. Numerical computations alone are not used as proof.

The statements are in the file \lean{ChallengeHybrid.lean}. Their proofs are in \lean{SolutionHybrid.lean}, in the
namespace \lean{TrigammaDifferences}. Each number $\xi$ is written as a series, or through \lean{Real.pi^2}. The
series is the classical one where there is one. Otherwise it is $\Psi(s,t)$, and $\zeta(3)$ is the series
$\sum_{n\ge0}(n+1)^{-3}$. Each point theorem states that $b^{-\mu_0}<|\xi-a/b|$ for all $a\in\Z$ and all sufficiently
large $b$. The plane theorems take the rational coefficients $a$ and $b$, not both zero, and write the fraction as $p/q$.
They state that $q^{-\mu_0}<|\xi-p/q|$ for all $p\in\Z$ and all sufficiently large $q$, where $\xi$ is $aL+b\pi^2$,
respectively $a\zeta(3)+b\pi^2$. The statements have no separate irrationality clause, since the bound for all large
denominators implies irrationality. The theorems are the following.
\begin{itemize}
\item \lean{plane1_measure}: Theorem~\ref{thm:plane1}, with $\mu_0=\MuPlaneA$.
\item \lean{chi3_measure}: Theorem~\ref{thm:points1}(i), with $\mu_0=\MuChi$.
\item \lean{chi3_plus_pi2_measure} and \lean{chi3_minus_pi2_measure}: Theorem~\ref{thm:points1}(ii), with
 $\mu_0=\MuThirdsU$.
\item \lean{dip_2_27_plus_measure}, \lean{dip_2_27_minus_measure}, \lean{dip_1_27_plus_measure},
 \lean{dip_1_27_minus_measure}, \lean{dip_1_18_plus_measure} and \lean{dip_1_18_minus_measure}:
 Theorem~\ref{thm:points1}(iii)--(v), with $\mu_0=\MuDipA$, $\MuDipB$ and $\MuDipC$.
\item \lean{plane2_measure}: Theorem~\ref{thm:plane2}, with $\mu_0=\MuPlaneB$.
\item \lean{zeta3_measure} and \lean{pi2_measure}: Theorem~\ref{thm:ends}, with $\mu_0=\MuZetaThree$ and
 $\MuHalvesU$.
\item \lean{lambda_plus_measure} and \lean{lambda_minus_measure}: Theorem~\ref{thm:offplane}(i), with
 $\mu_0=\MuLamU$.
\item \lean{fifths1_measure}, \dots, \lean{fifths4_measure}: Theorem~\ref{thm:offplane}(ii), with
 $\mu_0=\MuFifthsU$.
\item \lean{sevenths1_measure}, \dots, \lean{sevenths6_measure}: Theorem~\ref{thm:offplane}(iii), with
 $\mu_0=\MuSeventhsU$.
\end{itemize}
The shared results are in the namespace \lean{CatalanCore.Uniform}.
\begin{itemize}
\item \lean{Content.BankShared.Check.hasScaledContent_of_check} is Theorem~\ref{thm:content}. The local bounds of
 Theorem~\ref{thm:local} are proved in \lean{Check/Catalogue.lean}, \lean{Check/LP2W.lean},
 \lean{Check/Prefix.lean} and \lean{Check/UF.lean}, with the closed rank bound and the branches of
 Theorem~\ref{thm:wrap}. Its form for functions with several first poles is
 \lean{Content.BankShared.KTerm.hasScaledContent_of_kcheck}.
\item \lean{Content.BankShared.TwoTerm.LP2HypW.gle} is Theorem~\ref{thm:lp2}.
\item \lean{Content.BankShared.TwoTerm.Pre.prefix_gle} is Theorem~\ref{thm:window}.
\item \lean{Content.UFSimple.WrapCheck.catValid_wrap} and \lean{Content.UFSimple.WrapPrefixCheck.catValid_prefix_full}
 are Theorem~\ref{thm:wrap} for $n=0$ and for $n\ge1$. Lemma~\ref{lem:harmonic} is
 \lean{Content.UFSimple.WrapPrefixClosed.HSplit.det_ne} with \lean{HSplit.dir}.
\item \lean{Support.support_nv0_even} is Theorem~\ref{thm:nvframe}. \lean{BlockGeom.hasNV0} is
 Theorem~\ref{thm:selection}.
\item \lean{Decay.AltDecay.altGeom_hasScaledDecay} is Theorem~\ref{thm:decay}. It uses the node energy check
 \lean{Decay/Nodes/EnergyCheck.lean}, the contour check \lean{Decay/ContourCheck.lean} and the cell checks
 \lean{Decay/PolyWire3}. For the shifts $\frac k7$ it is applied in the form
 \lean{Decay.AltDecay.altGeom_hasScaledDecay_Tr}, with the same conclusion. There the cells on the contour at breakpoints use the
 truncated logarithm, proved in \lean{Decay/Breakpoint/TruncLog.lean} and \lean{Decay/Breakpoint/CellT.lean}.
\item \lean{measure_of_family_nv0_even} is Theorem~\ref{thm:deduction} at the even scales.
\item \lean{Content.BankShared.KTerm.hasScaledContent_of_periodic} is Theorem~\ref{thm:periodic} in the form used
 for Theorems~\ref{thm:planecontent} and~\ref{thm:content2}.
\item \lean{Bank.Plane1.Dips.dip_point} is Theorem~\ref{thm:dips} together with the deduction of
 Theorem~\ref{thm:points1}(iii)--(v).
\item \lean{Bank.Plane2.plane2_point_ab_of_content} is the proof of Theorem~\ref{thm:plane2}, for every rational
 $(a,b)\ne(0,0)$, from the denominator bounds of the three cases $v\ne0$ with $u\ne0$, $u=0$ and $v=0$.
\end{itemize}
The solution has no \lean{sorry} and no \lean{native_decide}. The command \texttt{\#print axioms} reports only the three
standard axioms \lean{propext}, \lean{Classical.choice} and \lean{Quot.sound}. The prime number theorem in arithmetic
progressions enters the asymptotics of the denominators. Its formal source is the PNT+ project \cite{PNTPlus}, through OpenAI's
workspace and compatibility additions. We use its Wiener--Ikehara theorem.

The proofs were checked with Lean \texttt{v4.35.0-rc2} and the matching Mathlib release. The vendored files from OpenAI's
repository and from the PNT+ project were converted to Lean's module system and ported to the current toolchain. The repository's README and provenance documents describe the changes. The certificate data were generated with
the interval arithmetic of Arb \cite{Arb}, through \texttt{python-flint}, and Lean re-checks them.

\section*{Disclosure of AI use}
The mathematics, the Lean formalization, the certificate generators and the text of this paper were produced by AI
systems at the request of the author. The author posed the problems, directed the work and is responsible for the
paper. The systems were Claude Opus~5.5, Claude Sonnet~5.5 and Claude Haiku~5.5 (Anthropic), running in Claude Code,
and GPT-6.1 Sol (OpenAI), running in Codex. The underlying irrationality proof for Catalan's constant and its Lean
formalization are due to OpenAI \cite{OpenAI2026}.

\subsection*{Acknowledgements}
We thank the authors of Mathlib and of the PNT+ project, and OpenAI for making its formalization available.

\providecommand{\bysame}{\leavevmode\hbox to3em{\hrulefill}\thinspace}
\providecommand{\MR}{\relax\ifhmode\unskip\space\fi MR }
% \MRhref is called by the amsart/book/proc definition of \MR.
\providecommand{\MRhref}[2]{%
  \href{http://www.ams.org/mathscinet-getitem?mr=#1}{#2}
}
\providecommand{\href}[2]{#2}
\begin{thebibliography}{10}

\bibitem{Apery79}
Roger Ap\'ery, \emph{Irrationalit\'e de {$\zeta(2)$} et {$\zeta(3)$}},
  Ast\'erisque \textbf{61} (1979), 11--13.

\bibitem{BallRivoal01}
Keith Ball and Tanguy Rivoal, \emph{Irrationalit\'e d'une infinit\'e de valeurs
  de la fonction z\^eta aux entiers impairs}, Invent. Math. \textbf{146}
  (2001), no.~1, 193--207.

\bibitem{Beukers79}
Frits Beukers, \emph{A note on the irrationality of {$\zeta(2)$} and
  {$\zeta(3)$}}, Bull. London Math. Soc. \textbf{11} (1979), no.~3, 268--272.

\bibitem{Bojnal2026}
Hrishik~Sai Bojnal, \emph{An irrationality exponent below 2402 for {C}atalan's
  constant}, GitHub repository, October 15, 2026,
  \url{https://github.com/proof-factory/catalan-exponent}; design search and
  certification by Claude Fable~5.1 (Anthropic), 2026.

\bibitem{Brown26}
Francis Brown, \emph{Mellin transforms, transfinite diameter and rational
  approximations of integrals}, 2026, Preprint, 22 April 2026,
  \href{https://arxiv.org/abs/2604.20741}{arXiv:2604.20741}.

\bibitem{CDT2024}
Frank Calegari, Vesselin Dimitrov, and Yunqing Tang, \emph{The linear
  independence of {$1$}, {$\zeta(2)$}, and {$L(2,\chi_{-3})$}}, 2024,
  \href{https://arxiv.org/abs/2408.15403}{arXiv:2408.15403}.

\bibitem{CDT2025}
\bysame, \emph{Arithmetic holonomy bounds and effective {D}iophantine
  approximation}, 2025, Theorem 5.1 gives $\mu(L(2,\chi_{-3}))\le 24781$, and
  the remark following it covers every $\mathbb{Q}$-linear combination of
  $L(2,\chi_{-3})$ and $\pi^2$,
  \href{https://arxiv.org/abs/2510.04156}{arXiv:2510.04156}.

\bibitem{Carreira2026b}
Marcos Costa Santos~Carreira, \emph{Determinant constructions for the
  {D}irichlet beta value at four}, Preprint, 9 October 2026,
  \href{https://doi.org/10.5281/zenodo.23326647}{doi:10.5281/zenodo.23326647},
  2026.

\bibitem{Carreira2026}
\bysame, \emph{General designs and an irrationality exponent for {C}atalan's
  constant}, Preprint, 9 October 2026, revised 21 October 2026,
  \href{https://doi.org/10.5281/zenodo.23532239}{doi:10.5281/zenodo.23532239};
  sources at \url{https://github.com/MarcosCarreira/catalan-designs}, 2026.

\bibitem{DauguetZudilin14}
Simon Dauguet and Wadim Zudilin, \emph{On simultaneous diophantine
  approximations to {$\zeta(2)$} and {$\zeta(3)$}}, J. Number Theory
  \textbf{145} (2014), 362--387,
  \href{https://arxiv.org/abs/1401.5322}{arXiv:1401.5322}.

\bibitem{Lean4}
Leonardo de~Moura and Sebastian Ullrich, \emph{The {L}ean 4 theorem prover and
  programming language}, Automated Deduction -- CADE 28, Lecture Notes in
  Computer Science, vol. 12699, Springer, 2021, pp.~625--635.

\bibitem{Eskandari2026}
Payman Eskandari, \emph{Improved rational approximations to {C}atalan's
  constant}, 2026, Preprint, 22 September 2026,
  \href{https://arxiv.org/abs/2609.26354}{arXiv:2609.26354}.

\bibitem{EMN2025}
Payman Eskandari, V.~Kumar Murty, and Yusuke Nemoto, \emph{Mixed motives and
  linear forms in the {C}atalan constant}, J. Lond. Math. Soc. (2) \textbf{114}
  (2026), no.~4, Paper No. e70708,
  \href{https://arxiv.org/abs/2510.20648}{arXiv:2510.20648}.

\bibitem{Fauzan2026}
Aabir Fauzan, \emph{{$\zeta(5)$} is irrational}, Preprint v1, September 17,
  2026, Zenodo,
  \href{https://doi.org/10.5281/zenodo.22826419}{doi:10.5281/zenodo.22826419},
  2026.

\bibitem{Forrester2018Andreief}
Peter~J. Forrester, \emph{Meet {A}ndr\'eief, {B}ordeaux 1886, and {A}ndreev,
  {K}harkov 1882--83}, 2018,
  \href{https://arxiv.org/abs/1806.10411}{arXiv:1806.10411}.

\bibitem{Giraudin26}
David~Niedbala Giraudin, \emph{A new upper bound for the irrationality exponent
  of {$\zeta(2)$}}, 2026, Preprint, 2 October 2026,
  \href{https://arxiv.org/abs/2610.02912}{arXiv:2610.02912}.

\bibitem{Giraudin26z3}
\bysame, \emph{A new upper bound for the irrationality exponent of
  {$\zeta(3)$}}, 2026, Preprint, 22 September 2026,
  \href{https://arxiv.org/abs/2609.26980}{arXiv:2609.26980}.

\bibitem{Jha2026}
R.~K. Jha, \emph{A finite irrationality measure for {C}atalan's constant},
  Preprint v1, October 8, 2026,
  \href{https://doi.org/10.5281/zenodo.23237269}{doi:10.5281/zenodo.23237269};
  Lean sources at \url{https://github.com/rohitkrjha/catalan-measure}, 2026.

\bibitem{Arb}
Fredrik Johansson, \emph{Arb: efficient arbitrary-precision midpoint-radius
  interval arithmetic}, IEEE Trans. Comput. \textbf{66} (2017), no.~8,
  1281--1292, \href{https://arxiv.org/abs/1611.02831}{arXiv:1611.02831}.

\bibitem{KLS}
Roelof Koekoek, Peter~A. Lesky, and Ren\'e~F. Swarttouw, \emph{Hypergeometric
  orthogonal polynomials and their {$q$}-analogues}, Springer Monographs in
  Mathematics, Springer, Berlin, 2010.

\bibitem{KrattRivoal08}
Christian Krattenthaler and Tanguy Rivoal, \emph{On a linear form for
  {C}atalan's constant}, South East Asian J. Math. Math. Sci. \textbf{6}
  (2008), no.~2, 3--15,
  \href{https://arxiv.org/abs/0810.1927}{arXiv:0810.1927}.

\bibitem{KrattZudilin19}
Christian Krattenthaler and Wadim Zudilin, \emph{Hypergeometry inspired by
  irrationality questions}, Kyushu J. Math. \textbf{73} (2019), no.~1,
  189--203, \href{https://arxiv.org/abs/1802.08856}{arXiv:1802.08856}.

\bibitem{Lai2026}
Li~Lai, \emph{The linear independence of {$1$}, {$\zeta(3)$}, and
  {$\zeta(4)$}}, 2026, Preprint, 5 October 2026 (v2: 15 October 2026),
  \href{https://arxiv.org/abs/2610.06838}{arXiv:2610.06838}.

\bibitem{LBKWSK26}
Elyasheev Leibtag, Hila Barkan, Rotem Kalisch, Shachar Weinbaum, Michael
  Shalyt, and Ido Kaminer, \emph{New directions in irrationality}, 2026,
  Preprint, 20 October 2026,
  \href{https://arxiv.org/abs/2610.27799}{arXiv:2610.27799}; supplement and
  code at
  \url{https://github.com/RamanujanMachine/new-directions-in-irrationality};
  the supplement on {$L(2,\chi_{-3})$} and {$\zeta(2)$} is the folder
  \texttt{additional/l2chi3\_zeta2}.

\bibitem{Litsey2026}
Sean Litsey, \emph{An irrationality measure for {C}atalan's constant}, Hexagon
  preprint hexagon:2610.00187, \url{https://hexagonmath.org/2610.00187}, 2026.

\bibitem{LiuJiangZhang26}
Jingwen Liu, Kai Jiang, and Pingwen Zhang, \emph{Irrationality exponents of
  logarithms of positive rational numbers}, 2026, Preprint, 7 October 2026,
  revised 5 November 2026 (v3),
  \href{https://arxiv.org/abs/2610.10192}{arXiv:2610.10192}.

\bibitem{LZZ25}
Junqi Liu, Jujian Zhang, and Lihong Zhi, \emph{A formal proof of the
  irrationality of {$\zeta(3)$} in {L}ean 4}, 2025,
  \href{https://arxiv.org/abs/2503.07625}{arXiv:2503.07625}.

\bibitem{Loewner34}
Karl L\"owner, \emph{\"{U}ber monotone {M}atrixfunktionen}, Math. Z.
  \textbf{38} (1934), 177--216.

\bibitem{Marcovecchio09}
Raffaele Marcovecchio, \emph{The {R}hin--{V}iola method for {$\log 2$}}, Acta
  Arith. \textbf{139} (2009), no.~2, 147--184.

\bibitem{Nesterenko85}
Yu.~V. Nesterenko, \emph{On the linear independence of numbers}, Vestnik
  Moskov. Univ. Ser. I Mat. Mekh. (1985), no.~1, 46--49, In Russian; English
  translation: Moscow Univ. Math. Bull. 40 (1985), no. 1, 69--74.

\bibitem{NikishinSorokin91}
E.~M. Nikishin and V.~N. Sorokin, \emph{Rational approximations and
  orthogonality}, Translations of Mathematical Monographs, vol.~92, American
  Mathematical Society, Providence, RI, 1991.

\bibitem{OpenAI2026}
{OpenAI}, \emph{Catalan's constant is irrational}, Preprint dated September 24,
  2026, released October 6, 2026, with a Lean formalization, available at
  \href{https://github.com/openai/math/blob/fd4aeeb2ee4fc729c18d98444fed42fd0529eeeb/preprints/Catalans-constant-is-irrational-September-24-2026/paper.pdf}{the
  original paper PDF} and
  \href{https://github.com/openai/math/tree/fd4aeeb2ee4fc729c18d98444fed42fd0529eeeb/lean/OAI/NumberTheory/Catalan}{the
  Lean formalization} (repository commit \texttt{fd4aeeb}, October 7, 2026),
  2026.

\bibitem{OpenAIPi2026}
\bysame, \emph{The irrationality exponent of {$\pi$} is 2}, Preprint dated
  September 24, 2026, released October 6, 2026,
  \url{https://github.com/openai/math/tree/main/preprints/The-irrationality-exponent-of-pi-is-2-September-24-2026},
  2026.

\bibitem{Prellberg26}
Thomas Prellberg, \emph{Three irrationality results for the dilogarithm:
  {R}hin--{V}iola and {V}iola--{Z}udilin constructions at {$-1/4$}, {$1/5$},
  and {$-1/3$}}, 2026, Preprint, 7 September 2026,
  \href{https://arxiv.org/abs/2609.06895}{arXiv:2609.06895}.

\bibitem{RhinViola96}
Georges Rhin and Carlo Viola, \emph{On a permutation group related to
  {$\zeta(2)$}}, Acta Arith. \textbf{77} (1996), no.~1, 23--56.

\bibitem{RhinViola01}
\bysame, \emph{The group structure for {$\zeta(3)$}}, Acta Arith. \textbf{97}
  (2001), no.~3, 269--293.

\bibitem{RivoalZudilin2003}
Tanguy Rivoal and Wadim Zudilin, \emph{Diophantine properties of numbers
  related to {C}atalan's constant}, Math. Ann. \textbf{326} (2003), no.~4,
  705--721.

\bibitem{SaffTotik}
Edward~B. Saff and Vilmos Totik, \emph{Logarithmic potentials with external
  fields}, Grundlehren der mathematischen Wissenschaften, vol. 316, Springer,
  Berlin, 1997.

\bibitem{Sun2026}
Zhi-Wei Sun, \emph{Catalan's constant is irrational}, 2026,
  \href{https://arxiv.org/abs/2609.04176}{arXiv:2609.04176}.

\bibitem{Tang26}
Q.~Tang, \emph{zeta2-zeta3-linear-independence: a {L}ean 4 formalization of the
  linear independence of {$1$}, {$\zeta(2)$} and {$\zeta(3)$} over
  {$\mathbb{Q}$}}, GitHub repository dtq1997/zeta2-zeta3-linear-independence,
  version of 30 September 2026,
  \href{https://github.com/dtq1997/zeta2-zeta3-linear-independence}{github.com/dtq1997/zeta2-zeta3-linear-independence},
  2026.

\bibitem{Mathlib}
{The mathlib Community}, \emph{The {L}ean mathematical library}, Proceedings of
  the 9th ACM SIGPLAN International Conference on Certified Programs and Proofs
  (CPP 2020), 2020, pp.~367--381.

\bibitem{PNTPlus}
{The PrimeNumberTheoremAnd contributors}, \emph{The {PNT+} project: a {L}ean
  formalization of the prime number theorem and related results},
  \url{https://github.com/AlexKontorovich/PrimeNumberTheoremAnd}, The vendored
  compatibility additions used here are recorded in the repository's
  \texttt{NOTICE} and \texttt{PORTING.md}.

\bibitem{Wachs2026}
Dane Wachs, \emph{A note on a recent claimed proof of the irrationality of
  {C}atalan's constant}, 2026,
  \href{https://arxiv.org/abs/2609.22339}{arXiv:2609.22339}.

\bibitem{Xu2026}
Ryan Xu, \emph{Improved {$\kappa=2$} trial: {$\mu(G)\le727.03$}}, GitHub issue,
  October 22, 2026,
  \url{https://github.com/slitsey/catalan-irrationality-measure/issues/5}; data
  and scripts at \url{https://github.com/ry-xu/slightlysupervised}, 2026.

\bibitem{ZeilbergerZudilin20}
Doron Zeilberger and Wadim Zudilin, \emph{The irrationality measure of {$\pi$}
  is at most {$7.103205334137\ldots$}}, Mosc. J. Comb. Number Theory \textbf{9}
  (2020), no.~4, 407--419,
  \href{https://arxiv.org/abs/1912.06345}{arXiv:1912.06345}.

\bibitem{Zudilin02Cheb}
Wadim Zudilin, \emph{A few remarks on linear forms involving {C}atalan's
  constant}, Chebyshevskii Sb. \textbf{3} (2002), no.~2(4), 60--70, In Russian,
  \href{https://arxiv.org/abs/math/0210423}{arXiv:math/0210423}.

\bibitem{Zudilin2003}
\bysame, \emph{An {A}p\'ery-like difference equation for {C}atalan's constant},
  Electron. J. Combin. \textbf{10} (2003), Research Paper 14,
  \href{https://arxiv.org/abs/math/0201024}{arXiv:math/0201024}.

\bibitem{Zudilin14}
\bysame, \emph{Two hypergeometric tales and a new irrationality measure of
  {$\zeta(2)$}}, Ann. Math. Qu\'ebec \textbf{38} (2014), no.~1, 101--117,
  \href{https://arxiv.org/abs/1310.1526}{arXiv:1310.1526}.

\bibitem{Zudilin}
\bysame, \emph{A determinantal approach to irrationality}, Constr. Approx.
  \textbf{45} (2017), no.~2, 301--310,
  \href{https://arxiv.org/abs/1507.05697}{arXiv:1507.05697}.

\bibitem{Zudilin19Survey}
\bysame, \emph{Arithmetic of {C}atalan's constant and its relatives}, Abh.
  Math. Semin. Univ. Hambg. \textbf{89} (2019), no.~1, 45--53,
  \href{https://arxiv.org/abs/1804.09922}{arXiv:1804.09922}.

\end{thebibliography}

\end{document}
